\section{Formal schemes}
\label{section:I.10}

\subsection{Affine formal schemes.}
\label{subsection:I.10.1}

\begin{env}[10.1.1]
\label{I.10.1.1}
Let $A$ be an \emph{admissible} topological ring
(\hyperref[0.7.1.2]{0, 7.1.2}); for every ideal of definition
$\mathfrak{J}$ of $A$, $\operatorname{Spec}(A/\mathfrak{J})$ is identified
with the closed subspace $V(\mathfrak{J})$ of $\operatorname{Spec}(A)$
(\hyperref[I.1.1.11]{1.1.11}), the set of \emph{open} prime ideals
of $A$; this topological space does not depend on
\oldpage[I]{181}
the ideal of definition $\mathfrak{J}$ considered; let us denote it by
$\mathfrak{X}$. Let $(\mathfrak{J}_\lambda)$ be a fundamental system of
neighbourhoods of 0 in $A$ consisting of ideals of definition, and for every
$\lambda$, let $\mathscr{O}_\lambda$ be the structural sheaf of
$\operatorname{Spec}(A/\mathfrak{J}_\lambda)$; this sheaf is induced on
$\mathfrak{X}$ by
$\widetilde{A}/\widetilde{\mathfrak{J}}_\lambda$ (which is zero outside
$\mathfrak{X}$).

For $\mathfrak{J}_\mu\subset\mathfrak{J}_\lambda$, the canonical
homomorphism $A/\mathfrak{J}_\mu\to A/\mathfrak{J}_\lambda$ therefore
defines a homomorphism
$u_{\lambda\mu}:\mathscr{O}_\mu\to\mathscr{O}_\lambda$ of sheaves of
rings (\hyperref[I.1.6.1]{1.6.1}), and
$(\mathscr{O}_\lambda)$ is a \emph{projective system of sheaves of
rings} for these homomorphisms. Since the topology of $\mathfrak{X}$
has a basis consisting of quasi-compact open subsets, we can associate to
every $\mathscr{O}_\lambda$ a \emph{pseudo-discrete sheaf of topological
rings} (\hyperref[0.3.8.1]{0, 3.8.1}) which has
$\mathscr{O}_\lambda$ as its underlying sheaf of rings (without topology),
and which we will again denote by $\mathscr{O}_\lambda$; and the
$\mathscr{O}_\lambda$ again form a \emph{projective system of sheaves of
topological rings} (\hyperref[0.3.8.2]{0, 3.8.2}). We will denote
by $\mathscr{O}_{\mathfrak{X}}$ the \emph{sheaf of topological rings}
on $\mathfrak{X}$ that is the projective limit of the system
$(\mathscr{O}_\lambda)$; for every \emph{quasi-compact} open subset $U$ of
$\mathfrak{X}$, $\Gamma(U,\mathscr{O}_{\mathfrak{X}})$ is therefore the
topological ring that is the projective limit of the system of
\emph{discrete} rings $\Gamma(U,\mathscr{O}_\lambda)$
(\hyperref[0.3.2.6]{0, 3.2.6}).
\end{env}

\begin{definition}[10.1.2]
\label{I.10.1.2}
Given an admissible topological ring $A$, we call formal spectrum
of $A$, and denote by $\operatorname{Spf}(A)$, the closed subspace
$\mathfrak{X}$ of $\operatorname{Spec}(A)$ consisting of the open prime
ideals of $A$. We say that a topologically ringed space is an affine formal
scheme if it is isomorphic to a formal spectrum
$\operatorname{Spf}(A)=\mathfrak{X}$ endowed with the sheaf of topological
rings $\mathscr{O}_{\mathfrak{X}}$ that is the projective limit of the
pseudo-discrete sheaves of rings
$(\widetilde{A}/\widetilde{\mathfrak{J}}_\lambda)|\mathfrak{X}$, where
$\mathfrak{J}_\lambda$ runs through the filtered set of ideals of
definition of $A$.
\end{definition}

When we speak of a \emph{formal spectrum}
$\mathfrak{X}=\operatorname{Spf}(A)$ as an affine formal scheme, we will
always mean the topologically ringed space
$(\mathfrak{X},\mathscr{O}_{\mathfrak{X}})$ where
$\mathscr{O}_{\mathfrak{X}}$ is defined as above.

Note that every \emph{affine scheme} $X=\operatorname{Spec}(A)$ can be
considered as an affine formal scheme in a unique way, by considering
$A$ as a discrete topological ring: the topological rings
$\Gamma(U,\mathscr{O}_X)$ are then discrete when $U$ is
quasi-compact (but not in general when $U$ is an arbitrary open subset of
$X$).

\begin{proposition}[10.1.3]
\label{I.10.1.3}
If $\mathfrak{X}=\operatorname{Spf}(A)$, where $A$ is an admissible ring,
$\Gamma(\mathfrak{X},\mathscr{O}_{\mathfrak{X}})$ is topologically
isomorphic to $A$.
\end{proposition}

\begin{proof}
Indeed, since $\mathfrak{X}$ is closed in $\operatorname{Spec}(A)$, it
is quasi-compact, and consequently
$\Gamma(\mathfrak{X},\mathscr{O}_{\mathfrak{X}})$ is topologically
isomorphic to the projective limit of the discrete rings
$\Gamma(\mathfrak{X},\mathscr{O}_\lambda)$; but
$\Gamma(\mathfrak{X},\mathscr{O}_\lambda)$ is isomorphic to
$A/\mathfrak{J}_\lambda$ (\hyperref[I.1.3.7]{1.3.7}); since $A$ is
separated and complete, it is topologically isomorphic to
$\varprojlim A/\mathfrak{J}_\lambda$
(\hyperref[0.7.2.1]{0, 7.2.1}), whence the proposition.
\end{proof}

\begin{proposition}[10.1.4]
\label{I.10.1.4}
Let $A$ be an admissible ring, $\mathfrak{X}=\operatorname{Spf}(A)$, and
for every $f\in A$, let
$\mathfrak{D}(f)=D(f)\cap\mathfrak{X}$; the topologically ringed space
$(\mathfrak{D}(f),\mathscr{O}_{\mathfrak{X}}|\mathfrak{D}(f))$ is
isomorphic to the affine formal spectrum
$\operatorname{Spf}(A_{\{f\}})$
(\hyperref[0.7.6.15]{0, 7.6.15}).
\end{proposition}

\begin{proof}
For every ideal of definition $\mathfrak{J}$ of $A$, the discrete ring
$S_f^{-1}A/S_f^{-1}\mathfrak{J}$ is canonically identified with
$A_{\{f\}}/\mathfrak{J}_{\{f\}}$
(\hyperref[0.7.6.9]{0, 7.6.9}), hence
(\hyperref[I.1.2.5]{1.2.5} and \hyperref[I.1.2.6]{1.2.6}) the
topological space $\operatorname{Spf}(A_{\{f\}})$ is canonically identified
with $\mathfrak{D}(f)$. Moreover, for every quasi-compact open subset $U$
of $\mathfrak{X}$ contained in $\mathfrak{D}(f)$,
$\Gamma(U,\mathscr{O}_\lambda)$ is identified with the module of sections
of the structural sheaf of
$\operatorname{Spec}(S_f^{-1}A/S_f^{-1}\mathfrak{J}_\lambda)$ over
$U$ (\hyperref[I.1.3.6]{1.3.6}); hence, if we set
$\mathfrak{Y}=\operatorname{Spf}(A_{\{f\}})$,
$\Gamma(U,\mathscr{O}_{\mathfrak{X}})$ is identified with the module of
sections $\Gamma(U,\mathscr{O}_{\mathfrak{Y}})$, which proves the
proposition.
\end{proof}

\oldpage[I]{182}
\begin{env}[10.1.5]
\label{I.10.1.5}
As a sheaf of rings \emph{without topology}, the structural sheaf
$\mathscr{O}_{\mathfrak{X}}$ of
$\operatorname{Spf}(A)$ has, for every $x\in\mathfrak{X}$, a stalk which,
by virtue of (\hyperref[I.10.1.4]{10.1.4}), is identified with the
inductive limit
$\varinjlim_{f\notin\mathfrak{j}_x}A_{\{f\}}$. Consequently
(\hyperref[0.7.6.17]{0, 7.6.17} and
\hyperref[0.7.6.18]{7.6.18}):
\end{env}

\begin{proposition}[10.1.6]
\label{I.10.1.6}
For every $x\in\mathfrak{X}=\operatorname{Spf}(A)$, the stalk
$\mathscr{O}_x$ is a local ring whose residue field is isomorphic to
$k(x)=A_x/\mathfrak{j}_xA_x$. If moreover $A$ is adic and Noetherian,
$\mathscr{O}_x$ is a Noetherian ring.
\end{proposition}

Since $k(x)$ does not consist of 0 alone, we conclude from this result
that the \emph{support} of the sheaf of rings $\mathscr{O}_{\mathfrak{X}}$
is \emph{equal to $\mathfrak{X}$}.

\subsection{Morphisms of affine formal schemes.}
\label{subsection:I.10.2}

\begin{env}[10.2.1]
\label{I.10.2.1}
Let $A$, $B$ be two admissible rings, and let
$\varphi:B\to A$ be a \emph{continuous} homomorphism. The continuous map
$ {}^a\varphi:\operatorname{Spec}(A)\to\operatorname{Spec}(B)$
(\hyperref[I.1.2.1]{1.2.1}) then maps
$\mathfrak{X}=\operatorname{Spf}(A)$ into
$\mathfrak{Y}=\operatorname{Spf}(B)$, since the inverse image under
$\varphi$ of an open prime ideal of $A$ is an open prime ideal of
$B$. On the other hand, for every $g\in B$, $\varphi$ defines a continuous
homomorphism
\[
  \Gamma(\mathfrak{D}(g),\mathscr{O}_{\mathfrak{Y}})
  \longrightarrow
  \Gamma(\mathfrak{D}(\varphi(g)),\mathscr{O}_{\mathfrak{X}})
\]
by virtue of (\hyperref[I.10.1.4]{10.1.4}),
(\hyperref[I.10.1.3]{10.1.3}) and
(\hyperref[0.7.6.7]{0, 7.6.7}); since these homomorphisms satisfy the
compatibility conditions for the restrictions corresponding to the
passage from $g$ to a multiple of $g$, and since
$\mathfrak{D}(\varphi(g))={}^a\varphi^{-1}(\mathfrak{D}(g))$, they
define a \emph{continuous} homomorphism of sheaves of topological
rings
$\mathscr{O}_{\mathfrak{Y}}\to
{}^a\varphi_*(\mathscr{O}_{\mathfrak{X}})$
(\hyperref[0.3.2.5]{0, 3.2.5}), which we will again denote by
$\widetilde{\varphi}$; we have thus obtained a morphism
$\Phi=({}^a\varphi,\widetilde{\varphi})$ of topologically ringed spaces
$\mathfrak{X}\to\mathfrak{Y}$. Note that, as a homomorphism of
sheaves of rings without topology, $\widetilde{\varphi}$ defines a
homomorphism
$\widetilde{\varphi}_x^\sharp:
\mathscr{O}_{{}^a\varphi(x)}\to\mathscr{O}_x$ on the stalks, for every
$x\in\mathfrak{X}$.
\end{env}

\begin{proposition}[10.2.2]
\label{I.10.2.2}
Let $A$, $B$ be two admissible topological rings, and let
$\mathfrak{X}=\operatorname{Spf}(A)$,
$\mathfrak{Y}=\operatorname{Spf}(B)$. For a morphism
$u=(\psi,\theta):\mathfrak{X}\to\mathfrak{Y}$ of topologically ringed
spaces to be of the form $({}^a\varphi,\widetilde{\varphi})$, where
$\varphi$ is a continuous ring homomorphism $B\to A$, it is necessary and
sufficient that for every $x\in\mathfrak{X}$,
$\theta_x^\sharp$ be a local homomorphism
$\mathscr{O}_{\psi(x)}\to\mathscr{O}_x$.
\end{proposition}

\begin{proof}
The condition is necessary: indeed, let
$\mathfrak{p}=\mathfrak{j}_x\in\operatorname{Spf}(A)$, and let
$\mathfrak{q}=\varphi^{-1}(\mathfrak{j}_x)$; if
$g\notin\mathfrak{q}$, we therefore have $\varphi(g)\notin\mathfrak{p}$, and it
is immediate that the homomorphism
$B_{\{g\}}\to A_{\{\varphi(g)\}}$ deduced from $\varphi$
(\hyperref[0.7.6.7]{0, 7.6.7}) transforms
$\mathfrak{q}_{\{g\}}$ into a subset of
$\mathfrak{p}_{\{\varphi(g)\}}$; passing to the inductive limit, we
therefore see (taking into account (\hyperref[I.10.1.5]{10.1.5}) and
(\hyperref[0.7.6.17]{0, 7.6.17})) that
$\widetilde{\varphi}_x^\sharp$ is a local homomorphism.

Conversely, let $(\psi,\theta)$ be a morphism satisfying the condition of
the statement; by virtue of (\hyperref[I.10.1.3]{10.1.3}), $\theta$
defines a continuous ring homomorphism
\[
  \varphi=\Gamma(\theta):
  B=\Gamma(\mathfrak{Y},\mathscr{O}_{\mathfrak{Y}})
  \longrightarrow
  \Gamma(\mathfrak{X},\mathscr{O}_{\mathfrak{X}})=A.
\]
By virtue of the hypothesis on $\theta$, for the section $\varphi(g)$ of
$\mathscr{O}_{\mathfrak{X}}$ over $\mathfrak{X}$ to have an invertible
germ at the point $x$, it is necessary and sufficient that $g$ have an
invertible germ at the point $\psi(x)$. But by virtue of
(\hyperref[0.7.6.17]{0, 7.6.17}), the sections of
$\mathscr{O}_{\mathfrak{X}}$ (resp. $\mathscr{O}_{\mathfrak{Y}}$)
over $\mathfrak{X}$ (resp. $\mathfrak{Y}$) whose germ is not
invertible at the point $x$ (resp. $\psi(x)$) are exactly the elements of
$\mathfrak{j}_x$
\oldpage[I]{183}
(resp. $\mathfrak{j}_{\psi(x)}$); the preceding remark therefore shows that
${}^a\varphi=\psi$. Finally, for every $g\in B$, the diagram
\[
  \xymatrix{
    B=\Gamma(\mathfrak{Y},\mathscr{O}_{\mathfrak{Y}})
      \ar[r]^\varphi\ar[d] &
    \Gamma(\mathfrak{X},\mathscr{O}_{\mathfrak{X}})=A\ar[d]\\
    B_{\{g\}}=\Gamma(\mathfrak{D}(g),\mathscr{O}_{\mathfrak{Y}})
      \ar[r]_{\Gamma(\theta_{\mathfrak{D}(g)})} &
    \Gamma(\mathfrak{D}(\varphi(g)),\mathscr{O}_{\mathfrak{X}})
      =A_{\{\varphi(g)\}}
  }
\]
is commutative; by the universal property of complete rings of
fractions (\hyperref[0.7.6.6]{0, 7.6.6}),
$\theta_{\mathfrak{D}(g)}$ is equal to
$\widetilde{\varphi}_{\mathfrak{D}(g)}$ for every $g\in B$, hence
(\hyperref[0.3.2.5]{0, 3.2.5}) we have
$\theta=\widetilde{\varphi}$.
\end{proof}

We will say that a morphism $(\psi,\theta)$ of topologically ringed
spaces satisfying the condition of (\hyperref[I.10.2.2]{10.2.2}) is a
\emph{morphism of affine formal schemes}. We can say that the functors
$\operatorname{Spf}(A)$ in $A$ and
$\Gamma(\mathfrak{X},\mathscr{O}_{\mathfrak{X}})$ in $\mathfrak{X}$
define an \emph{equivalence} between the category of admissible rings
and the dual category of the category of affine formal schemes
(T, I, 1.2).

\begin{env}[10.2.3]
\label{I.10.2.3}
As a particular case of (\hyperref[I.10.2.2]{10.2.2}), note that, for
$f\in A$, the canonical injection of the affine formal scheme induced by
$\mathfrak{X}$ on $\mathfrak{D}(f)$ corresponds to the canonical continuous
homomorphism $A\to A_{\{f\}}$. Under the hypotheses of
(\hyperref[I.10.2.2]{10.2.2}), let $h$ be an element of $B$ and $g$ an
element of $A$ that is a multiple of $\varphi(h)$; we then have
$\psi(\mathfrak{D}(g))\subset\mathfrak{D}(h)$; the restriction of $u$ to
$\mathfrak{D}(g)$, considered as a morphism from $\mathfrak{D}(g)$ to
$\mathfrak{D}(h)$, is the unique morphism $v$ making the
diagram
\[
  \xymatrix{
    \mathfrak{D}(g)\ar[r]^v\ar[d] &
    \mathfrak{D}(h)\ar[d]\\
    \mathfrak{X}\ar[r]_u &
    \mathfrak{Y}
  }
\]
commutative. This morphism corresponds to the unique continuous homomorphism
$\varphi':B_{\{h\}}\to A_{\{g\}}$
(\hyperref[0.7.6.7]{0, 7.6.7}) making the diagram
\[
  \xymatrix{
    A\ar[d] &
    B\ar[l]_\varphi\ar[d]\\
    A_{\{g\}} &
    B_{\{h\}}\ar[l]_{\varphi'}
  }
\]
commutative.
\end{env}

\subsection{Ideals of definition of an affine formal scheme.}
\label{subsection:I.10.3}

\begin{env}[10.3.1]
\label{I.10.3.1}
Let $A$ be an admissible ring, $\mathfrak{J}$ an open ideal of $A$, and
$\mathfrak{X}$ the affine formal scheme $\operatorname{Spf}(A)$. Let
$(\mathfrak{J}_\lambda)$ be the set of ideals of definition of $A$
contained in $\mathfrak{J}$; then
$\widetilde{\mathfrak{J}}/\widetilde{\mathfrak{J}}_\lambda$ is a
sheaf of ideals of
$\widetilde{A}/\widetilde{\mathfrak{J}}_\lambda$. Let us denote by
$\mathfrak{J}^\Delta$ the projective limit of the sheaves induced on
$\mathfrak{X}$ by
$\widetilde{\mathfrak{J}}/\widetilde{\mathfrak{J}}_\lambda$, which
is identified with a \emph{sheaf of ideals} of
$\mathscr{O}_{\mathfrak{X}}$ (\hyperref[0.3.2.6]{0, 3.2.6}). For every
$f\in A$, $\Gamma(\mathfrak{D}(f),\mathfrak{J}^\Delta)$ is the projective
limit of
$S_f^{-1}\mathfrak{J}/S_f^{-1}\mathfrak{J}_\lambda$, in other words it
is identified with the open ideal $\mathfrak{J}_{\{f\}}$ of the ring
$A_{\{f\}}$ (\hyperref[0.7.6.9]{0, 7.6.9}), and in particular
$\Gamma(\mathfrak{X},\mathfrak{J}^\Delta)=\mathfrak{J}$; we conclude
(the $\mathfrak{D}(f)$ forming a base of the topology of $\mathfrak{X}$)
that we have
\[
  \mathfrak{J}^\Delta|\mathfrak{D}(f)
  =(\mathfrak{J}_{\{f\}})^\Delta.
  \tag{10.3.1.1}\label{I.10.3.1.1}
\]
\end{env}

\oldpage[I]{184}
\begin{env}[10.3.2]
\label{I.10.3.2}
With the notation of (\hyperref[I.10.3.1]{10.3.1}), for every $f\in A$,
the canonical map of
$A_{\{f\}}=\Gamma(\mathfrak{D}(f),\mathscr{O}_{\mathfrak{X}})$ into
$\Gamma\bigl(\mathfrak{D}(f),
(\widetilde{A}/\widetilde{\mathfrak{J}})|\mathfrak{X}\bigr)
=S_f^{-1}A/S_f^{-1}\mathfrak{J}$
is \emph{surjective} and has kernel
$\Gamma(\mathfrak{D}(f),\mathfrak{J}^\Delta)=\mathfrak{J}_{\{f\}}$
(\hyperref[0.7.6.9]{0, 7.6.9}); these maps therefore define a
continuous \emph{surjective} homomorphism, called \emph{canonical}, of the
sheaf of topological rings $\mathscr{O}_{\mathfrak{X}}$ onto the sheaf of
discrete rings
$(\widetilde{A}/\widetilde{\mathfrak{J}})|\mathfrak{X}$, whose kernel is
$\mathfrak{J}^\Delta$; this homomorphism is moreover none other than
$\widetilde{\varphi}$ (\hyperref[I.10.2.1]{10.2.1}), where $\varphi$ is
the continuous homomorphism $A\to A/\mathfrak{J}$; the morphism
$({}^a\varphi,\widetilde{\varphi}):
\operatorname{Spec}(A/\mathfrak{J})\to\mathfrak{X}$ of affine formal
schemes (where ${}^a\varphi$ is a homeomorphism onto the closed subspace $V(\mathfrak{J})$ of $\mathfrak{X}$)\footnote{Editorial correction: for an arbitrary open ideal, the image is not necessarily all of $\mathfrak{X}$. For example, if $A=k[t]$ is discrete and $\mathfrak{J}=(t)$, the image is the proper closed subset $V(t)$. The printed assertion holds when $\mathfrak{J}$ is an ideal of definition.} is also called \emph{canonical}. We therefore have,
by what precedes, a \emph{canonical isomorphism}
\[
  \mathscr{O}_{\mathfrak{X}}/\mathfrak{J}^\Delta
  \xrightarrow{\sim}
  (\widetilde{A}/\widetilde{\mathfrak{J}})|\mathfrak{X}.
  \tag{10.3.2.1}\label{I.10.3.2.1}
\]

It is clear (by virtue of
$\Gamma(\mathfrak{X},\mathfrak{J}^\Delta)=\mathfrak{J}$) that
the map $\mathfrak{J}\to\mathfrak{J}^\Delta$ is
\emph{strictly increasing}; by what precedes, for
$\mathfrak{J}\subset\mathfrak{J}'$, the sheaf
${\mathfrak{J}'}^\Delta/\mathfrak{J}^\Delta$ is canonically isomorphic to
$\widetilde{\mathfrak{J}'}/\widetilde{\mathfrak{J}}
=\widetilde{\mathfrak{J}'/\mathfrak{J}}$.
\end{env}

\begin{env}[10.3.3]
\label{I.10.3.3}
The hypotheses and notation still being those of
(\hyperref[I.10.3.1]{10.3.1}), we will say that a sheaf of ideals
$\mathscr{J}$ of $\mathscr{O}_{\mathfrak{X}}$ is a
\emph{sheaf of ideals of definition of $\mathfrak{X}$} (or an
\emph{ideal of definition of $\mathfrak{X}$}) if, for every
$x\in\mathfrak{X}$, there exists an open neighbourhood of $x$ of the form
$\mathfrak{D}(f)$, where $f\in A$, such that
$\mathscr{J}|\mathfrak{D}(f)$ is of the form $\mathfrak{H}^\Delta$, where
$\mathfrak{H}$ is an ideal of definition of $A_{\{f\}}$.
\end{env}

\begin{proposition}[10.3.4]
\label{I.10.3.4}
For every $f\in A$, every ideal of definition of $\mathfrak{X}$ induces an
ideal of definition of $\mathfrak{D}(f)$.
\end{proposition}

\begin{proof}
This follows from (\hyperref[I.10.3.1.1]{10.3.1.1}).
\end{proof}

\begin{proposition}[10.3.5]
\label{I.10.3.5}
If $A$ is an admissible ring, every ideal of definition of
$\mathfrak{X}=\operatorname{Spf}(A)$ is of the form
$\mathfrak{J}^\Delta$, where $\mathfrak{J}$ is an ideal of definition of
$A$, uniquely determined.
\end{proposition}

\begin{proof}
Indeed, let $\mathscr{J}$ be an ideal of definition of $\mathfrak{X}$;
by hypothesis, and since $\mathfrak{X}$ is quasi-compact, there are finitely
many elements $f_i\in A$ such that the $\mathfrak{D}(f_i)$ cover
$\mathfrak{X}$ and
$\mathscr{J}|\mathfrak{D}(f_i)=\mathfrak{H}_i^\Delta$, where
$\mathfrak{H}_i$ is an ideal of definition of $A_{\{f_i\}}$. For every
$i$, there therefore exists an open ideal $\mathfrak{K}_i$ of $A$ such that
$(\mathfrak{K}_i)_{\{f_i\}}=\mathfrak{H}_i$
(\hyperref[0.7.6.9]{0, 7.6.9}); let $\mathfrak{K}$ be an ideal of
definition of $A$ contained in all the $\mathfrak{K}_i$. The canonical image
of $\mathscr{J}/\mathfrak{K}^\Delta$ in the structural sheaf
$\widetilde{A/\mathfrak{K}}$ of
$\operatorname{Spec}(A/\mathfrak{K})$
(\hyperref[I.10.3.2]{10.3.2}) is therefore such that its restriction to
$\mathfrak{D}(f_i)$ is equal to that of
$\widetilde{\mathfrak{K}_i/\mathfrak{K}}$; we conclude that this canonical
image is a quasi-coherent sheaf on
$\operatorname{Spec}(A/\mathfrak{K})$, hence of the form
$\widetilde{\mathfrak{J}/\mathfrak{K}}$, where $\mathfrak{J}$ is an ideal of
$A$ containing $\mathfrak{K}$ (\hyperref[I.1.4.1]{1.4.1}), whence
$\mathscr{J}=\mathfrak{J}^\Delta$
(\hyperref[I.10.3.2]{10.3.2}); moreover, since for every $i$ there
exists an integer $n_i$ such that
$\mathfrak{H}_i^{n_i}\subset\mathfrak{K}_{\{f_i\}}$, we will have, denoting
by $n$ the largest of the $n_i$,
$(\mathscr{J}/\mathfrak{K}^\Delta)^n=0$, and consequently
(\hyperref[I.10.3.2]{10.3.2})
$(\widetilde{\mathfrak{J}/\mathfrak{K}})^n=0$, whence finally
$(\mathfrak{J}/\mathfrak{K})^n=0$
(\hyperref[I.1.3.13]{1.3.13}), which proves that $\mathfrak{J}$ is an
ideal of definition of $A$ (\hyperref[0.7.1.4]{0, 7.1.4}).
\end{proof}

\begin{proposition}[10.3.6]
\label{I.10.3.6}
Let $A$ be an adic ring and $\mathfrak{J}$ an ideal of definition of $A$
such that $\mathfrak{J}/\mathfrak{J}^2$ is an
$A/\mathfrak{J}$-module of finite type. Then for every integer $n>0$ we
have $(\mathfrak{J}^\Delta)^n=(\mathfrak{J}^n)^\Delta$.
\end{proposition}

\begin{proof}
Indeed, for every $f\in A$, we have (since $\mathfrak{J}^n$ is an open
ideal)
\[
  \bigl(\Gamma(\mathfrak{D}(f),\mathfrak{J}^\Delta)\bigr)^n
  =(\mathfrak{J}_{\{f\}})^n
  =(\mathfrak{J}^n)_{\{f\}}
  =\Gamma\bigl(\mathfrak{D}(f^n),(\mathfrak{J}^n)^\Delta\bigr)
\]
\oldpage[I]{185}
by virtue of (\hyperref[I.10.3.1.1]{10.3.1.1}) and
(\hyperref[0.7.6.12]{0, 7.6.12}). Since
$(\mathfrak{J}^\Delta)^n$ is associated to the presheaf
$U\mapsto(\Gamma(U,\mathfrak{J}^\Delta))^n$
(\hyperref[0.4.1.6]{0, 4.1.6}), the conclusion follows, since the
$\mathfrak{D}(f)$ form a base of the topology of $\mathfrak{X}$.
\end{proof}

\begin{env}[10.3.7]
\label{I.10.3.7}
We say that a family $(\mathscr{J}_\lambda)$ of ideals of definition of
$\mathfrak{X}$ is a \emph{fundamental system of ideals of definition}
if every ideal of definition of $\mathfrak{X}$ contains one of the
$\mathscr{J}_\lambda$; since
$\mathscr{J}_\lambda=\mathfrak{J}_\lambda^\Delta$, it amounts to the same to
say that the $\mathfrak{J}_\lambda$ form a \emph{fundamental system
of neighbourhoods of 0} in $A$. Let $(f_\alpha)$ be a family of elements of
$A$ such that the $\mathfrak{D}(f_\alpha)$ cover $\mathfrak{X}$. If
$(\mathscr{J}_\lambda)$ is a decreasing filtered family of ideals of
$\mathscr{O}_{\mathfrak{X}}$ such that for every $\alpha$ the family
$(\mathscr{J}_\lambda|\mathfrak{D}(f_\alpha))$ is a fundamental system
of ideals of definition of $\mathfrak{D}(f_\alpha)$, then
$(\mathscr{J}_\lambda)$ is a fundamental system of ideals of definition
of $\mathfrak{X}$. Indeed, for every ideal of definition $\mathscr{J}$
of $\mathfrak{X}$, there is a finite covering of $\mathfrak{X}$ by sets
$\mathfrak{D}(f_i)$ such that, for every $i$,
$\mathscr{J}_{\lambda_i}|\mathfrak{D}(f_i)$ is an ideal of definition
of $\mathfrak{D}(f_i)$ contained in
$\mathscr{J}|\mathfrak{D}(f_i)$. If $\mu$ is an index such that
$\mathscr{J}_\mu\subset\mathscr{J}_{\lambda_i}$ for every $i$, it follows
from (\hyperref[I.10.3.3]{10.3.3}) that $\mathscr{J}_\mu$ is an ideal of
definition of $\mathfrak{X}$, evidently contained in $\mathscr{J}$, whence
our assertion.
\end{env}

\subsection{Formal preschemes and morphisms of formal preschemes.}
\label{subsection:I.10.4}

\begin{env}[10.4.1]
\label{I.10.4.1}
Given a topologically ringed space $\mathfrak{X}$, we say that an
open subset $U\subset\mathfrak{X}$ is an \emph{affine formal open subset}
(resp. an \emph{adic affine formal open subset}, resp. a \emph{Noetherian
affine formal open subset}) if the topologically ringed space induced by
$\mathfrak{X}$ on $U$ is an affine formal scheme (resp. such a scheme whose
ring is adic, resp. adic and Noetherian).
\end{env}

\begin{definition}[10.4.2]
\label{I.10.4.2}
We call formal prescheme a topologically ringed space $\mathfrak{X}$
every point of which has an affine formal open neighbourhood. We say that
the formal prescheme $\mathfrak{X}$ is adic (resp. locally Noetherian) if
every point of $\mathfrak{X}$ has an adic (resp. Noetherian) affine formal
open neighbourhood. We say that $\mathfrak{X}$ is Noetherian if it is
locally Noetherian and its underlying space is quasi-compact (hence
Noetherian).
\end{definition}

\begin{proposition}[10.4.3]
\label{I.10.4.3}
If $\mathfrak{X}$ is a formal prescheme (resp. a locally Noetherian one),
the affine formal open sets (resp. Noetherian affine formal open sets) form
a base of the topology of $\mathfrak{X}$.
\end{proposition}

\begin{proof}
This follows from (\hyperref[I.10.4.2]{10.4.2}) and
(\hyperref[I.10.1.4]{10.1.4}), taking into account that if $A$ is a
Noetherian adic ring, so is $A_{\{f\}}$ for every
$f\in A$ (\hyperref[0.7.6.11]{0, 7.6.11}).
\end{proof}

\begin{corollary}[10.4.4]
\label{I.10.4.4}
If $\mathfrak{X}$ is a formal prescheme (resp. a locally Noetherian formal
prescheme, resp. a Noetherian one), the topologically ringed space
induced on any open subset of $\mathfrak{X}$ is again a formal prescheme
(resp. a locally Noetherian formal prescheme, resp. a Noetherian one).
\end{corollary}

\begin{definition}[10.4.5]
\label{I.10.4.5}
Given two formal preschemes $\mathfrak{X}$, $\mathfrak{Y}$, we call
morphism (of formal preschemes) from $\mathfrak{X}$ to
$\mathfrak{Y}$ every morphism $(\psi,\theta)$ of topologically ringed
spaces such that, for every $x\in\mathfrak{X}$, $\theta_x^\sharp$ is a
local homomorphism
$\mathscr{O}_{\psi(x)}\to\mathscr{O}_x$.
\end{definition}

It is immediate that the composite of two morphisms of formal preschemes is
again such a morphism; the formal preschemes therefore form a
\emph{category}, and we will denote by
$\operatorname{Hom}(\mathfrak{X},\mathfrak{Y})$ the set of morphisms
from a formal prescheme $\mathfrak{X}$ to a formal prescheme
$\mathfrak{Y}$.

\oldpage[I]{186}
If $U$ is an open subset of $\mathfrak{X}$, the canonical injection
into $\mathfrak{X}$ of the formal prescheme induced by $\mathfrak{X}$ on $U$
is a morphism of formal preschemes (and even a
\emph{monomorphism} of topologically ringed spaces
(\hyperref[0.4.1.1]{0, 4.1.1})).

\begin{proposition}[10.4.6]
\label{I.10.4.6}
Let $\mathfrak{X}$ be a formal prescheme and
$\mathfrak{S}=\operatorname{Spf}(A)$ an affine formal scheme. There is
a canonical bijective correspondence between the morphisms of the formal
prescheme $\mathfrak{X}$ into the formal prescheme $\mathfrak{S}$ and the
continuous homomorphisms of the ring $A$ into the topological ring
$\Gamma(\mathfrak{X},\mathscr{O}_{\mathfrak{X}})$.
\end{proposition}

The proof is the same as that of
(\hyperref[I.2.2.4]{2.2.4}), replacing ``homomorphism'' by
``continuous homomorphism'' and ``affine open subset'' by ``affine formal
open subset'', and using (\hyperref[I.10.2.2]{10.2.2}) instead of
(\hyperref[I.1.7.3]{1.7.3}); we leave the details to the reader.

\begin{env}[10.4.7]
\label{I.10.4.7}
Given a formal prescheme $\mathfrak{S}$, we say that the data of a
formal prescheme $\mathfrak{X}$ and of a morphism
$\varphi:\mathfrak{X}\to\mathfrak{S}$ define a formal prescheme
$\mathfrak{X}$ \emph{over $\mathfrak{S}$}, or a
\emph{formal $\mathfrak{S}$-prescheme}, $\varphi$ being called the
\emph{structural morphism} of the $\mathfrak{S}$-prescheme
$\mathfrak{X}$. If $\mathfrak{S}=\operatorname{Spf}(A)$, where $A$ is an
admissible ring, we also say that the formal $\mathfrak{S}$-prescheme
$\mathfrak{X}$ is a \emph{formal $A$-prescheme} or a formal prescheme
\emph{over $A$}. An arbitrary formal prescheme can always
be considered as a formal prescheme over $\mathbf{Z}$ (endowed
with the discrete topology).

If $\mathfrak{X}$, $\mathfrak{Y}$ are two formal $\mathfrak{S}$-preschemes,
we say that a morphism
$u:\mathfrak{X}\to\mathfrak{Y}$ is an
\emph{$\mathfrak{S}$-morphism} if the diagram
\[
  \xymatrix{
    \mathfrak{X}\ar[rr]^u\ar[rd] & &
    \mathfrak{Y}\ar[ld]\\
    & \mathfrak{S}
  }
\]
(where the oblique arrows are the structural morphisms) is
commutative. With this definition, the formal $\mathfrak{S}$-preschemes
(for fixed $\mathfrak{S}$) form a \emph{category}. We denote by
$\operatorname{Hom}_{\mathfrak{S}}(\mathfrak{X},\mathfrak{Y})$
the set of $\mathfrak{S}$-morphisms from the formal $\mathfrak{S}$-prescheme
$\mathfrak{X}$ to the formal $\mathfrak{S}$-prescheme
$\mathfrak{Y}$. When $\mathfrak{S}=\operatorname{Spf}(A)$, we also
say \emph{$A$-morphism} instead of
\emph{$\mathfrak{S}$-morphism}.
\end{env}

\begin{env}[10.4.8]
\label{I.10.4.8}
Since every affine scheme can be considered as an affine formal scheme
(\hyperref[I.10.1.2]{10.1.2}), every (ordinary) prescheme can be
considered as a formal prescheme. It follows moreover from
(\hyperref[I.10.4.5]{10.4.5}) that for \emph{ordinary} preschemes,
the morphisms (resp. $S$-morphisms) of \emph{formal} preschemes
coincide with the morphisms (resp. $S$-morphisms) defined in
\S~\hyperref[section:I.2]{2}.
\end{env}

\subsection{Ideals of definition of formal preschemes.}
\label{subsection:I.10.5}

\begin{env}[10.5.1]
\label{I.10.5.1}
Let $\mathfrak{X}$ be a formal prescheme; we say that an
$\mathscr{O}_{\mathfrak{X}}$-ideal $\mathscr{J}$ is a
\emph{sheaf of ideals of definition} (or an
\emph{ideal of definition}) of $\mathfrak{X}$ if every
$x\in\mathfrak{X}$ has an affine formal open neighbourhood $U$ such that
$\mathscr{J}|U$ is an ideal of definition of the affine formal scheme
induced by $\mathfrak{X}$ on $U$
(\hyperref[I.10.3.3]{10.3.3}); by virtue of
(\hyperref[I.10.3.1.1]{10.3.1.1}) and
(\hyperref[I.10.4.3]{10.4.3}), for every open subset
$V\subset\mathfrak{X}$, $\mathscr{J}|V$ is then an ideal of definition
of the formal prescheme induced by $\mathfrak{X}$ on $V$.

We say that a family $(\mathscr{J}_\lambda)$ of ideals of definition of
$\mathfrak{X}$ is a \emph{fundamental system}
\oldpage[I]{187}
\emph{of ideals of definition} if there exists a covering
$(U_\alpha)$ of $\mathfrak{X}$ by affine formal open subsets such that,
for every $\alpha$, the family of the
$\mathscr{J}_\lambda|U_\alpha$ is a fundamental system of ideals of
definition (\hyperref[I.10.3.7]{10.3.7}) of the affine formal scheme
induced by $\mathfrak{X}$ on $U_\alpha$. It follows from the final remark
of (\hyperref[I.10.3.7]{10.3.7}) that when $\mathfrak{X}$ is an
affine formal scheme, this definition coincides with the definition given
in (\hyperref[I.10.3.7]{10.3.7}). For every open subset $V$ of
$\mathfrak{X}$, the restrictions $\mathscr{J}_\lambda|V$ then form a
fundamental system of ideals of definition of the formal prescheme induced
on $V$, by virtue of (\hyperref[I.10.3.1.1]{10.3.1.1}). If
$\mathfrak{X}$ is a \emph{locally Noetherian} formal prescheme and
$\mathscr{J}$ an ideal of definition of $\mathfrak{X}$, it follows from
(\hyperref[I.10.3.6]{10.3.6}) that the powers $\mathscr{J}^n$
form a fundamental system of ideals of definition of $\mathfrak{X}$.
\end{env}

\begin{env}[10.5.2]
\label{I.10.5.2}
Let $\mathfrak{X}$ be a formal prescheme and $\mathscr{J}$ an ideal of
definition of $\mathfrak{X}$. Then the ringed space
$(\mathfrak{X},\mathscr{O}_{\mathfrak{X}}/\mathscr{J})$ is an (ordinary)
\emph{prescheme}, which is affine (resp. locally Noetherian,
resp. Noetherian) when $\mathfrak{X}$ is an affine formal scheme
(resp. a locally Noetherian formal prescheme, resp. a Noetherian one);
indeed, we are immediately reduced to the affine case, and then the
proposition has already been proved in (\hyperref[I.10.3.2]{10.3.2}).
Moreover, if
$\theta:\mathscr{O}_{\mathfrak{X}}\to
\mathscr{O}_{\mathfrak{X}}/\mathscr{J}$ is the canonical homomorphism,
$u=(1_{\mathfrak{X}},\theta)$ is a \emph{morphism} (called
\emph{canonical}) of formal preschemes
$(\mathfrak{X},\mathscr{O}_{\mathfrak{X}}/\mathscr{J})\to
(\mathfrak{X},\mathscr{O}_{\mathfrak{X}})$, since here again this was seen
in the affine case (\hyperref[I.10.3.2]{10.3.2}), to which we
immediately reduce.
\end{env}

\begin{proposition}[10.5.3]
\label{I.10.5.3}
Let $\mathfrak{X}$ be a formal prescheme and $(\mathscr{J}_\lambda)$ a
fundamental system of ideals of definition of $\mathfrak{X}$. Then the
sheaf of topological rings $\mathscr{O}_{\mathfrak{X}}$ is the projective
limit of the pseudo-discrete sheaves of rings
(\hyperref[0.3.8.1]{0, 3.8.1})
$\mathscr{O}_{\mathfrak{X}}/\mathscr{J}_\lambda$.
\end{proposition}

\begin{proof}
Since the topology of $\mathfrak{X}$ has a basis of quasi-compact affine
formal open subsets (\hyperref[I.10.4.3]{10.4.3}), we are reduced to the
affine case, where the proposition is a consequence of
(\hyperref[I.10.3.5]{10.3.5}),
(\hyperref[I.10.3.2]{10.3.2}) and of the definition
(\hyperref[I.10.1.1]{10.1.1}).
\end{proof}

It is not certain that every formal prescheme admits ideals of
definition. However:

\begin{proposition}[10.5.4]
\label{I.10.5.4}
Let $\mathfrak{X}$ be a locally Noetherian formal prescheme. There exists
a largest ideal of definition $\mathscr{T}$ of $\mathfrak{X}$;
it is the only ideal of definition $\mathscr{J}$ such that the prescheme
$(\mathfrak{X},\mathscr{O}_{\mathfrak{X}}/\mathscr{J})$ is reduced. If
$\mathscr{J}$ is an ideal of definition of $\mathfrak{X}$,
$\mathscr{T}$ is the inverse image under
$\mathscr{O}_{\mathfrak{X}}\to
\mathscr{O}_{\mathfrak{X}}/\mathscr{J}$ of the nilradical of
$\mathscr{O}_{\mathfrak{X}}/\mathscr{J}$.
\end{proposition}

\begin{proof}
Suppose first that $\mathfrak{X}=\operatorname{Spf}(A)$, where $A$ is
a Noetherian adic ring. The existence and the properties of
$\mathscr{T}$ follow immediately from
(\hyperref[I.10.3.4]{10.3.4}) and
(\hyperref[I.5.1.1]{5.1.1}), taking into account the existence and the
properties of the largest ideal of definition of $A$
(\hyperref[0.7.1.6]{0, 7.1.6} and
\hyperref[0.7.1.7]{7.1.7}).

To prove the existence and the properties of $\mathscr{T}$ in the general
case, it suffices to show that if $U\supset V$ are two Noetherian affine
formal open subsets of $\mathfrak{X}$, the largest ideal of
definition $\mathscr{T}_U$ of $U$ induces the largest ideal of definition
$\mathscr{T}_V$ of $V$; but since
$(V,(\mathscr{O}_{\mathfrak{X}}|V)/(\mathscr{T}_U|V))$ is reduced, this
follows from what precedes.
\end{proof}

We denote by $\mathfrak{X}_{\mathrm{red}}$ the reduced (ordinary)
prescheme $(\mathfrak{X},\mathscr{O}_{\mathfrak{X}}/\mathscr{T})$.

\begin{corollary}[10.5.5]
\label{I.10.5.5}
Let $\mathfrak{X}$ be a locally Noetherian formal prescheme and
$\mathscr{T}$ the largest ideal of definition of $\mathfrak{X}$; for
every open subset $V$ of $\mathfrak{X}$, $\mathscr{T}|V$ is the largest
ideal of definition of the formal prescheme induced by $\mathfrak{X}$ on
$V$.
\end{corollary}

\oldpage[I]{188}

\begin{proposition}[10.5.6]
\label{I.10.5.6}
Let $\mathfrak{X}$, $\mathfrak{Y}$ be two formal preschemes,
$\mathscr{J}$ (resp. $\mathscr{K}$) an ideal of definition of
$\mathfrak{X}$ (resp. $\mathfrak{Y}$), and
$f:\mathfrak{X}\to\mathfrak{Y}$ a morphism of formal preschemes.
\begin{enumerate}
  \item[\rm{(i)}] If
    $f^*(\mathscr{K})\mathscr{O}_{\mathfrak{X}}\subset\mathscr{J}$,
    there exists a unique morphism
    $f':(\mathfrak{X},\mathscr{O}_{\mathfrak{X}}/\mathscr{J})\to
    (\mathfrak{Y},\mathscr{O}_{\mathfrak{Y}}/\mathscr{K})$ of
    ordinary preschemes making the diagram
    \[
      \label{I.10.5.6.1}
      \xymatrix{
        (\mathfrak{X},\mathscr{O}_{\mathfrak{X}})\ar[r]^f &
        (\mathfrak{Y},\mathscr{O}_{\mathfrak{Y}})\\
        (\mathfrak{X},\mathscr{O}_{\mathfrak{X}}/\mathscr{J})
          \ar[r]_{f'}\ar[u] &
        (\mathfrak{Y},\mathscr{O}_{\mathfrak{Y}}/\mathscr{K})\ar[u]
      }
      \tag{10.5.6.1}
    \]
    commutative, where the vertical arrows are the canonical morphisms.
  \item[\rm{(ii)}] Suppose that
    $\mathfrak{X}=\operatorname{Spf}(A)$,
    $\mathfrak{Y}=\operatorname{Spf}(B)$ are affine formal
    schemes, $\mathscr{J}=\mathfrak{J}^{\Delta}$,
    $\mathscr{K}=\mathfrak{K}^{\Delta}$, where $\mathfrak{J}$ (resp.
    $\mathfrak{K}$) is an ideal of definition of $A$ (resp. $B$), and
    $f=({}^a\varphi,\widetilde{\varphi})$, where
    $\varphi:B\to A$ is a continuous homomorphism; for
    $f^*(\mathscr{K})\mathscr{O}_{\mathfrak{X}}\subset\mathscr{J}$,
    it is necessary and sufficient that
    $\varphi(\mathfrak{K})\subset\mathfrak{J}$, and $f'$ is then the
    morphism $({}^a\varphi',\widetilde{\varphi'})$, where
    $\varphi':B/\mathfrak{K}\to A/\mathfrak{J}$ is the homomorphism
    deduced from $\varphi$ by passing to the quotients.
\end{enumerate}
\end{proposition}

\begin{proof}
\medskip\noindent
\begin{enumerate}
  \item[\rm{(i)}] If $f=(\psi,\theta)$, the hypothesis implies that
    the image under
    $\theta^\sharp:\psi^*(\mathscr{O}_{\mathfrak{Y}})\to
    \mathscr{O}_{\mathfrak{X}}$ of the sheaf of ideals
    $\psi^*(\mathscr{K})$ of $\psi^*(\mathscr{O}_{\mathfrak{Y}})$ is
    contained in $\mathscr{J}$
    (\hyperref[0.4.3.5]{0, 4.3.5}). By passing to the quotients, we
    therefore deduce from $\theta^\sharp$ a homomorphism of sheaves of rings
    \[
      \omega:\psi^*(\mathscr{O}_{\mathfrak{Y}}/\mathscr{K})
      =\psi^*(\mathscr{O}_{\mathfrak{Y}})/\psi^*(\mathscr{K})
      \to\mathscr{O}_{\mathfrak{X}}/\mathscr{J}~;
    \]
    moreover, since for every $x\in\mathfrak{X}$,
    $\theta_x^\sharp$ is a \emph{local} homomorphism, the same is true
    of $\omega_x$. The morphism of ringed spaces
    $(\psi,\omega^b)$ is therefore
    (\hyperref[I.2.2.1]{2.2.1}) the unique morphism $f'$ of ringed
    spaces answering the question.
  \item[\rm{(ii)}] The canonical functorial correspondence between
    morphisms of affine formal preschemes and continuous ring
    homomorphisms (\hyperref[I.10.2.2]{10.2.2}) shows that in the case
    under consideration, the relation
    $f^*(\mathscr{K})\mathscr{O}_{\mathfrak{X}}\subset\mathscr{J}$
    implies that we have
    $f'=({}^a\varphi',\widetilde{\varphi'})$, where
    $\varphi':B/\mathfrak{K}\to A/\mathfrak{J}$ is the unique
    homomorphism making the diagram
    \[
      \label{I.10.5.6.2}
      \xymatrix{
        B\ar[r]^\varphi\ar[d] & A\ar[d]\\
        B/\mathfrak{K}\ar[r]_{\varphi'} & A/\mathfrak{J}
      }
      \tag{10.5.6.2}
    \]
    commutative. The existence of $\varphi'$ therefore implies that
    $\varphi(\mathfrak{K})\subset\mathfrak{J}$. Conversely, if this
    condition is satisfied, denoting by $\varphi'$ the unique
    homomorphism making the diagram
    (\hyperref[I.10.5.6.2]{10.5.6.2}) commutative and setting
    $f'=({}^a\varphi',\widetilde{\varphi'})$, it is clear that the
    diagram (\hyperref[I.10.5.6.1]{10.5.6.1}) is commutative; the
    consideration of the homomorphisms
    ${}^a\varphi^*(\mathscr{O}_{\mathfrak{Y}})\to
    \mathscr{O}_{\mathfrak{X}}$ and
    ${}^a{\varphi'}^*(\mathscr{O}_{\mathfrak{Y}}/\mathscr{K})\to
    \mathscr{O}_{\mathfrak{X}}/\mathscr{J}$ corresponding to $f$ and
    $f'$ respectively then shows that this implies the relation
    $f^*(\mathscr{K})\mathscr{O}_{\mathfrak{X}}\subset\mathscr{J}$.
\end{enumerate}
\end{proof}

It is clear that the correspondence $f\mapsto f'$ defined above is
\emph{functorial}.


\subsection{Formal preschemes as inductive limits of preschemes.}
\label{subsection:I.10.6}

\begin{env}[10.6.1]
\label{I.10.6.1}
Let $\mathfrak{X}$ be a formal prescheme,
$(\mathscr{J}_\lambda)$ a fundamental system of ideals of definition of
$\mathfrak{X}$; for every $\lambda$, let $f_\lambda$ be the canonical
morphism
$(\mathfrak{X},\mathscr{O}_{\mathfrak{X}}/\mathscr{J}_\lambda)
\to\mathfrak{X}$ (\hyperref[I.10.5.2]{10.5.2}); for
$\mathscr{J}_\mu\subset\mathscr{J}_\lambda$, the canonical homomorphism
$\mathscr{O}_{\mathfrak{X}}/\mathscr{J}_\mu\to
\mathscr{O}_{\mathfrak{X}}/\mathscr{J}_\lambda$ defines a canonical
morphism

\oldpage[I]{189}

$f_{\mu\lambda}:(\mathfrak{X},
\mathscr{O}_{\mathfrak{X}}/\mathscr{J}_\lambda)\to
(\mathfrak{X},\mathscr{O}_{\mathfrak{X}}/\mathscr{J}_\mu)$ of
(ordinary) preschemes such that $f_\lambda=f_\mu\circ f_{\mu\lambda}$.
The preschemes
$X_\lambda=(\mathfrak{X},
\mathscr{O}_{\mathfrak{X}}/\mathscr{J}_\lambda)$ and the morphisms
$f_{\mu\lambda}$ therefore form (by virtue of
\hyperref[I.10.4.8]{10.4.8}) an \emph{inductive system} in the
category of formal preschemes.
\end{env}

\begin{proposition}[10.6.2]
\label{I.10.6.2}
With the notation of (\hyperref[I.10.6.1]{10.6.1}), the formal
prescheme $\mathfrak{X}$ and the morphisms $f_\lambda$ form an
inductive limit (T, I, 1.8) of the system $(X_\lambda,f_{\mu\lambda})$
in the category of formal preschemes.
\end{proposition}

\begin{proof}
Let $\mathfrak{Y}$ be a formal prescheme, and for each index $\lambda$,
let
\[
  g_\lambda=(\psi_\lambda,\theta_\lambda):X_\lambda\to\mathfrak{Y}
\]
be a morphism such that
$g_\lambda=g_\mu\circ f_{\mu\lambda}$ for
$\mathscr{J}_\mu\subset\mathscr{J}_\lambda$. This last condition
and the definition of the $X_\lambda$ first of all imply that all the
$\psi_\lambda$ are identical to one and the same continuous map
$\psi:\mathfrak{X}\to\mathfrak{Y}$ of the underlying spaces; moreover,
the homomorphisms
$\theta_\lambda^\sharp:\psi^*(\mathscr{O}_{\mathfrak{Y}})\to
\mathscr{O}_{X_\lambda}=
\mathscr{O}_{\mathfrak{X}}/\mathscr{J}_\lambda$ form a
\emph{projective system} of homomorphisms of sheaves of rings. By
passing to the projective limit, we therefore deduce from them a
homomorphism
\[
  \omega:\psi^*(\mathscr{O}_{\mathfrak{Y}})\to
  \varprojlim\mathscr{O}_{\mathfrak{X}}/\mathscr{J}_\lambda
  =\mathscr{O}_{\mathfrak{X}},
\]
and it is clear that the morphism $g=(\psi,\omega^b)$ of \emph{ringed
spaces} is the \emph{only} one making the diagrams
\[
  \label{I.10.6.2.1}
  \xymatrix{
    X_\lambda\ar[rr]^{g_\lambda}\ar[rd]_{f_\lambda} &&
    \mathfrak{Y}\\
    & \mathfrak{X}\ar[ru]_g
  }
  \tag{10.6.2.1}
\]
commutative. It therefore remains to prove that $g$ is a morphism of
\emph{formal preschemes}; the question being local on $\mathfrak{X}$ and
$\mathfrak{Y}$, we may therefore suppose
$\mathfrak{X}=\operatorname{Spf}(A)$,
$\mathfrak{Y}=\operatorname{Spf}(B)$, $A$ and $B$ being admissible
rings, with
$\mathscr{J}_\lambda=\mathfrak{J}_\lambda^{\Delta}$, where
$(\mathfrak{J}_\lambda)$ is a fundamental system of ideals of
definition of $A$ (\hyperref[I.10.3.5]{10.3.5}); since
$A=\varprojlim A/\mathfrak{J}_\lambda$, the existence of a morphism $g$
of affine formal schemes making the diagrams
(\hyperref[I.10.6.2.1]{10.6.2.1}) commutative then follows from the
bijective correspondence (\hyperref[I.10.2.2]{10.2.2}) between
morphisms of affine formal schemes and continuous ring homomorphisms,
and from the definition of the projective limit. But the uniqueness of
$g$ as a morphism of ringed spaces shows that it coincides with the
morphism denoted in the same way at the beginning of the proof.
\end{proof}

Under certain additional conditions, the next proposition shows that a given
inductive system of (ordinary) preschemes has an inductive limit in the
category of formal preschemes:

\begin{proposition}[10.6.3]
\label{I.10.6.3}
Let $\mathfrak{X}$ be a topological space,
$(\mathscr{O}_i,u_{ji})$ a projective system of sheaves of rings on
$\mathfrak{X}$, having $\mathbf{N}$ as index set. Let
$\mathscr{J}_i$ be the kernel of
$u_{0i}:\mathscr{O}_i\to\mathscr{O}_0$. Suppose that:
\begin{enumerate}
  \item[a)] The ringed space $(\mathfrak{X},\mathscr{O}_i)$ is a
    prescheme $X_i$.
  \item[b)] For every $x\in X$ and every $i$, there exists an open
    neighbourhood $U_i$ of $x$ in $\mathfrak{X}$ such that the
    restriction $\mathscr{J}_i|U_i$ is nilpotent.
  \item[c)] The homomorphisms $u_{ji}$ are surjective.
\end{enumerate}

\oldpage[I]{190}

Let $\mathscr{O}_{\mathfrak{X}}$ be the sheaf of topological rings
that is the projective limit of the pseudo-discrete sheaves of rings
$\mathscr{O}_i$, and let
$u_i:\mathscr{O}_{\mathfrak{X}}\to\mathscr{O}_i$ be the canonical
homomorphism. Then the topologically ringed space
$(\mathfrak{X},\mathscr{O}_{\mathfrak{X}})$ is a formal prescheme;
the homomorphisms $u_i$ are surjective; their kernels
$\mathscr{J}^{(i)}$ form a fundamental system of ideals of definition
of $\mathfrak{X}$, and $\mathscr{J}^{(0)}$ is the projective limit of the
sheaves of ideals $\mathscr{J}_i$.
\end{proposition}

\begin{proof}
Let us first note that on each stalk, $u_{ji}$ is a surjective
homomorphism and \emph{a fortiori} a local homomorphism; hence
$v_{ij}=(1_{\mathfrak{X}},u_{ji})$ is a morphism of preschemes
$X_j\to X_i$ ($i\geq j$) (\hyperref[I.2.2.1]{2.2.1}). Suppose
first that each $X_i$ is an affine scheme with ring $A_i$. There exists
a homomorphism \emph{of rings}
$\varphi_{ji}:A_i\to A_j$ such that
$u_{ji}=\widetilde{\varphi_{ji}}$
(\hyperref[I.1.7.3]{1.7.3}); consequently
(\hyperref[I.1.6.3]{1.6.3}), the sheaf $\mathscr{O}_j$ is a
quasi-coherent $\mathscr{O}_i$-module on $X_i$ (for the external law
defined by $u_{ji}$), associated to $A_j$ regarded as an $A_i$-module by
means of $\varphi_{ji}$. For every $f\in A_i$, let
$f'=\varphi_{ji}(f)$; by hypothesis, the open subsets $D(f)$ and $D(f')$
are identical in $\mathfrak{X}$, and the homomorphism from
$\Gamma(D(f),\mathscr{O}_i)=(A_i)_f$ to
$\Gamma(D(f),\mathscr{O}_j)=(A_j)_{f'}$ corresponding to $u_{ji}$ is
none other than $(\varphi_{ji})_f$ (\hyperref[I.1.6.1]{1.6.1}). But
when $A_j$ is regarded as an $A_i$-module, $(A_j)_{f'}$ is the
$(A_i)_f$-module $(A_j)_f$, so we also have
$u_{ji}=\widetilde{\varphi_{ji}}$ when $\varphi_{ji}$ is this time
regarded as a homomorphism of $A_i$-modules. Then, since $u_{ji}$ is
surjective, we conclude that $\varphi_{ji}$ is surjective as well
(\hyperref[I.1.3.9]{1.3.9}), and if $\mathfrak{J}_{ji}$ is the kernel of
$\varphi_{ji}$, the kernel of $u_{ji}$ is a quasi-coherent
$\mathscr{O}_i$-module equal to $\widetilde{\mathfrak{J}_{ji}}$. In
particular, we have $\mathscr{J}_i=\widetilde{\mathfrak{J}_i}$, where
$\mathfrak{J}_i$ is the kernel of $\varphi_{0i}:A_i\to A_0$. Hypothesis
b) implies that $\mathscr{J}_i$ is \emph{nilpotent}: indeed, since
$\mathfrak{X}$ is quasi-compact, we can cover $\mathfrak{X}$ by finitely
many open subsets $U_k$ such that $(\mathscr{J}_i|U_k)^{n_k}=0$, and,
taking for $n$ the largest of the $n_k$, we have $\mathscr{J}_i^n=0$. We
conclude that $\mathfrak{J}_i$ is nilpotent
(\hyperref[I.1.3.13]{1.3.13}). Then the ring $A=\varprojlim A_i$ is
admissible (\hyperref[0.7.2.2]{0, 7.2.2}), the canonical homomorphism
$\varphi_i:A\to A_i$ is surjective and its kernel
$\mathfrak{J}^{(i)}$ is equal to the projective limit of the
$\mathfrak{J}_{ik}$ for $k\geq i$; the $\mathfrak{J}^{(i)}$ form
a fundamental system of neighbourhoods of 0 in $A$. The assertions of
(\hyperref[I.10.6.3]{10.6.3}) follow in this case from
(\hyperref[I.10.1.1]{10.1.1}) and
(\hyperref[I.10.3.2]{10.3.2}),
$(\mathfrak{X},\mathscr{O}_{\mathfrak{X}})$ being none other than
$\operatorname{Spf}(A)$.

Still in this same particular case, let us note that if $f=(f_i)$ is an
element of the projective limit $A=\varprojlim A_i$, all the open subsets
$D(f_i)$ (an affine open subset in $X_i$) are identified with the open
subset $\mathfrak{D}(f)$ of $\mathfrak{X}$, the prescheme induced by
$X_i$ on $\mathfrak{D}(f)$ therefore being identified with the affine
scheme $\operatorname{Spec}((A_i)_{f_i})$.

In the general case, let us first remark that for every quasi-compact
open subset $U$ of $\mathfrak{X}$, each of the $\mathscr{J}_i|U$ is
nilpotent, as the argument given above shows. We will see that for every
$x\in\mathfrak{X}$, there is an open neighbourhood $U$ of $x$ in
$\mathfrak{X}$ which is an \emph{affine open subset for all the $X_i$}.
Indeed, take for $U$ an affine open subset for $X_0$, and observe that
$\mathscr{O}_{X_0}=\mathscr{O}_{X_i}/\mathscr{J}_i$. Since
$\mathscr{J}_i|U$ is nilpotent, by virtue of the foregoing, $U$ is also
an affine open subset for each $X_i$ by virtue of
(\hyperref[I.5.1.9]{5.1.9}). This being so, for every $U$ satisfying
the preceding conditions, the study of the affine case made above shows
that $(U,\mathscr{O}_{\mathfrak{X}}|U)$ is a formal prescheme of which
the $\mathscr{J}^{(i)}|U$ form a fundamental system of ideals of
definition, and $\mathscr{J}^{(0)}|U$ is the projective limit of the
$\mathscr{J}_i|U$; whence the conclusion.
\end{proof}

\begin{corollary}[10.6.4]
\label{I.10.6.4}
Suppose that, for $i\geq j$, the kernel of $u_{ji}$ is
$\mathscr{J}_i^{j+1}$ and that $\mathscr{J}_1/\mathscr{J}_1^2$

\oldpage[I]{191}

is of finite type over
$\mathscr{O}_0=\mathscr{O}_1/\mathscr{J}_1$. Then $\mathfrak{X}$ is
an adic formal prescheme, and if $\mathscr{J}^{(n)}$ is the kernel of
$\mathscr{O}_{\mathfrak{X}}\to\mathscr{O}_n$, we have
$\mathscr{J}^{(n)}=\mathscr{J}^{n+1}$ and
$\mathscr{J}/\mathscr{J}^2$ is isomorphic to $\mathscr{J}_1$. If
moreover $X_0$ is locally Noetherian (resp. Noetherian),
$\mathfrak{X}$ is locally Noetherian (resp. Noetherian).
\end{corollary}

\begin{proof}
Since the underlying spaces of $\mathfrak{X}$ and $X_0$ are the same,
the question is local and we may suppose all the $X_i$ affine;
taking into account the relations
$\mathscr{J}_{ij}=\widetilde{\mathfrak{J}_{ji}}$ (with the notation
of (\hyperref[I.10.6.3]{10.6.3})), we are immediately reduced to the
corresponding assertions of
(\hyperref[0.7.2.7]{0, 7.2.7} and
\hyperref[0.7.2.8]{7.2.8}), noting that
$\mathfrak{J}_1/\mathfrak{J}_1^2$ is then an $A_0$-module of finite
type (\hyperref[I.1.3.9]{1.3.9}).
\end{proof}

In particular, \emph{every locally Noetherian formal prescheme
$\mathfrak{X}$} is the inductive limit of a sequence $(X_n)$ of locally
Noetherian (ordinary) preschemes satisfying the conditions of
(\hyperref[I.10.6.3]{10.6.3}) and
(\hyperref[I.10.6.4]{10.6.4}): it suffices to consider an ideal of
definition $\mathscr{J}$ of $\mathfrak{X}$
(\hyperref[I.10.5.4]{10.5.4}) and to take
$X_n=(\mathfrak{X},
\mathscr{O}_{\mathfrak{X}}/\mathscr{J}^{n+1})$
((\hyperref[I.10.5.1]{10.5.1}) and
(\hyperref[I.10.6.2]{10.6.2})).

\begin{corollary}[10.6.5]
\label{I.10.6.5}
Let $A$ be an admissible ring. In order that the affine formal scheme
$\mathfrak{X}=\operatorname{Spf}(A)$ be Noetherian, it is necessary and
sufficient that $A$ be adic and Noetherian.
\end{corollary}

\begin{proof}
The condition is evidently sufficient. Conversely, suppose that
$\mathfrak{X}$ is Noetherian, and let $\mathfrak{J}$ be an ideal of
definition of $A$ and $\mathscr{J}=\mathfrak{J}^\Delta$ the corresponding
ideal of definition of $\mathfrak{X}$. The (ordinary) preschemes
$X_n=(\mathfrak{X},
\mathscr{O}_{\mathfrak{X}}/\mathscr{J}^{n+1})$ are then affine and
Noetherian, so the rings
$A_n=A/\mathfrak{J}^{n+1}$ are Noetherian
(\hyperref[I.6.1.3]{6.1.3}), from which we conclude that
$\mathfrak{J}/\mathfrak{J}^2$ is an $A/\mathfrak{J}$-module of finite
type. Since the $\mathscr{J}^n$ form a fundamental system of ideals
of definition of $\mathfrak{X}$
(\hyperref[I.10.5.1]{10.5.1}), we have
$\mathscr{O}_{\mathfrak{X}}
=\varprojlim(\mathscr{O}_{\mathfrak{X}}/\mathscr{J}^n)$
(\hyperref[I.10.5.3]{10.5.3}); we conclude from this
(\hyperref[I.10.1.3]{10.1.3}) that $A$ is topologically isomorphic to
$\varprojlim A/\mathfrak{J}^n$, and is therefore adic and Noetherian
(\hyperref[0.7.2.8]{0, 7.2.8}).
\end{proof}

\begin{remark}[10.6.6]
\label{I.10.6.6}
With the notation of (\hyperref[I.10.6.3]{10.6.3}), let
$\mathscr{F}_i$ be an $\mathscr{O}_i$-module, and suppose given, for
$i\geq j$, a $v_{ij}$-morphism
$\theta_{ji}:\mathscr{F}_i\to\mathscr{F}_j$, in such a way that
$\theta_{kj}\circ\theta_{ji}=\theta_{ki}$ for $k\leq j\leq i$.
Since the continuous map underlying $v_{ij}$ is the identity,
$\theta_{ji}$ is a homomorphism of sheaves of abelian groups on
the space $\mathfrak{X}$; moreover, if $\mathscr{F}$ is the projective
limit of the projective system $(\mathscr{F}_i)$ of sheaves of abelian
groups, the fact that the $\theta_{ji}$ are $v_{ij}$-morphisms
makes it possible to define on $\mathscr{F}$ a structure of
$\mathscr{O}_{\mathfrak{X}}$-module by passing to the projective limit;
we will say that $\mathscr{F}$, endowed with this structure, is the
\emph{projective limit} (for the $\theta_{ji}$) of the system of
$\mathscr{O}_i$-modules $(\mathscr{F}_i)$. In the particular case where
$v_{ij}^*(\mathscr{F}_i)=\mathscr{F}_j$ and $\theta_{ji}$ is
the \emph{identity}, we will say for brevity that $\mathscr{F}$ is the
projective limit of a system $(\mathscr{F}_i)$ such that
$v_{ij}^*(\mathscr{F}_i)=\mathscr{F}_j$ for $j\leq i$ (without
mentioning the $\theta_{ji}$).
\end{remark}

\begin{env}[10.6.7]
\label{I.10.6.7}
Let $\mathfrak{X}$, $\mathfrak{Y}$ be two formal preschemes,
$\mathscr{J}$ (resp. $\mathscr{K}$) an ideal of definition of
$\mathfrak{X}$ (resp. $\mathfrak{Y}$), and
$f:\mathfrak{X}\to\mathfrak{Y}$ a morphism such that
$f^*(\mathscr{K})\mathscr{O}_{\mathfrak{X}}\subset\mathscr{J}$. We
then have, for every integer $n>0$,
$f^*(\mathscr{K}^n)\mathscr{O}_{\mathfrak{X}}
=(f^*(\mathscr{K})\mathscr{O}_{\mathfrak{X}})^n
\subset\mathscr{J}^n$; we can therefore
(\hyperref[I.10.5.6]{10.5.6}) deduce from $f$ a morphism of
(ordinary) preschemes $f_n:X_n\to Y_n$, by setting
$X_n=(\mathfrak{X},
\mathscr{O}_{\mathfrak{X}}/\mathscr{J}^{n+1})$,
$Y_n=(\mathfrak{Y},
\mathscr{O}_{\mathfrak{Y}}/\mathscr{K}^{n+1})$, and it follows
immediately from the definitions that the diagrams
\[
  \label{I.10.6.7.1}
  \xymatrix{
    X_m\ar[r]^{f_m}\ar[d] &
    Y_m\ar[d]\\
    X_n\ar[r]_{f_n} &
    Y_n
  }
  \tag{10.6.7.1}
\]

\oldpage[I]{192}

are commutative for $m\leq n$; in other words, $(f_n)$ is an
\emph{inductive system} of morphisms.
\end{env}

\begin{env}[10.6.8]
\label{I.10.6.8}
Conversely, let $(X_n)$ (resp. $(Y_n)$) be an inductive system of
(ordinary) preschemes satisfying conditions b) and c) of
(\hyperref[I.10.6.3]{10.6.3}), and let $\mathfrak{X}$ (resp.
$\mathfrak{Y}$) be its inductive limit. By definition of inductive
limits, every sequence $(f_n)$ of morphisms $X_n\to Y_n$ forming an
inductive system has an \emph{inductive limit}
$f:\mathfrak{X}\to\mathfrak{Y}$, which is the unique morphism of
formal preschemes making the diagrams
\[
  \xymatrix{
    X_n\ar[r]^{f_n}\ar[d] &
    Y_n\ar[d]\\
    \mathfrak{X}\ar[r]_f &
    \mathfrak{Y}
  }
\]
commutative.
\end{env}

\begin{proposition}[10.6.9]
\label{I.10.6.9}
Let $\mathfrak{X}$, $\mathfrak{Y}$ be two locally Noetherian formal
preschemes, and $\mathscr{J}$ (resp. $\mathscr{K}$) an ideal
of definition of $\mathfrak{X}$ (resp. $\mathfrak{Y}$);
the map $f\mapsto(f_n)$ defined in
(\hyperref[I.10.6.7]{10.6.7}) is a bijection of the set of
morphisms $f:\mathfrak{X}\to\mathfrak{Y}$ such that
$f^*(\mathscr{K})\mathscr{O}_{\mathfrak{X}}\subset\mathscr{J}$ onto
the set of sequences $(f_n)$ of morphisms making the
diagrams (\hyperref[I.10.6.7.1]{10.6.7.1}) commutative.
\end{proposition}

\begin{proof}
If $f$ is the inductive limit of such a sequence, we must show that
$f^*(\mathscr{K})\mathscr{O}_{\mathfrak{X}}\subset\mathscr{J}$. The
question being local on $\mathfrak{X}$ and $\mathfrak{Y}$, we may
restrict to the case where $\mathfrak{X}=\operatorname{Spf}(A)$,
$\mathfrak{Y}=\operatorname{Spf}(B)$ are affine, $A$ and $B$ being
Noetherian adic rings,
$\mathscr{J}=\mathfrak{J}^\Delta$,
$\mathscr{K}=\mathfrak{K}^\Delta$, where $\mathfrak{J}$ (resp.
$\mathfrak{K}$) is an ideal of definition of $A$ (resp. $B$). We then
have $X_n=\operatorname{Spec}(A_n)$,
$Y_n=\operatorname{Spec}(B_n)$, with
$A_n=A/\mathfrak{J}^{n+1}$ and $B_n=B/\mathfrak{K}^{n+1}$, by virtue of
(\hyperref[I.10.3.6]{10.3.6}) and
(\hyperref[I.10.3.2]{10.3.2});
$f_n=({}^a\varphi_n,\widetilde{\varphi_n})$, where the homomorphisms
$\varphi_n:B_n\to A_n$ form a projective system, hence
$f=({}^a\varphi,\widetilde{\varphi})$, where
$\varphi=\varprojlim\varphi_n$. The commutativity of the diagram
(\hyperref[I.10.6.7.1]{10.6.7.1}) for $m=0$ then gives the
condition
$\varphi_n(\mathfrak{K}/\mathfrak{K}^{n+1})
\subset\mathfrak{J}/\mathfrak{J}^{n+1}$ for every $n$, hence, by
passing to the projective limit, $\varphi(\mathfrak{K})\subset
\mathfrak{J}$, and this implies
$f^*(\mathscr{K})\mathscr{O}_{\mathfrak{X}}\subset\mathscr{J}$
(\hyperref[I.10.5.6]{10.5.6, (ii)}).
\end{proof}

\begin{corollary}[10.6.10]
\label{I.10.6.10}
Let $\mathfrak{X}$, $\mathfrak{Y}$ be two locally Noetherian formal
preschemes, and $\mathscr{T}$ the largest ideal of definition
of $\mathfrak{X}$ (\hyperref[I.10.5.4]{10.5.4}).
\begin{enumerate}
  \item[\rm{(i)}] For every ideal of definition $\mathscr{K}$ of
  $\mathfrak{Y}$ and every morphism
  $f:\mathfrak{X}\to\mathfrak{Y}$, we have
  $f^*(\mathscr{K})\mathscr{O}_{\mathfrak{X}}\subset\mathscr{T}$.
  \item[\rm{(ii)}] There is a canonical bijective correspondence between
  $\operatorname{Hom}(\mathfrak{X},\mathfrak{Y})$ and the set of
  sequences $(f_n)$ of morphisms making the diagrams
  (\hyperref[I.10.6.7.1]{10.6.7.1}) commutative, where
  $X_n=(\mathfrak{X},
  \mathscr{O}_{\mathfrak{X}}/\mathscr{T}^{n+1})$,
  $Y_n=(\mathfrak{Y},
  \mathscr{O}_{\mathfrak{Y}}/\mathscr{K}^{n+1})$.
\end{enumerate}
\end{corollary}

\begin{proof}
(ii) follows immediately from (i) and from
(\hyperref[I.10.6.9]{10.6.9}). To prove (i), we may restrict
to the case where $\mathfrak{X}=\operatorname{Spf}(A)$,
$\mathfrak{Y}=\operatorname{Spf}(B)$, $A$ and $B$ being Noetherian,
$\mathscr{T}=\mathfrak{T}^\Delta$,
$\mathscr{K}=\mathfrak{K}^\Delta$, where $\mathfrak{T}$ is the largest
ideal of definition of $A$ and $\mathfrak{K}$ an ideal of definition of
$B$. Let $f=({}^a\varphi,\widetilde{\varphi})$, where
$\varphi:B\to A$ is a continuous homomorphism; since the elements of
$\mathfrak{K}$ are topologically nilpotent
(\hyperref[0.7.1.4]{0, 7.1.4, (ii)}), the same is true of those of
$\varphi(\mathfrak{K})$, hence
$\varphi(\mathfrak{K})\subset\mathfrak{T}$, since $\mathfrak{T}$ is
the set of topologically nilpotent elements of $A$
(\hyperref[0.7.1.6]{0, 7.1.6}); whence the conclusion by virtue of
(\hyperref[I.10.5.6]{10.5.6, (ii)}).
\end{proof}

\begin{corollary}[10.6.11]
\label{I.10.6.11}
Let $\mathfrak{S}$, $\mathfrak{X}$, $\mathfrak{Y}$ be three locally
Noetherian formal preschemes, and
$f:\mathfrak{X}\to\mathfrak{S}$,
$g:\mathfrak{Y}\to\mathfrak{S}$ morphisms making
$\mathfrak{X}$ and $\mathfrak{Y}$ into formal $\mathfrak{S}$-preschemes.
Let $\mathscr{J}$ (resp. $\mathscr{K}$, $\mathscr{L}$) be an
ideal of definition of $\mathfrak{S}$ (resp. $\mathfrak{X}$,
$\mathfrak{Y}$), and suppose that
$f^*(\mathscr{J})\mathscr{O}_{\mathfrak{X}}\subset\mathscr{K}$,
$g^*(\mathscr{J})\mathscr{O}_{\mathfrak{Y}}=\mathscr{L}$; set
$S_n=(\mathfrak{S},
\mathscr{O}_{\mathfrak{S}}/\mathscr{J}^{n+1})$,
$X_n=(\mathfrak{X},
\mathscr{O}_{\mathfrak{X}}/\mathscr{K}^{n+1})$,
$Y_n=(\mathfrak{Y},
\mathscr{O}_{\mathfrak{Y}}/\mathscr{L}^{n+1})$. There is then a
canonical bijective

\oldpage[I]{193}

correspondence between
$\operatorname{Hom}_{\mathfrak{S}}(\mathfrak{X},\mathfrak{Y})$ and
the set of sequences $(u_n)$ of $S_n$-morphisms
$u_n:X_n\to Y_n$ making the diagrams
(\hyperref[I.10.6.7.1]{10.6.7.1}) commutative.
\end{corollary}

\begin{proof}
For every $\mathfrak{S}$-morphism
$u:\mathfrak{X}\to\mathfrak{Y}$, we have $f=g\circ u$ by definition, hence
\[
  u^*(\mathscr{L})\mathscr{O}_{\mathfrak{X}}
  =u^*(g^*(\mathscr{J})\mathscr{O}_{\mathfrak{Y}})
    \mathscr{O}_{\mathfrak{X}}
  =f^*(\mathscr{J})\mathscr{O}_{\mathfrak{X}}
  \subset\mathscr{K}
\]
the corollary therefore follows from
(\hyperref[I.10.6.9]{10.6.9}).
\end{proof}

Note that, for $m\leq n$, the datum of a morphism
$f_n:X_n\to Y_n$ determines one and only one morphism
$f_m:X_m\to Y_m$ making the diagram
(\hyperref[I.10.6.7.1]{10.6.7.1}) commutative, as one sees
immediately by reducing to the affine case; we have thus defined a map
$\varphi_{mn}:\operatorname{Hom}_{S_n}(X_n,Y_n)
\to\operatorname{Hom}_{S_m}(X_m,Y_m)$, and the
$\operatorname{Hom}_{S_n}(X_n,Y_n)$ form a
\emph{projective system of sets} for the $\varphi_{mn}$;
(\hyperref[I.10.6.11]{10.6.11}) can also be stated by saying that
there exists a canonical bijection
\[
  \operatorname{Hom}_{\mathfrak{S}}(\mathfrak{X},\mathfrak{Y})
  \xrightarrow{\sim}
  \varprojlim_n\operatorname{Hom}_{S_n}(X_n,Y_n).
\]

\subsection{Product of formal preschemes.}
\label{subsection:I.10.7}

\begin{env}[10.7.1]
\label{I.10.7.1}
Let $\mathfrak{S}$ be a formal prescheme; since the formal
$\mathfrak{S}$-preschemes form a category, we can define the notion of
\emph{product} of formal $\mathfrak{S}$-preschemes.
\end{env}

\begin{proposition}[10.7.2]
\label{I.10.7.2}
Let $\mathfrak{X}=\operatorname{Spf}(B)$,
$\mathfrak{Y}=\operatorname{Spf}(C)$ be two affine formal schemes
over an affine formal scheme
$\mathfrak{S}=\operatorname{Spf}(A)$. Let
$\mathfrak{Z}=\operatorname{Spf}(B\widehat{\otimes}_A C)$, and let
$p_1$, $p_2$ be the $\mathfrak{S}$-morphisms corresponding
(\hyperref[I.10.2.2]{10.2.2}) to the canonical (continuous)
$A$-homomorphisms $\rho$ and $\sigma$ from $B$ and $C$ to
$B\widehat{\otimes}_A C$; then $(\mathfrak{Z},p_1,p_2)$ is a product
of the affine formal $\mathfrak{S}$-schemes $\mathfrak{X}$ and
$\mathfrak{Y}$.
\end{proposition}

\begin{proof}
By virtue of (\hyperref[I.10.4.6]{10.4.6}), everything comes down to
verifying that if, to every continuous $A$-homomorphism
$\varphi:B\widehat{\otimes}_A C\to D$, where $D$ is an admissible ring
which is a topological $A$-algebra, we associate the pair
$(\varphi\circ\rho,\varphi\circ\sigma)$, we define a bijection
\[
  \operatorname{Hom}_A(B\widehat{\otimes}_A C,D)
  \xrightarrow{\sim}
  \operatorname{Hom}_A(B,D)\times\operatorname{Hom}_A(C,D)
\]
which is none other than the universal property of the completed
tensor product (\hyperref[0.7.7.6]{0, 7.7.6}).
\end{proof}

\begin{proposition}[10.7.3]
\label{I.10.7.3}
Given two formal $\mathfrak{S}$-preschemes
$\mathfrak{X}$, $\mathfrak{Y}$, the product
$\mathfrak{X}\times_{\mathfrak{S}}\mathfrak{Y}$ exists.
\end{proposition}

\begin{proof}
The proof is identical to that of
(\hyperref[I.3.2.6]{3.2.6}), replacing in it the affine schemes
(resp. the affine open subsets) by the affine formal schemes (resp. the
affine formal open subsets), and Proposition
(\hyperref[I.3.2.2]{3.2.2}) by
(\hyperref[I.10.7.2]{10.7.2}).
\end{proof}

All the formal properties of the product of preschemes
(\hyperref[I.3.2.7]{3.2.7} and
\hyperref[I.3.2.8]{3.2.8},
\hyperref[I.3.3.1]{3.3.1} to
\hyperref[I.3.3.12]{3.3.12}) are valid without any modification
for the product of formal preschemes.

\begin{env}[10.7.4]
\label{I.10.7.4}
Let $\mathfrak{S}$, $\mathfrak{X}$, $\mathfrak{Y}$ be three formal
preschemes and let $f:\mathfrak{X}\to\mathfrak{S}$,
$g:\mathfrak{Y}\to\mathfrak{S}$ be two morphisms. Suppose that there
exist in $\mathfrak{S}$, $\mathfrak{X}$, $\mathfrak{Y}$ respectively
three fundamental systems of ideals of definition
$(\mathscr{J}_\lambda)$, $(\mathscr{K}_\lambda)$,
$(\mathscr{L}_\lambda)$ respectively, having the same index set
$I$, such that
$f^*(\mathscr{J}_\lambda)\mathscr{O}_{\mathfrak{X}}
\subset\mathscr{K}_\lambda$ and
$g^*(\mathscr{J}_\lambda)\mathscr{O}_{\mathfrak{Y}}
\subset\mathscr{L}_\lambda$ for every $\lambda$. Set
$S_\lambda=(\mathfrak{S},
\mathscr{O}_{\mathfrak{S}}/\mathscr{J}_\lambda)$,
$X_\lambda=(\mathfrak{X},
\mathscr{O}_{\mathfrak{X}}/\mathscr{K}_\lambda)$,
$Y_\lambda=(\mathfrak{Y},
\mathscr{O}_{\mathfrak{Y}}/\mathscr{L}_\lambda)$; for
$\mathscr{J}_\mu\subset\mathscr{J}_\lambda$,
$\mathscr{K}_\mu\subset\mathscr{K}_\lambda$,
$\mathscr{L}_\mu\subset\mathscr{L}_\lambda$, note that
$S_\lambda$ (resp. $X_\lambda$, $Y_\lambda$) is a closed sub-prescheme
of $S_\mu$ (resp. $X_\mu$, $Y_\mu$) having the same

\oldpage[I]{194}

underlying space (\hyperref[I.10.6.1]{10.6.1}). Since
$S_\lambda\to S_\mu$ is a monomorphism of preschemes, we therefore see
first that the products
$X_\lambda\times_{S_\lambda}Y_\lambda$ and
$X_\lambda\times_{S_\mu}Y_\lambda$ are identical
(\hyperref[I.3.2.4]{3.2.4}), and then that
$X_\lambda\times_{S_\mu}Y_\lambda$ is identified with a closed
sub-prescheme of $X_\mu\times_{S_\mu}Y_\mu$ having the same underlying
space (\hyperref[I.4.3.1]{4.3.1}). This being so, the product
$\mathfrak{X}\times_{\mathfrak{S}}\mathfrak{Y}$ is the
\emph{inductive limit} of the ordinary preschemes
$X_\lambda\times_{S_\lambda}Y_\lambda$: indeed, we see as in
(\hyperref[I.10.6.2]{10.6.2}) that we can reduce to the case where
$\mathfrak{S}$, $\mathfrak{X}$ and $\mathfrak{Y}$ are affine formal
schemes. Taking into account (\hyperref[I.10.5.6]{10.5.6, (ii)}) and
the hypothesis on the fundamental systems of ideals of definition of
$\mathfrak{S}$, $\mathfrak{X}$ and $\mathfrak{Y}$, we see immediately
that our assertion follows from the definition of the completed tensor
product of two algebras (\hyperref[0.7.7.1]{0, 7.7.1}).

Moreover, let $\mathfrak{Z}$ be a formal $\mathfrak{S}$-prescheme,
$(\mathscr{M}_\lambda)$ a fundamental system of ideals of definition of
$\mathfrak{Z}$ having $I$ as index set, and
$u:\mathfrak{Z}\to\mathfrak{X}$,
$v:\mathfrak{Z}\to\mathfrak{Y}$ two $\mathfrak{S}$-morphisms such that
$u^*(\mathscr{K}_\lambda)\mathscr{O}_{\mathfrak{Z}}
\subset\mathscr{M}_\lambda$ and
$v^*(\mathscr{L}_\lambda)\mathscr{O}_{\mathfrak{Z}}
\subset\mathscr{M}_\lambda$. If we set
$Z_\lambda=(\mathfrak{Z},
\mathscr{O}_{\mathfrak{Z}}/\mathscr{M}_\lambda)$, and if
$u_\lambda:Z_\lambda\to X_\lambda$ and
$v_\lambda:Z_\lambda\to Y_\lambda$ are the $S_\lambda$-morphisms
corresponding to $u$ and $v$ (\hyperref[I.10.5.6]{10.5.6}), we
verify immediately that $(u,v)_{\mathfrak{S}}$ is the inductive limit
of the $S_\lambda$-morphisms $(u_\lambda,v_\lambda)_{S_\lambda}$.

The considerations of this subsection apply in particular when
$\mathfrak{S}$, $\mathfrak{X}$ and $\mathfrak{Y}$ are locally
Noetherian, by taking as fundamental systems of ideals of definition
the systems formed by the powers of an ideal of definition
(\hyperref[I.10.5.1]{10.5.1}). But note that
$\mathfrak{X}\times_{\mathfrak{S}}\mathfrak{Y}$ is not necessarily
locally Noetherian (see however
(\hyperref[I.10.13.5]{10.13.5})).
\end{env}


\subsection{Formal completion of a prescheme along a closed subset.}
\label{subsection:I.10.8}

\begin{env}[10.8.1]
\label{I.10.8.1}
Let $X$ be a locally Noetherian (ordinary) prescheme, and $X'$ a closed
subset of the underlying space of $X$; denote by $\Phi$ the set of
coherent sheaves of ideals $\mathscr{J}$ in $\mathscr{O}_X$ such that
the support of $\mathscr{O}_X/\mathscr{J}$ is $X'$. The set $\Phi$
is not empty (\hyperref[I.5.2.1]{5.2.1},
\hyperref[I.4.1.4]{4.1.4} and
\hyperref[I.6.1.1]{6.1.1}); we order it by the relation
$\supset$.
\end{env}

\begin{lemma}[10.8.2]
\label{I.10.8.2}
The ordered set $\Phi$ is filtered; if $X$ is Noetherian, then for every
$\mathscr{J}_0\in\Phi$, the set of powers
$\mathscr{J}_0^n$ ($n>0$) is cofinal in $\Phi$.
\end{lemma}

\begin{proof}
Indeed, if $\mathscr{J}_1$ and $\mathscr{J}_2$ belong to $\Phi$, and
if we set $\mathscr{J}=\mathscr{J}_1\cap\mathscr{J}_2$, then
$\mathscr{J}$ is coherent since $\mathscr{O}_X$ is coherent
(\hyperref[I.6.1.1]{6.1.1} and
\hyperref[0.5.3.4]{0, 5.3.4}), and we have
$\mathscr{J}_x=(\mathscr{J}_1)_x\cap(\mathscr{J}_2)_x$ for every
$x\in X$, hence $\mathscr{J}_x=\mathscr{O}_x$ for $x\notin X'$ and
$\mathscr{J}_x\neq\mathscr{O}_x$ for $x\in X'$, which proves that
$\mathscr{J}\in\Phi$. On the other hand, if $X$ is Noetherian, and if
$\mathscr{J}_0$ and $\mathscr{J}$ belong to $\Phi$, there exists an
integer $n>0$ such that
$\mathscr{J}_0^n(\mathscr{O}_X/\mathscr{J})=0$
(\hyperref[I.9.3.4]{9.3.4}), which means that
$\mathscr{J}_0^n\subset\mathscr{J}$.
\end{proof}

\begin{env}[10.8.3]
\label{I.10.8.3}
Now let $\mathscr{F}$ be a coherent $\mathscr{O}_X$-module; for
every $\mathscr{J}\in\Phi$,
$\mathscr{F}\otimes_{\mathscr{O}_X}(\mathscr{O}_X/\mathscr{J})$ is a
coherent $\mathscr{O}_X$-module
(\hyperref[I.9.1.1]{9.1.1}), with support contained in $X'$, which we
will most often identify with its restriction to $X'$. As
$\mathscr{J}$ runs through $\Phi$, these sheaves form a
\emph{projective system} of sheaves of abelian groups.
\end{env}

\begin{definition}[10.8.4]
\label{I.10.8.4}
Given a closed subset $X'$ of a locally Noetherian prescheme
$X$ and a coherent $\mathscr{O}_X$-module $\mathscr{F}$, we call
\emph{completion of $\mathscr{F}$ along $X'$}, and denote by
$\mathscr{F}_{/X'}$ or by $\widehat{\mathscr{F}}$ (when no confusion
is possible), the restriction to $X'$ of the sheaf

\oldpage[I]{195}

$\varprojlim_{\Phi}
(\mathscr{F}\otimes_{\mathscr{O}_X}
(\mathscr{O}_X/\mathscr{J}))$; we say that its sections over $X'$ are
the \emph{formal sections of $\mathscr{F}$ along $X'$}.
\end{definition}

It is immediate that for every open subset $U\subset X$, we have
$(\mathscr{F}|U)_{/(U\cap X')}=(\mathscr{F}_{/X'})|(U\cap X')$.

By passing to the projective limit, it is clear that
$(\mathscr{O}_X)_{/X'}$ is a sheaf of rings, and that
$\mathscr{F}_{/X'}$ can be regarded as an
$(\mathscr{O}_X)_{/X'}$-module. Moreover, since there exists a basis of
the topology of $X'$ formed by quasi-compact open subsets, we can regard
$(\mathscr{O}_X)_{/X'}$ (resp. $\mathscr{F}_{/X'}$) as a
\emph{sheaf of topological rings} (resp. of
\emph{topological groups}) that is the projective limit of the
\emph{pseudo-discrete} sheaves of rings (resp. groups)
$\mathscr{O}_X/\mathscr{J}$ (resp.
$\mathscr{F}\otimes_{\mathscr{O}_X}
(\mathscr{O}_X/\mathscr{J})=\mathscr{F}/\mathscr{J}\mathscr{F}$), and, by
passing to the projective limit, $\mathscr{F}_{/X'}$ then becomes a
\emph{topological $(\mathscr{O}_X)_{/X'}$-module}
(\hyperref[0.3.8.1]{0, 3.8.1} and
\hyperref[0.3.8.2]{3.8.2}); recall that for every
\emph{quasi-compact} open subset $U\subset X$,
$\Gamma(U\cap X',(\mathscr{O}_X)_{/X'})$ (resp.
$\Gamma(U\cap X',\mathscr{F}_{/X'})$) is then the projective limit of
the discrete rings (resp. groups) $\Gamma(U,\mathscr{O}_X/\mathscr{J})$
(resp. $\Gamma(U,\mathscr{F}/\mathscr{J}\mathscr{F})$).

If now $u:\mathscr{F}\to\mathscr{G}$ is a homomorphism of
$\mathscr{O}_X$-modules, we canonically deduce from it homomorphisms
$u_{\mathscr{J}}:
\mathscr{F}\otimes_{\mathscr{O}_X}(\mathscr{O}_X/\mathscr{J})
\to
\mathscr{G}\otimes_{\mathscr{O}_X}(\mathscr{O}_X/\mathscr{J})$
for every $\mathscr{J}\in\Phi$, and these homomorphisms form a
projective system. By passing to the projective limit and restricting
to $X'$, they therefore give a continuous
$(\mathscr{O}_X)_{/X'}$-homomorphism
$\mathscr{F}_{/X'}\to\mathscr{G}_{/X'}$, denoted $u_{/X'}$ or $\widehat u$,
and called the \emph{completion of the homomorphism $u$ along $X'$}. It
is clear that if $v:\mathscr{G}\to\mathscr{H}$ is a second homomorphism
of $\mathscr{O}_X$-modules, we have
$(v\circ u)_{/X'}=(v_{/X'})\circ(u_{/X'})$, so
$\mathscr{F}_{/X'}$ is a \emph{covariant additive functor} in
$\mathscr{F}$, from the category of coherent $\mathscr{O}_X$-modules to
the category of topological $(\mathscr{O}_X)_{/X'}$-modules.

\begin{proposition}[10.8.5]
\label{I.10.8.5}
The support of $(\mathscr{O}_X)_{/X'}$ is $X'$; the topologically
ringed space $(X',(\mathscr{O}_X)_{/X'})$ is a locally Noetherian
formal prescheme, and if $\mathscr{J}\in\Phi$, $\mathscr{J}_{/X'}$ is an
ideal of definition of this formal prescheme. If
$X=\operatorname{Spec}(A)$ is an affine scheme with Noetherian ring,
$\mathscr{J}=\widetilde{\mathfrak{J}}$, where $\mathfrak{J}$ is an ideal
of $A$, and $X'=V(\mathfrak{J})$, then
$(X',(\mathscr{O}_X)_{/X'})$ is canonically identified with
$\operatorname{Spf}(\widehat A)$, where $\widehat A$ is the separated
completion of $A$ for the $\mathfrak{J}$-preadic topology.
\end{proposition}

\begin{proof}
We may evidently restrict ourselves to proving the last assertion. We
know (\hyperref[0.7.3.3]{0, 7.3.3}) that the separated completion
$\widehat{\mathfrak{J}}$ of $\mathfrak{J}$ for the
$\mathfrak{J}$-preadic topology is identified with the ideal
$\mathfrak{J}\widehat A$ of $\widehat A$, and that $\widehat A$ is a
Noetherian $\widehat{\mathfrak{J}}$-adic ring such that
$\widehat A/\widehat{\mathfrak{J}}^n=A/\mathfrak{J}^n$
(\hyperref[0.7.2.6]{0, 7.2.6}). This last relation shows that the
open prime ideals of $\widehat A$ are the ideals
$\widehat{\mathfrak p}=\mathfrak p\widehat A$, where $\mathfrak p$ is a
prime ideal of $A$ containing $\mathfrak{J}$, and that we have
$\widehat{\mathfrak p}\cap A=\mathfrak p$, whence
$\operatorname{Spf}(\widehat A)=X'$. Since
$\mathscr{O}_X/\mathscr{J}^n=(A/\mathfrak{J}^n)^{\sim}$, the proposition
follows immediately from the definitions.
\end{proof}

We say that the formal prescheme thus defined is the
\emph{completion of $X$ along $X'$}, and we denote it by $X_{/X'}$, or by
$\widehat X$ if no confusion is to be feared. When we take
$X'=X$, we can take $\mathscr{J}=0$, and we therefore have $X_{/X}=X$.

It is clear that if $U$ is a sub-prescheme induced on an open subset of
$X$, then $U_{/(U\cap X')}$ is canonically identified with the formal
sub-prescheme induced by $X_{/X'}$ on the open subset $U\cap X'$ of $X'$.

\begin{corollary}[10.8.6]
\label{I.10.8.6}
The (ordinary) prescheme $\widehat X_{\mathrm{red}}$ is the unique
reduced sub-prescheme of $X$ having $X'$ as underlying space
(\hyperref[I.5.2.1]{5.2.1}). In order that $\widehat X$ be Noetherian,
it is necessary and sufficient that $\widehat X_{\mathrm{red}}$ be so,
and it suffices that $X$ be so.
\end{corollary}

\oldpage[I]{196}

The determination of $\widehat X_{\mathrm{red}}$ being local
(\hyperref[I.10.5.4]{10.5.4}), we may again suppose that $X$ is an
affine scheme with Noetherian ring; with the notation of
(\hyperref[I.10.8.5]{10.8.5}), the ideal $\mathfrak{T}$ of
topologically nilpotent elements of $\widehat A$ is the inverse image,
under the canonical map
$\widehat A\to\widehat A/\widehat{\mathfrak{J}}=A/\mathfrak{J}$, of the
nilradical of $A/\mathfrak{J}$
(\hyperref[0.7.1.3]{0, 7.1.3}), so $\widehat A/\mathfrak{T}$ is
isomorphic to the quotient of $A/\mathfrak{J}$ by its nilradical. The
first assertion therefore follows from (\hyperref[I.10.5.4]{10.5.4})
and (\hyperref[I.5.1.1]{5.1.1}). If $\widehat X_{\mathrm{red}}$ is
Noetherian, its underlying space $X'$ is also Noetherian, hence the
$X'_n=\operatorname{Spec}(\mathscr{O}_X/\mathscr{J}^n)$ are Noetherian
(\hyperref[I.6.1.2]{6.1.2}), and the same holds for $\widehat X$
(\hyperref[I.10.6.4]{10.6.4}); the converse is immediate, by virtue
of (\hyperref[I.6.1.2]{6.1.2}).

\begin{env}[10.8.7]
\label{I.10.8.7}
The canonical homomorphisms
$\mathscr{O}_X\to\mathscr{O}_X/\mathscr{J}$ (for
$\mathscr{J}\in\Phi$) form a projective system and therefore give, by
passing to the projective limit, a homomorphism of sheaves of rings
$\theta:\mathscr{O}_X\to\psi_*((\mathscr{O}_X)_{/X'})
=\varprojlim_\Phi(\mathscr{O}_X/\mathscr{J})$, where $\psi$ denotes
the canonical injection $X'\to X$ of the underlying spaces. We denote
by $i$ (or $i_X$) the (so-called canonical) morphism
\[
  (\psi,\theta):X_{/X'}\to X
\]
of ringed spaces.

By tensoring, for every coherent $\mathscr{O}_X$-module
$\mathscr{F}$, the canonical homomorphisms
$\mathscr{O}_X\to\mathscr{O}_X/\mathscr{J}$ give homomorphisms
$\mathscr{F}\to
\mathscr{F}\otimes_{\mathscr{O}_X}(\mathscr{O}_X/\mathscr{J})$ of
$\mathscr{O}_X$-modules which again form a projective system, and
therefore give, by passing to the projective limit, a canonical
functorial homomorphism
$\gamma:\mathscr{F}\to\psi_*(\mathscr{F}_{/X'})$ of
$\mathscr{O}_X$-modules.
\end{env}

\begin{proposition}[10.8.8]
\label{I.10.8.8}
\medskip\noindent
\begin{enumerate}
  \item[\rm{(i)}] The functor $\mathscr{F}_{/X'}$ (in $\mathscr{F}$) is
    exact.
  \item[\rm{(ii)}] The functorial homomorphism
    $\gamma^\sharp:i^*(\mathscr{F})\to\mathscr{F}_{/X'}$ of
    $(\mathscr{O}_X)_{/X'}$-modules is an isomorphism.
\end{enumerate}
\end{proposition}

\medskip\noindent
\begin{enumerate}
  \item[\rm{(i)}] It suffices to prove that if
    $0\to\mathscr{F}'\to\mathscr{F}\to\mathscr{F}''\to0$ is an exact
    sequence of coherent $\mathscr{O}_X$-modules, and $U$ an affine
    open subset of $X$ with Noetherian ring $A$, the sequence
    \[
      0\to\Gamma(U\cap X',\mathscr{F}'_{/X'})
      \to\Gamma(U\cap X',\mathscr{F}_{/X'})
      \to\Gamma(U\cap X',\mathscr{F}''_{/X'})\to0
    \]
    is exact. We then have $\mathscr{F}|U=\widetilde M$,
    $\mathscr{F}'|U=\widetilde{M'}$, $\mathscr{F}''|U=\widetilde{M''}$,
    where $M$, $M'$, $M''$ are three $A$-modules of finite type such that
    the sequence $0\to M'\to M\to M''\to0$ is exact
    (\hyperref[I.1.5.1]{1.5.1} and
    \hyperref[I.1.3.11]{1.3.11}); let $\mathscr{J}\in\Phi$ and let
    $\mathfrak{J}$ be an ideal of $A$ such that
    $\mathscr{J}|U=\widetilde{\mathfrak{J}}$. We then have
    \[
      \Gamma(U\cap X',
      \mathscr{F}\otimes_{\mathscr{O}_X}(\mathscr{O}_X/\mathscr{J}^n))
      =M\otimes_A(A/\mathfrak{J}^n)
    \]
    (\hyperref[I.1.3.12]{1.3.12}); hence, by definition of the
    projective limit, we have
    \[
      \Gamma(U\cap X',\mathscr{F}_{/X'})
      =\varprojlim_n(M\otimes_A(A/\mathfrak{J}^n))=\widehat M
    \]
    the separated completion of $M$ for the $\mathfrak{J}$-preadic
    topology, and similarly
    \[
      \Gamma(U\cap X',\mathscr{F}'_{/X'})=\widehat{M'},\qquad
      \Gamma(U\cap X',\mathscr{F}''_{/X'})=\widehat{M''}~;
    \]
    our assertion then follows from the fact that, when $A$ is
    Noetherian, the functor $\widehat M$ in $M$ is exact on the category
    of $A$-modules of finite type (\hyperref[0.7.3.3]{0, 7.3.3}).

\oldpage[I]{197}
  \item[\rm{(ii)}] The question being local, we may suppose that we have
    an exact sequence
    $\mathscr{O}_X^m\to\mathscr{O}_X^n\to\mathscr{F}\to0$
    (\hyperref[0.5.3.2]{0, 5.3.2}); since $\gamma^\sharp$ is
    functorial, and the functors $i^*(\mathscr{F})$ and
    $\mathscr{F}_{/X'}$ are right exact (by (i) and
    (\hyperref[0.4.3.1]{0, 4.3.1})), we have the commutative diagram
    \[
      \xymatrix{
        i^*(\mathscr{O}_X^m)\ar[r]\ar[d]_{\gamma^\sharp} &
        i^*(\mathscr{O}_X^n)\ar[r]\ar[d]_{\gamma^\sharp} &
        i^*(\mathscr{F})\ar[r]\ar[d]_{\gamma^\sharp} &
        0\\
        (\mathscr{O}_X^m)_{/X'}\ar[r] &
        (\mathscr{O}_X^n)_{/X'}\ar[r] &
        \mathscr{F}_{/X'}\ar[r] &
        0
      }
      \tag{10.8.8.1}\label{I.10.8.8.1}
    \]
    whose rows are exact. Moreover, the two functors
    $i^*(\mathscr{F})$ and $\mathscr{F}_{/X'}$ commute with finite direct
    sums (\hyperref[0.3.2.6]{0, 3.2.6} and
    \hyperref[0.4.3.2]{4.3.2}), and we are therefore reduced to proving
    our assertion for $\mathscr{F}=\mathscr{O}_X$. We then have
    $i^*(\mathscr{O}_X)=(\mathscr{O}_X)_{/X'}=
    \mathscr{O}_{\widehat X}$
    (\hyperref[0.4.3.4]{0, 4.3.4}), and $\gamma^\sharp$ is a
    homomorphism of $\mathscr{O}_{\widehat X}$-modules; it therefore
    suffices to verify that $\gamma^\sharp$ transforms the unit section of
    $\mathscr{O}_{\widehat X}$ over an open subset of $X'$ into itself,
    which is immediate and therefore shows that in this case
    $\gamma^\sharp$ is the identity.
\end{enumerate}

\begin{corollary}[10.8.9]
\label{I.10.8.9}
The morphism of ringed spaces $i:X_{/X'}\to X$ is flat.
\end{corollary}

Indeed, this follows from (\hyperref[0.6.7.3]{0, 6.7.3}) and from
(\hyperref[I.10.8.8]{10.8.8, (i)}).

\begin{corollary}[10.8.10]
\label{I.10.8.10}
If $\mathscr{F}$ and $\mathscr{G}$ are coherent
$\mathscr{O}_X$-modules, there exist canonical isomorphisms, functorial
(in $\mathscr{F}$ and $\mathscr{G}$),
\[
  (\mathscr{F}_{/X'})\otimes_{(\mathscr{O}_X)_{/X'}}
  (\mathscr{G}_{/X'})
  \xrightarrow{\sim}
  (\mathscr{F}\otimes_{\mathscr{O}_X}\mathscr{G})_{/X'}
  \tag{10.8.10.1}\label{I.10.8.10.1}
\]
\[
  (\mathscr{H}\!om_{\mathscr{O}_X}(\mathscr{F},\mathscr{G}))_{/X'}
  \xrightarrow{\sim}
  \mathscr{H}\!om_{(\mathscr{O}_X)_{/X'}}
  (\mathscr{F}_{/X'},\mathscr{G}_{/X'})
  \tag{10.8.10.2}\label{I.10.8.10.2}
\]
\end{corollary}

This follows from the canonical identification of $i^*(\mathscr{F})$ and
$\mathscr{F}_{/X'}$; the existence of the first isomorphism is then a
result valid for all morphisms of ringed spaces
(\hyperref[0.4.3.3.1]{0, 4.3.3.1}), and that of the second is a result
valid for all flat morphisms
(\hyperref[0.6.7.6]{0, 6.7.6}), and therefore follows from
(\hyperref[I.10.8.9]{10.8.9}).

\begin{proposition}[10.8.11]
\label{I.10.8.11}
For every coherent $\mathscr{O}_X$-module $\mathscr{F}$, the kernel of
the canonical homomorphism
$\Gamma(X,\mathscr{F})\to\Gamma(X',\mathscr{F}_{/X'})$ deduced from
$\mathscr{F}\to\mathscr{F}_{/X'}$ consists of the sections that are zero
in a neighbourhood of $X'$.
\end{proposition}

It follows from the definition of $\mathscr{F}_{/X'}$ that the canonical
image of such a section is zero. Conversely, if
$s\in\Gamma(X,\mathscr{F})$ has zero image in
$\Gamma(X',\mathscr{F}_{/X'})$, it suffices to see that every $x\in X'$
has a neighbourhood in $X$ in which $s$ is zero, and we can therefore
reduce to the case where $X=\operatorname{Spec}(A)$ is affine, $A$
Noetherian, $X'=V(\mathfrak{J})$, where $\mathfrak{J}$ is an ideal of
$A$, and $\mathscr{F}=\widetilde M$, where $M$ is an $A$-module of finite
type. Then $\Gamma(X',\mathscr{F}_{/X'})$ is the separated completion
$\widehat M$ of $M$ for the $\mathfrak{J}$-preadic topology, and the
homomorphism $\Gamma(X,\mathscr{F})\to\Gamma(X',\mathscr{F}_{/X'})$ is
the canonical homomorphism $M\to\widehat M$. We know
(\hyperref[0.7.3.7]{0, 7.3.7}) that the kernel of this homomorphism is
the set of $z\in M$ annihilated by an element of $1+\mathfrak{J}$. We
therefore have $(1+f)s=0$ for some $f\in\mathfrak{J}$; for every
$x\in X'$ we deduce from this that $(1_x+f_x)s_x=0$, and since $1_x+f_x$
is invertible in $\mathscr{O}_x$ ($\mathfrak{J}_x\mathscr{O}_x$ being
contained in the maximal ideal of $\mathscr{O}_x$), we have $s_x=0$,
which proves the proposition.

\begin{corollary}[10.8.12]
\label{I.10.8.12}
The support of $\mathscr{F}_{/X'}$ is equal to
$\operatorname{Supp}(\mathscr{F})\cap X'$.
\end{corollary}

It is clear that $\mathscr{F}_{/X'}$ is an
$(\mathscr{O}_X)_{/X'}$-module of finite type
(\hyperref[I.10.8.8]{10.8.8, (ii)}) and
(\hyperref[0.5.2.4]{0, 5.2.4}),

\oldpage[I]{198}

so its support is closed (\hyperref[0.5.2.2]{0, 5.2.2}) and
evidently contained in $\operatorname{Supp}(\mathscr{F})\cap X'$. To
show that it is equal to this latter set, we are immediately reduced to
proving that the relation $\Gamma(X',\mathscr{F}_{/X'})=0$ implies
$\operatorname{Supp}(\mathscr{F})\cap X'=\varnothing$; now this follows
from (\hyperref[I.10.8.11]{10.8.11}) and from
(\hyperref[I.1.4.1]{1.4.1}).

\begin{corollary}[10.8.13]
\label{I.10.8.13}
Let $u:\mathscr{F}\to\mathscr{G}$ be a homomorphism of coherent
$\mathscr{O}_X$-modules. In order that
$u_{/X'}:\mathscr{F}_{/X'}\to\mathscr{G}_{/X'}$ be zero, it is necessary
and sufficient that $u$ be zero in a neighbourhood of $X'$.
\end{corollary}

Indeed, by (\hyperref[I.10.8.8]{10.8.8, (ii)}), $u_{/X'}$
is identified with $i^*(u)$, so if we regard $u$ as a section over
$X$ of the sheaf
$\mathscr{H}=\mathscr{H}\!om_{\mathscr{O}_X}(\mathscr{F},\mathscr{G})$,
then $u_{/X'}$ is the section of
$i^*(\mathscr{H})=\mathscr{H}_{/X'}$ over $X'$ that corresponds to it
canonically ((\hyperref[I.10.8.10.2]{10.8.10.2}) and
(\hyperref[0.4.4.6]{0, 4.4.6})). It therefore suffices to apply
(\hyperref[I.10.8.11]{10.8.11}) to the coherent $\mathscr{O}_X$-module
$\mathscr{H}$.

\begin{corollary}[10.8.14]
\label{I.10.8.14}
Let $u:\mathscr{F}\to\mathscr{G}$ be a homomorphism of coherent
$\mathscr{O}_X$-modules. In order that $u_{/X'}$ be a monomorphism
(resp. an epimorphism), it is necessary and sufficient that $u$ be a
monomorphism (resp. an epimorphism) in a neighbourhood of $X'$.
\end{corollary}

Let $\mathscr{P}$ and $\mathscr{N}$ be the cokernel and the kernel of
$u$, so that we have the exact sequence
\[
  0\to\mathscr{N}\xrightarrow{v}\mathscr{F}\xrightarrow{u}
  \mathscr{G}\xrightarrow{w}\mathscr{P}\to0,
\]
whence (\hyperref[I.10.8.8]{10.8.8, (i)}) the exact sequence
\[
  0\to\mathscr{N}_{/X'}\xrightarrow{v_{/X'}}
  \mathscr{F}_{/X'}\xrightarrow{u_{/X'}}
  \mathscr{G}_{/X'}\xrightarrow{w_{/X'}}
  \mathscr{P}_{/X'}\to0.
\]
If $u_{/X'}$ is a monomorphism (resp. an epimorphism), we have
$v_{/X'}=0$ (resp. $w_{/X'}=0$), so there is a neighbourhood of $X'$ in
which $v=0$ (resp. $w=0$), by virtue of
(\hyperref[I.10.8.13]{10.8.13}).

\subsection{Extension of a morphism to the completions.}
\label{subsection:I.10.9}

\begin{env}[10.9.1]
\label{I.10.9.1}
Let $X$, $Y$ be two locally Noetherian (ordinary) preschemes,
$f:X\to Y$ a morphism, and $X'$ (resp. $Y'$) a closed subset of the
underlying space $X$ (resp. $Y$), such that $f(X')\subset Y'$. Let
$\mathscr{J}$ (resp. $\mathscr{K}$) be a sheaf of ideals of
$\mathscr{O}_X$ (resp. $\mathscr{O}_Y$) such that the support of
$\mathscr{O}_X/\mathscr{J}$ (resp. $\mathscr{O}_Y/\mathscr{K}$) is
$X'$ (resp. $Y'$) and that
$f^*(\mathscr{K})\mathscr{O}_X\subset\mathscr{J}$; note that such
sheaves of ideals always exist, since we can for example take for
$\mathscr{J}$ the largest sheaf of ideals of $\mathscr{O}_X$
defining a sub-prescheme of $X$ having $X'$ as underlying space
(\hyperref[I.5.2.1]{5.2.1}), and the hypothesis $f(X')\subset Y'$ then
implies $f^*(\mathscr{K})\mathscr{O}_X\subset\mathscr{J}$
(\hyperref[I.5.2.4]{5.2.4}). We therefore have, for every integer
$n>0$, $f^*(\mathscr{K}^n)\mathscr{O}_X\subset\mathscr{J}^n$
(\hyperref[0.4.3.5]{0, 4.3.5}); consequently
(\hyperref[I.4.4.6]{4.4.6}), if we set
\[
  X'_n=(X',\mathscr{O}_X/\mathscr{J}^{n+1}),\qquad
  Y'_n=(Y',\mathscr{O}_Y/\mathscr{K}^{n+1}),
\]
we deduce from $f$ a morphism $f_n:X'_n\to Y'_n$, and it is immediate
that the $f_n$ form an inductive system. We denote its inductive limit
(\hyperref[I.10.6.8]{10.6.8}) by
$\widehat f:X_{/X'}\to Y_{/Y'}$, and we will say (by abuse of language)
that $\widehat f$ is the \emph{extension of $f$ to the completions of
$X$ and $Y$ along $X'$ and $Y'$}. It is immediate to verify that this
morphism does not depend on the choice of the sheaves of ideals
$\mathscr{J}$, $\mathscr{K}$ satisfying the above conditions. Indeed, it
suffices to see this when $X$ and $Y$ are Noetherian affine schemes with
rings $A$, $B$; then
$\mathscr{J}=\widetilde{\mathfrak{J}}$,
$\mathscr{K}=\widetilde{\mathfrak{K}}$, where $\mathfrak{J}$ (resp.
$\mathfrak{K}$) is an ideal of $A$ (resp. $B$), and $f$ corresponds to a
ring homomorphism $\varphi:B\to A$ such that
$\varphi(\mathfrak{K})\subset\mathfrak{J}$
(\hyperref[I.4.4.6]{4.4.6} and
\hyperref[I.1.7.4]{1.7.4}); \emph{$\widehat f$ is then the morphism
which corresponds (\hyperref[I.10.2.2]{10.2.2}) to the continuous
homomorphism $\widehat\varphi:\widehat B\to\widehat A$}, where
$\widehat A$ (resp. $\widehat B$) is the separated completion of $A$
(resp. $B$) for the $\mathfrak{J}$-preadic (resp. $\mathfrak{K}$-preadic)
topology (\hyperref[I.10.6.8]{10.6.8}); and we know that if we
replace $\mathscr{J}$ by another

\oldpage[I]{199}

sheaf of ideals
$\mathscr{J}'=\widetilde{\mathfrak{J}'}$ such that the support of
$\mathscr{O}_X/\mathscr{J}'$ is still $X'$, the
$\mathfrak{J}$-preadic and $\mathfrak{J}'$-preadic topologies on $A$
are the same (\hyperref[I.10.8.2]{10.8.2}).

Note that, according to this definition, the continuous map
$X'\to Y'$ of the underlying spaces of $X_{/X'}$ and $Y_{/Y'}$ which
corresponds to $\widehat f$ is none other than the restriction of $f$ to
$X'$.
\end{env}

\begin{env}[10.9.2]
\label{I.10.9.2}
It follows immediately from the preceding definition that the diagram
of morphisms of ringed spaces
\[ \xymatrix{
    \widehat X\ar[r]^{\widehat f}\ar[d]_{i_X} &
    \widehat Y\ar[d]^{i_Y}\\
    X\ar[r]_f &
    Y
  }
\tag{10.9.2.1}\label{I.10.9.2.1}
\]
is commutative, the vertical arrows being the canonical morphisms
(\hyperref[I.10.8.7]{10.8.7}).
\end{env}

\begin{env}[10.9.3]
\label{I.10.9.3}
Let $Z$ be a third prescheme, $g:Y\to Z$ a morphism, and $Z'$ a
closed subset of $Z$ such that $g(Y')\subset Z'$. If $\widehat g$
denotes the completion along $Y'$ and $Z'$ of the morphism $g$, it
follows immediately from (\hyperref[I.10.9.1]{10.9.1}) that
$\widehat{(g\circ f)}=\widehat g\circ\widehat f$.
\end{env}

\begin{proposition}[10.9.4]
\label{I.10.9.4}
Let $X$, $Y$ be two locally Noetherian $S$-preschemes, $Y$ being of
finite type over $S$. Let $f$, $g$ be two $S$-morphisms from $X$ to $Y$
such that $f(X')\subset Y'$, $g(X')\subset Y'$. In order that
$\widehat f=\widehat g$, it is necessary and sufficient that $f$ and $g$
coincide in a neighbourhood of $X'$.
\end{proposition}

The condition is evidently sufficient (without any finiteness hypothesis
on $Y$). To see that it is necessary, let us first remark that the
hypothesis $\widehat f=\widehat g$ implies $f(x)=g(x)$ for every
$x\in X'$. On the other hand, the question being local, we may suppose
that $X$ and $Y$ are affine open neighbourhoods of $x$ and of
$y=f(x)=g(x)$ respectively, with Noetherian rings, that $S$ is affine and
that $\Gamma(Y,\mathscr{O}_Y)$ is a $\Gamma(S,\mathscr{O}_S)$-algebra of
finite type (\hyperref[I.6.3.3]{6.3.3}). Then $f$ and $g$ correspond
to two $\Gamma(S,\mathscr{O}_S)$-homomorphisms $\rho$, $\sigma$ from
$\Gamma(Y,\mathscr{O}_Y)$ to $\Gamma(X,\mathscr{O}_X)$
(\hyperref[I.1.7.3]{1.7.3}), and by hypothesis, the extensions by
continuity of these homomorphisms to the separated completion of
$\Gamma(Y,\mathscr{O}_Y)$ are the same. We conclude from
(\hyperref[I.10.8.11]{10.8.11}) that for every section
$s\in\Gamma(Y,\mathscr{O}_Y)$, the sections $\rho(s)$ and $\sigma(s)$
coincide in a neighbourhood of $X'$ (depending on $s$); since
$\Gamma(Y,\mathscr{O}_Y)$ is an algebra of finite type over
$\Gamma(S,\mathscr{O}_S)$, we deduce immediately that there exists a
neighbourhood $V$ of $X'$ such that $\rho(s)$ and $\sigma(s)$ coincide in
$V$ for \emph{every} section $s\in\Gamma(Y,\mathscr{O}_Y)$. If
$h\in\Gamma(X,\mathscr{O}_X)$ is such that $D(h)$ is a neighbourhood of
$x$ contained in $V$, we conclude from the foregoing and from
(\hyperref[I.1.4.1]{1.4.1, d)}) that $f$ and $g$ coincide in $D(h)$.

\begin{proposition}[10.9.5]
\label{I.10.9.5}
Under the hypotheses of (\hyperref[I.10.9.1]{10.9.1}), for every
coherent $\mathscr{O}_Y$-module $\mathscr{G}$, there exists a canonical
functorial isomorphism of $(\mathscr{O}_X)_{/X'}$-modules
\[
  (f^*(\mathscr{G}))_{/X'}\xrightarrow{\sim}
  \widehat f^*(\mathscr{G}_{/Y'})
\]
\end{proposition}

If we canonically identify $(f^*(\mathscr{G}))_{/X'}$ with
$i_X^*(f^*(\mathscr{G}))$ and $\widehat f^*(\mathscr{G}_{/Y'})$ with
$\widehat f^*(i_Y^*(\mathscr{G}))$
(\hyperref[I.10.8.8]{10.8.8}), the proposition follows immediately
from the commutativity of the diagram of
(\hyperref[I.10.9.2]{10.9.2}).

\begin{env}[10.9.6]
\label{I.10.9.6}
Now let $\mathscr{F}$ be a coherent $\mathscr{O}_X$-module and
$\mathscr{G}$ a coherent $\mathscr{O}_Y$-module. If
$u:\mathscr{G}\to\mathscr{F}$ is an $f$-morphism from $\mathscr{G}$ to
$\mathscr{F}$, there corresponds to it an $\mathscr{O}_X$-homomorphism
$u^\sharp:f^*(\mathscr{G})\to\mathscr{F}$, hence by completion a
continuous $(\mathscr{O}_X)_{/X'}$-homomorphism
$(u^\sharp)_{/X'}:(f^*(\mathscr{G}))_{/X'}\to\mathscr{F}_{/X'}$, and by
virtue of (\hyperref[I.10.9.5]{10.9.5}) there exists an
$\widehat f$-morphism $v:\mathscr{G}_{/Y'}\to\mathscr{F}_{/X'}$,
\oldpage[I]{200}

and only one, such that $v^\sharp=(u^\sharp)_{/X'}$. If we regard the
triples $(\mathscr{F},X,X')$ ($\mathscr{F}$ being a coherent
$\mathscr{O}_X$-module and $X'$ a closed subset of $X$) as
a \emph{category}, the morphisms
$(\mathscr{F},X,X')\to(\mathscr{G},Y,Y')$ consisting of a morphism of
preschemes $f:X\to Y$ such that $f(X')\subset Y'$ and an
$f$-morphism $u:\mathscr{G}\to\mathscr{F}$, we can therefore say that
$(X_{/X'},\mathscr{F}_{/X'})$ is a \emph{functor} in
$(\mathscr{F},X,X')$, taking its values in the category of pairs
$(\mathfrak{Z},\mathscr{H})$ formed by a locally Noetherian formal
prescheme $\mathfrak{Z}$ and an
$\mathscr{O}_{\mathfrak{Z}}$-module $\mathscr{H}$, the morphisms of
this latter category consisting of the pairs formed by a morphism
$g$ of formal preschemes and a $g$-morphism.
\end{env}

\begin{proposition}[10.9.7]
\label{I.10.9.7}
Let $S$, $X$, $Y$ be three locally Noetherian preschemes,
$g:X\to S$, $h:Y\to S$ two morphisms, $S'$ a closed subset of $S$,
and $X'$ (resp. $Y'$) a closed subset of $X$ (resp. $Y$) such that
$g(X')\subset S'$ (resp. $h(Y')\subset S'$); let
$Z=X\times_S Y$; suppose $Z$ locally Noetherian, and let
$Z'=p^{-1}(X')\cap q^{-1}(Y')$, where $p$ and $q$ are the projections of
$X\times_S Y$. Under these conditions, the completion $Z_{/Z'}$ is
identified with the product of formal $S_{/S'}$-preschemes
$(X_{/X'})\times_{S_{/S'}}(Y_{/Y'})$, the structural morphisms
being identified with $\widehat g$ and $\widehat h$, and the projections
with $\widehat p$ and $\widehat q$.
\end{proposition}

It is immediate that the question is local for $S$, $X$ and $Y$, and we
are therefore reduced to the case where $S=\operatorname{Spec}(A)$,
$X=\operatorname{Spec}(B)$, $Y=\operatorname{Spec}(C)$,
$S'=V(\mathfrak{J})$, $X'=V(\mathfrak{K})$, $Y'=V(\mathfrak{L})$, where
$\mathfrak{J}$, $\mathfrak{K}$, $\mathfrak{L}$ are three ideals such
that $\varphi(\mathfrak{J})\subset\mathfrak{K}$ and
$\psi(\mathfrak{J})\subset\mathfrak{L}$, denoting by $\varphi$ and
$\psi$ the homomorphisms $A\to B$ and $A\to C$ which correspond to $g$
and $h$. Then we know that $Z=\operatorname{Spec}(B\otimes_A C)$ and that
$Z'=V(\mathfrak{M})$, where $\mathfrak{M}$ is the ideal
$\operatorname{Im}(\mathfrak{K}\otimes_A C)
+\operatorname{Im}(B\otimes_A\mathfrak{L})$. The conclusion follows
(\hyperref[I.10.7.2]{10.7.2}) from the fact that the completed tensor
product $(\widehat B\otimes_{\widehat A}\widehat C)^\wedge$ (where
$\widehat A$, $\widehat B$, $\widehat C$ are respectively the separated
completions of $A$, $B$, $C$ for the $\mathfrak{J}$-, $\mathfrak{K}$-
and $\mathfrak{L}$-preadic topologies) is the separated completion of
the tensor product $B\otimes_A C$ for the $\mathfrak{M}$-preadic
topology (\hyperref[0.7.7.2]{0, 7.7.2}).

Note moreover that if $T$ is a locally Noetherian $S$-prescheme,
$u:T\to X$, $v:T\to Y$ two $S$-morphisms, and $T'$ a closed subset of
$T$ such that $u(T')\subset X'$, $v(T')\subset Y'$, then the extension
to the completions $((u,v)_S)^\wedge$ is identified with
$(\widehat u,\widehat v)_{S_{/S'}}$.

\begin{corollary}[10.9.8]
\label{I.10.9.8}
Let $X$, $Y$ be two locally Noetherian $S$-preschemes such that
$X\times_S Y$ is locally Noetherian; let $S'$ be a closed subset of
$S$, and $X'$ (resp. $Y'$) a closed subset of $X$ (resp. $Y$) whose
image in $S$ is contained in $S'$. For every $S$-morphism
$f:X\to Y$ such that $f(X')\subset Y'$, the graph morphism
$\Gamma_{\widehat f}$ is identified with the extension
$(\Gamma_f)^\wedge$ of the graph morphism of $f$.
\end{corollary}

\begin{corollary}[10.9.9]
\label{I.10.9.9}
Let $X$, $Y$ be two locally Noetherian preschemes, $f:X\to Y$ a
morphism, $Y'$ a closed subset of $Y$, and $X'=f^{-1}(Y')$. Then the
formal prescheme $X_{/X'}$ is identified, by means of the commutative
diagram
\[ \xymatrix{
    X\ar[d]_f &
    X_{/X'}\ar[l]\ar[d]^{\widehat f}\\
    Y &
    Y_{/Y'}\ar[l]
  }
\tag{10.9.9.1}\label{I.10.9.9.1}
\]
with the product $X\times_Y(Y_{/Y'})$ of formal preschemes.
\end{corollary}

It suffices to apply (\hyperref[I.10.9.7]{10.9.7}), replacing
$S$ and $S'$ by $Y$, and $X$ and $X'$ by $X$.
\oldpage[I]{201}

\begin{remark}[10.9.10]
\label{I.10.9.10}
If $X$ is the sum $X_1\amalg X_2$
(\hyperref[subsection:I.3.1]{3.1}) and $X'$ the union
$X'_1\cup X'_2$, where $X'_i$ is a closed subset of $X_i$ ($i=1,2$),
we see immediately that
$X_{/X'}={X_1}_{/X'_1}\amalg {X_2}_{/X'_2}$.
\end{remark}

\subsection{Application to coherent sheaves on affine formal
schemes.}
\label{subsection:I.10.10}

\begin{env}[10.10.1]
\label{I.10.10.1}
Throughout this section, $A$ will denote a \emph{Noetherian adic
ring}, and $\mathfrak{J}$ an ideal of definition of $A$. Let
$X=\operatorname{Spec}(A)$, and $\mathfrak{X}=\operatorname{Spf}(A)$,
which is identified with the closed subset $V(\mathfrak{J})$ of $X$
(\hyperref[I.10.1.2]{10.1.2}). Moreover, the definition
(\hyperref[I.10.1.2]{10.1.2}) and the definition
(\hyperref[I.10.8.4]{10.8.4}) show that the
\emph{affine formal scheme $\mathfrak{X}$} is identical with the
\emph{completion $X_{/\mathfrak{X}}$} of the affine scheme $X$ along the
closed subset $\mathfrak{X}$ of its underlying space. To every
coherent $\mathscr{O}_X$-module $\mathscr{F}$ there therefore corresponds
an $\mathscr{O}_{\mathfrak{X}}$-module of finite type
$\mathscr{F}_{/\mathfrak{X}}$, which is moreover a sheaf of topological
modules over the sheaf of topological rings
$\mathscr{O}_{\mathfrak{X}}$. But every coherent $\mathscr{O}_X$-module
$\mathscr{F}$ is of the form $\widetilde M$, where $M$ is an $A$-module
of finite type (\hyperref[I.1.5.1]{1.5.1}); we set
$(\widetilde M)_{/\mathfrak{X}}=M^\Delta$. Moreover, if $u:M\to N$ is
an $A$-homomorphism of $A$-modules of finite type, there corresponds to
it a homomorphism $\widetilde u:\widetilde M\to\widetilde N$, and
consequently also a continuous homomorphism
$\widetilde u_{/\mathfrak{X}}:(\widetilde M)_{/\mathfrak{X}}\to
(\widetilde N)_{/\mathfrak{X}}$, which we will denote by $u^\Delta$. It is
immediate that
$(v\circ u)^\Delta=v^\Delta\circ u^\Delta$; we have thus defined a
\emph{covariant additive functor $M^\Delta$} from the category of
$A$-modules of finite type to that of
$\mathscr{O}_{\mathfrak{X}}$-modules of finite type. When $A$ is a
\emph{discrete} ring, we have $M^\Delta=\widetilde M$.
\end{env}

\begin{proposition}[10.10.2]
\label{I.10.10.2}
\medskip\noindent
\begin{enumerate}
  \item[\rm{(i)}] $M^\Delta$ is an exact functor in $M$, and there exists
    a canonical functorial isomorphism of $A$-modules
    $\Gamma(\mathfrak{X},M^\Delta)\xrightarrow{\sim}M$.
  \item[\rm{(ii)}] If $M$ and $N$ are two $A$-modules of finite type,
    there exist canonical functorial isomorphisms
    \[
      (M\otimes_A N)^\Delta
      \xrightarrow{\sim}
      M^\Delta\otimes_{\mathscr{O}_{\mathfrak{X}}}N^\Delta
      \tag{10.10.2.1}\label{I.10.10.2.1}
    \]
    \[
      (\operatorname{Hom}_A(M,N))^\Delta
      \xrightarrow{\sim}
      \mathscr{H}\!om_{\mathscr{O}_{\mathfrak{X}}}
      (M^\Delta,N^\Delta)
      \tag{10.10.2.2}\label{I.10.10.2.2}
    \]
  \item[\rm{(iii)}] The map $u\to u^\Delta$ is a functorial
    isomorphism
    \[
      \operatorname{Hom}_A(M,N)
      \xrightarrow{\sim}
      \operatorname{Hom}_{\mathscr{O}_{\mathfrak{X}}}
      (M^\Delta,N^\Delta)
      \tag{10.10.2.3}\label{I.10.10.2.3}
    \]
\end{enumerate}
\end{proposition}

The exactness of $M^\Delta$ follows from the exactness of the functors
$\widetilde M$ (\hyperref[I.1.3.5]{1.3.5}) and
$\mathscr{F}_{/\mathfrak{X}}$ (\hyperref[I.10.8.8]{10.8.8}). By
definition, $\Gamma(\mathfrak{X},M^\Delta)$ is the separated completion
of the $A$-module $\Gamma(X,\widetilde M)=M$ for the
$\mathfrak{J}$-preadic topology; but since $A$ is complete and $M$ of
finite type, we know (\hyperref[0.7.3.6]{0, 7.3.6}) that $M$ is
separated and complete, which completes the proof of (i). The isomorphism
(\hyperref[I.10.10.2.1]{10.10.2.1}) (resp.
\hyperref[I.10.10.2.2]{10.10.2.2}) comes from the composition of the
isomorphisms (\hyperref[I.1.3.12]{1.3.12, (i)}) and
(\hyperref[I.10.8.10.1]{10.8.10.1}) (resp.
(\hyperref[I.1.3.12]{1.3.12, (ii)}) and
(\hyperref[I.10.8.10.2]{10.8.10.2})). Finally, since
$\operatorname{Hom}_A(M,N)$ is an $A$-module of finite type, we can
apply (i) to it, which identifies
$\Gamma(\mathfrak{X},(\operatorname{Hom}_A(M,N))^\Delta)$ with
$\operatorname{Hom}_A(M,N)$, and use
(\hyperref[I.10.10.2.2]{10.10.2.2}), which proves that the
homomorphism (\hyperref[I.10.10.2.3]{10.10.2.3}) is an
isomorphism.

From (\hyperref[I.10.10.2]{10.10.2}) we deduce a whole series of
consequences analogous to those deduced from
(\hyperref[I.1.3.7]{1.3.7}) and
(\hyperref[I.1.3.12]{1.3.12}); we leave it to the reader to
formulate them.
\oldpage[I]{202}

Note that the exactness property of $M^\Delta$, applied to the exact
sequence
$0\to\mathfrak{J}\to A\to A/\mathfrak{J}\to0$, shows that the sheaf
of ideals of $\mathscr{O}_{\mathfrak{X}}$ denoted here by
$\mathfrak{J}^\Delta$ coincides with the one that was denoted in the same
way in (\hyperref[I.10.3.1]{10.3.1}), by virtue of
(\hyperref[I.10.3.2]{10.3.2}).

\begin{proposition}[10.10.3]
\label{I.10.10.3}
Under the hypotheses of (\hyperref[I.10.10.1]{10.10.1}),
$\mathscr{O}_{\mathfrak{X}}$ is a coherent sheaf of rings.
\end{proposition}

If $f\in A$, we know that $A_{\{f\}}$ is a Noetherian adic ring
(\hyperref[0.7.6.11]{0, 7.6.11}), and since the question is local,
we are reduced (\hyperref[I.10.1.4]{10.1.4}) to proving that the
kernel of a homomorphism
$v:\mathscr{O}_{\mathfrak{X}}^n\to\mathscr{O}_{\mathfrak{X}}$ is an
$\mathscr{O}_{\mathfrak{X}}$-module of finite type. We then have
$v=u^\Delta$, where $u$ is an $A$-homomorphism $A^n\to A$
(\hyperref[I.10.10.2]{10.10.2}); since $A$ is Noetherian, the kernel
of $u$ is of finite type, in other words we have a homomorphism
$A^m\xrightarrow{w}A^n$ such that the sequence
$A^m\xrightarrow{w}A^n\xrightarrow{u}A$ is exact. We conclude from this
(\hyperref[I.10.10.2]{10.10.2}) that the sequence
$\mathscr{O}_{\mathfrak{X}}^m\xrightarrow{w^\Delta}
\mathscr{O}_{\mathfrak{X}}^n\xrightarrow{v}
\mathscr{O}_{\mathfrak{X}}$ is exact, which proves that the kernel of
$v$ is of finite type.

\begin{env}[10.10.4]
\label{I.10.10.4}
With the preceding notation, set
$A_n=A/\mathfrak{J}^{n+1}$, and let $X_n$ be the affine scheme
$\operatorname{Spec}(A_n)=(\mathfrak{X},
\mathscr{O}_{\mathfrak{X}}/\mathscr{J}^{n+1})$,
$\mathscr{J}=\mathfrak{J}^\Delta$ being the sheaf of ideals of
definition of $\mathscr{O}_{\mathfrak{X}}$ corresponding to the ideal
$\mathfrak{J}$. Let $u_{mn}$ be the morphism of preschemes
$X_m\to X_n$ corresponding to the canonical homomorphism
$A_n\to A_m$ for $m\leq n$; the formal scheme $\mathfrak{X}$ is the
inductive limit of the $X_n$ for the $u_{mn}$
(\hyperref[I.10.6.3]{10.6.3}).
\end{env}

\begin{proposition}[10.10.5]
\label{I.10.10.5}
Under the hypotheses of (\hyperref[I.10.10.1]{10.10.1}), let
$\mathscr{F}$ be an $\mathscr{O}_{\mathfrak{X}}$-module. The following
conditions are equivalent:
\begin{enumerate}
  \item[a)] $\mathscr{F}$ is a coherent $\mathscr{O}_{\mathfrak{X}}$-module.
  \item[b)] $\mathscr{F}$ is isomorphic to the projective limit
    (\hyperref[I.10.6.6]{10.6.6}) of a sequence $(\mathscr{F}_n)$ of
    coherent $\mathscr{O}_{X_n}$-modules such that
    $u_{mn}^*(\mathscr{F}_n)=\mathscr{F}_m$.
  \item[c)] There exists an $A$-module of finite type $M$ (determined up
    to canonical isomorphism by
    (\hyperref[I.10.10.2]{10.10.2, (i)})) such that $\mathscr{F}$ is
    isomorphic to $M^\Delta$.
\end{enumerate}
\end{proposition}

Let us first show that b) implies c). We have
$\mathscr{F}_n=\widetilde{M_n}$, where $M_n$ is an $A_n$-module of
finite type, and the hypothesis implies that
$M_m=M_n\otimes_{A_n}A_m$ for $m\leq n$
(\hyperref[I.1.6.5]{1.6.5}); the $M_n$ therefore form a projective
system for the canonical di-homomorphisms $M_n\to M_m$
($m\leq n$), and it follows immediately from the definition of the $A_n$
that this projective system satisfies the conditions of
(\hyperref[0.7.2.9]{0, 7.2.9}); its projective limit $M$ is
consequently an $A$-module of finite type such that
$M_n=M\otimes_A A_n$ for every $n$. We deduce from this that
$\mathscr{F}_n$ is induced on $X_n$ by
$\widetilde M\otimes_{\mathscr{O}_X}
(\mathscr{O}_X/\widetilde{\mathfrak{J}}^{n+1})$, hence
$\mathscr{F}=M^\Delta$ by definition
(\hyperref[I.10.8.4]{10.8.4}).

Conversely, c) implies b); indeed, if $u_n$ is the immersion morphism
$X_n\to X$, then
$u_n^*(\widetilde M)=(M\otimes_A A_n)^{\sim}$ is induced on $X_n$ by
$\widetilde M\otimes_{\mathscr{O}_X}
(\mathscr{O}_X/\widetilde{\mathfrak{J}}^{n+1})$, and
$M^\Delta=\varprojlim u_n^*(\widetilde M)$ by definition
(\hyperref[I.10.8.4]{10.8.4}); since
$u_m=u_n\circ u_{mn}$ for $m\leq n$, the
$\mathscr{F}_n=u_n^*(\widetilde M)$ satisfy the conditions of b),
whence our assertion.

Let us now show that c) implies a): indeed, we have by definition
$\mathscr{O}_{\mathfrak{X}}=A^\Delta$; $M$ being the cokernel of a
homomorphism $A^m\to A^n$, it follows from
(\hyperref[I.10.10.2]{10.10.2}) that $M^\Delta$ is the cokernel of a
homomorphism
$\mathscr{O}_{\mathfrak{X}}^m\to\mathscr{O}_{\mathfrak{X}}^n$, and
since the sheaf of rings $\mathscr{O}_{\mathfrak{X}}$ is coherent
(\hyperref[I.10.10.3]{10.10.3}), so is $M^\Delta$
(\hyperref[0.5.3.4]{0, 5.3.4}).
\oldpage[I]{203}

Finally, a) implies b). As an
$\mathscr{O}_{\mathfrak{X}}$-module, we have
$\mathscr{O}_{X_n}=\mathscr{O}_{\mathfrak{X}}/
\mathscr{J}^{n+1}=A_n^\Delta$; $\mathscr{F}_n=\mathscr{F}
\otimes_{\mathscr{O}_{\mathfrak{X}}}\mathscr{O}_{X_n}$ is a coherent
$\mathscr{O}_{\mathfrak{X}}$-module
(\hyperref[0.5.3.5]{0, 5.3.5}), and since it is also an
$\mathscr{O}_{X_n}$-module and $\mathscr{J}^{n+1}$ is coherent,
we conclude that $\mathscr{F}_n$ is a coherent
$\mathscr{O}_{X_n}$-module
(\hyperref[0.5.3.10]{0, 5.3.10}), and it is immediate that
$u_{mn}^*(\mathscr{F}_n)=\mathscr{F}_m$ for $m\leq n$ (remembering
that the continuous map $X_m\to X_n$ of the underlying spaces is the
identity of $\mathfrak{X}$). The sheaf
$\mathscr{G}=\varprojlim\mathscr{F}_n$ is therefore a coherent
$\mathscr{O}_{\mathfrak{X}}$-module, since we have seen that b)
implies a). The canonical homomorphisms
$\mathscr{F}\to\mathscr{F}_n$ form a projective system, which by
passing to the limit gives a canonical homomorphism
$w:\mathscr{F}\to\mathscr{G}$, and everything comes down to proving that
$w$ is bijective. The question now being local, we may restrict to the
case where $\mathscr{F}$ is the cokernel of a homomorphism
$\mathscr{O}_{\mathfrak{X}}^p\to\mathscr{O}_{\mathfrak{X}}^q$; this
homomorphism being of the form $v^\Delta$, where $v$ is a homomorphism
$A^p\to A^q$ (\hyperref[I.10.10.2]{10.10.2}), $\mathscr{F}$ is
isomorphic to $M^\Delta$, where $M=\operatorname{Coker}v$
(\hyperref[I.10.10.2]{10.10.2}). We then have, by virtue of
(\hyperref[I.10.10.2]{10.10.2}),
$\mathscr{F}_n=M^\Delta\otimes_{\mathscr{O}_{\mathfrak{X}}}A_n^\Delta
=(M\otimes_A A_n)^\Delta$, and since the $\mathfrak{J}$-adic topology
on $M\otimes_A A_n$ is discrete, we have
$(M\otimes_A A_n)^\Delta=(M\otimes_A A_n)^{\sim}$ (as an
$\mathscr{O}_{X_n}$-module); we have seen above that
$M^\Delta=\varprojlim\mathscr{F}_n$, and $w$ is therefore indeed the
identity in this case.
\par\hfill Q.E.D.\par

\begin{corollary}[10.10.6]
\label{I.10.10.6}
If $\mathscr{F}$ satisfies condition b) of
(\hyperref[I.10.10.5]{10.10.5}), the projective system
$(\mathscr{F}_n)$ is isomorphic to the system of the
$\mathscr{F}\otimes_{\mathscr{O}_{\mathfrak{X}}}\mathscr{O}_{X_n}$.
\end{corollary}

\begin{env}[10.10.7]
\label{I.10.10.7}
Now let $A$, $B$ be two Noetherian adic rings, and
$\varphi:B\to A$ a continuous homomorphism; we denote by
$\mathfrak{J}$ (resp. $\mathfrak{K}$) an ideal of definition of $A$
(resp. $B$), such that $\varphi(\mathfrak{K})\subset\mathfrak{J}$, and
we set $X=\operatorname{Spec}(A)$, $Y=\operatorname{Spec}(B)$,
$\mathfrak{X}=\operatorname{Spf}(A)$,
$\mathfrak{Y}=\operatorname{Spf}(B)$. Let $f:X\to Y$ be the morphism
of preschemes corresponding to $\varphi$
(\hyperref[I.1.6.1]{1.6.1}), and $\widehat f:\mathfrak{X}\to
\mathfrak{Y}$ its extension to the completions
(\hyperref[I.10.9.1]{10.9.1}), which is also the morphism of
formal preschemes corresponding to $\varphi$
(\hyperref[I.10.2.2]{10.2.2}).
\end{env}

\begin{proposition}[10.10.8]
\label{I.10.10.8}
For every $B$-module $N$ of finite type, there exists a canonical
functorial isomorphism of $\mathscr{O}_{\mathfrak{X}}$-modules
\[
  \widehat f^*(N^\Delta)
  \xrightarrow{\sim}
  (N\otimes_B A)^\Delta.
\]
\end{proposition}

Indeed, denoting by $i_X:\mathfrak{X}\to X$ and
$i_Y:\mathfrak{Y}\to Y$ the canonical morphisms, we have
(\hyperref[I.10.8.8]{10.8.8}), up to canonical functorial
isomorphisms, $N^\Delta=i_Y^*(\widetilde N)$ and
\[
  (N\otimes_B A)^\Delta
  =i_X^*((N\otimes_B A)^{\sim})
  =i_X^*(f^*(\widetilde N))
\]
(\hyperref[I.1.6.5]{1.6.5}); the proposition therefore follows from
the commutativity of the diagram (\hyperref[I.10.9.2]{10.9.2}).

\begin{corollary}[10.10.9]
\label{I.10.10.9}
For every ideal $\mathfrak b$ of $B$, we have
$\widehat f^*(\mathfrak b^\Delta)\mathscr{O}_{\mathfrak{X}}
=(\mathfrak bA)^\Delta$.
\end{corollary}

Indeed, let $j$ be the canonical injection $\mathfrak b\to B$, to which
there corresponds the canonical injection
$j^\Delta:\mathfrak b^\Delta\to\mathscr{O}_{\mathfrak{Y}}$ of
sheaves of $\mathscr{O}_{\mathfrak{Y}}$-modules; by definition,
$\widehat f^*(\mathfrak b^\Delta)\mathscr{O}_{\mathfrak{X}}$ is
the image of the homomorphism
$\widehat f^*(j^\Delta):\widehat f^*(\mathfrak b^\Delta)\to
\mathscr{O}_{\mathfrak{X}}=\widehat f^*(\mathscr{O}_{\mathfrak{Y}})$;
but this homomorphism is identified with
$(j\otimes1)^\Delta:(\mathfrak b\otimes_B A)^\Delta\to
\mathscr{O}_{\mathfrak{X}}=(B\otimes_B A)^\Delta$ by
(\hyperref[I.10.10.8]{10.10.8}). Since the image of $j\otimes1$ is
the ideal $\mathfrak bA$ of $A$, the image of $(j\otimes1)^\Delta$ is
therefore $(\mathfrak bA)^\Delta$ by virtue of
(\hyperref[I.10.10.2]{10.10.2}), whence the conclusion.
\oldpage[I]{204}


\subsection{Coherent sheaves on formal preschemes.}
\label{subsection:I.10.11}

\begin{proposition}[10.11.1]
\label{I.10.11.1}
If $\mathfrak{X}$ is a locally Noetherian formal prescheme, the
sheaf of rings $\mathscr{O}_{\mathfrak{X}}$ is coherent and every
sheaf of ideals of definition of $\mathfrak{X}$ is coherent.
\end{proposition}

The question being local, we are reduced to the case of a Noetherian
affine formal scheme, and the proposition therefore follows from
(\hyperref[I.10.10.3]{10.10.3}) and
(\hyperref[I.10.10.5]{10.10.5}).

\begin{env}[10.11.2]
\label{I.10.11.2}
Let $\mathfrak{X}$ be a locally Noetherian formal prescheme,
$\mathscr{J}$ a sheaf of ideals of definition of $\mathfrak{X}$,
$X_n$ the locally Noetherian (ordinary) prescheme
$(\mathfrak{X},\mathscr{O}_{\mathfrak{X}}/\mathscr{J}^{n+1})$, so
that $\mathfrak{X}$ is the inductive limit of the sequence $(X_n)$ for
the canonical morphisms $u_{mn}:X_m\to X_n$
(\hyperref[I.10.6.3]{10.6.3}). With this notation:
\end{env}

\begin{theorem}[10.11.3]
\label{I.10.11.3}
For an $\mathscr{O}_{\mathfrak{X}}$-module $\mathscr{F}$ to be
coherent, it is necessary and sufficient that it be isomorphic to a
projective limit of a sequence $(\mathscr{F}_n)$, where $\mathscr{F}_n$ is a
coherent $\mathscr{O}_{X_n}$-module, such that
$u_{mn}^*(\mathscr{F}_n)=\mathscr{F}_m$ for $m\leq n$
(\hyperref[I.10.6.6]{10.6.6}). The projective system
$(\mathscr{F}_n)$ is then isomorphic to the system of the
$u_n^*(\mathscr{F})=\mathscr{F}
\otimes_{\mathscr{O}_{\mathfrak{X}}}\mathscr{O}_{X_n}$, $u_n$ being
the canonical morphism $X_n\to\mathfrak{X}$.
\end{theorem}

The question being local, we are reduced to the case where $\mathfrak{X}$ is
a Noetherian affine formal scheme, and the theorem is then a consequence
of (\hyperref[I.10.10.5]{10.10.5}) and
(\hyperref[I.10.10.6]{10.10.6}).

We can therefore say that \emph{giving a coherent
$\mathscr{O}_{\mathfrak{X}}$-module is equivalent to giving a
projective system $(\mathscr{F}_n)$ of
coherent $\mathscr{O}_{X_n}$-modules such that
$u_{mn}^*(\mathscr{F}_n)=\mathscr{F}_m$ for $m\leq n$}.

\begin{corollary}[10.11.4]
\label{I.10.11.4}
If $\mathscr{F}$ and $\mathscr{G}$ are two coherent
$\mathscr{O}_{\mathfrak{X}}$-modules, we can (with the
notation of (\hyperref[I.10.11.3]{10.11.3})) define a
canonical functorial isomorphism
\[
  \label{I.10.11.4.1}
  \operatorname{Hom}_{\mathscr{O}_{\mathfrak{X}}}
  (\mathscr{F},\mathscr{G})
  \xrightarrow{\sim}
  \varprojlim_n
  \operatorname{Hom}_{\mathscr{O}_{X_n}}
  (\mathscr{F}_n,\mathscr{G}_n).
  \tag{10.11.4.1}
\]
\end{corollary}

The projective limit on the right-hand side is to be understood with
respect to the maps
$\theta_n\mapsto u_{mn}^*(\theta_n)$ ($m\leq n$) from
$\operatorname{Hom}_{\mathscr{O}_{X_n}}
(\mathscr{F}_n,\mathscr{G}_n)$ to
$\operatorname{Hom}_{\mathscr{O}_{X_m}}
(\mathscr{F}_m,\mathscr{G}_m)$. The homomorphism
(\hyperref[I.10.11.4.1]{10.11.4.1}) assigns to an element
$\theta\in\operatorname{Hom}_{\mathscr{O}_{\mathfrak{X}}}
(\mathscr{F},\mathscr{G})$ the sequence $(u_n^*(\theta))$; we see
immediately that we define a homomorphism inverse to the preceding one by
assigning to the projective system
$(\theta_n)\in\varprojlim_n
\operatorname{Hom}_{\mathscr{O}_{X_n}}
(\mathscr{F}_n,\mathscr{G}_n)$
its projective limit in
$\operatorname{Hom}_{\mathscr{O}_{\mathfrak{X}}}
(\mathscr{F},\mathscr{G})$, taking into account
(\hyperref[I.10.11.3]{10.11.3}).

\begin{corollary}[10.11.5]
\label{I.10.11.5}
For a homomorphism $\theta:\mathscr{F}\to\mathscr{G}$ to be
surjective, it is necessary and sufficient that the corresponding homomorphism
$\theta_0=u_0^*(\theta):\mathscr{F}_0\to\mathscr{G}_0$ be so.
\end{corollary}

The question being local, we are reduced to the case where
$\mathfrak{X}=\operatorname{Spf}(A)$, $A$ being a Noetherian adic ring,
$\mathscr{F}=M^\Delta$, $\mathscr{G}=N^\Delta$ and $\theta=u^\Delta$,
where $M$ and $N$ are $A$-modules of finite type and $u$ a
homomorphism $M\to N$; moreover, we then have
$\theta_0=\widetilde{u_0}$, where $u_0$ is the homomorphism
$u\otimes1:M\otimes_A A/\mathfrak{J}\to
N\otimes_A A/\mathfrak{J}$;
the conclusion follows from the fact that $\theta$ and $u$ (resp. $\theta_0$ and
$u_0$) are simultaneously surjective
(\hyperref[I.1.3.9]{1.3.9} and
\hyperref[I.10.10.2]{10.10.2}) and from the fact that $u$ and $u_0$ are
simultaneously surjective (\hyperref[0.7.1.14]{0, 7.1.14}).

\begin{env}[10.11.6]
\label{I.10.11.6}
Theorem (\hyperref[I.10.11.3]{10.11.3}) shows that we can consider
every coherent $\mathscr{O}_{\mathfrak{X}}$-module $\mathscr{F}$ as
a \emph{topological} $\mathscr{O}_{\mathfrak{X}}$-module, by
regarding it as the projective limit of the \emph{pseudo-discrete}
sheaves of groups $\mathscr{F}_n$
(\hyperref[0.3.8.1]{0, 3.8.1}). It then follows from
(\hyperref[I.10.11.4]{10.11.4}) that every homomorphism
$u:\mathscr{F}\to\mathscr{G}$ of
coherent $\mathscr{O}_{\mathfrak{X}}$-modules is automatically
\emph{continuous}

\oldpage[I]{205}

(\hyperref[0.3.8.2]{0, 3.8.2}). Moreover, if $\mathscr{H}$ is a
coherent sub-$\mathscr{O}_{\mathfrak{X}}$-module of a
coherent $\mathscr{O}_{\mathfrak{X}}$-module $\mathscr{F}$, then for every
open subset $U\subset\mathfrak{X}$, $\Gamma(U,\mathscr{H})$ is a \emph{closed}
subgroup of the topological group $\Gamma(U,\mathscr{F})$, because, the functor
$\Gamma$ being left exact, $\Gamma(U,\mathscr{H})$ is the kernel of
the homomorphism
$\Gamma(U,\mathscr{F})\to\Gamma(U,\mathscr{F}/\mathscr{H})$, which is
\emph{continuous} by the above, since $\mathscr{F}/\mathscr{H}$
is coherent (\hyperref[0.5.3.4]{0, 5.3.4}); our assertion follows
from the fact that $\Gamma(U,\mathscr{F}/\mathscr{H})$ is a separated
topological group.
\end{env}

\begin{proposition}[10.11.7]
\label{I.10.11.7}
Let $\mathscr{F}$ and $\mathscr{G}$ be two coherent
$\mathscr{O}_{\mathfrak{X}}$-modules. We can define (with the
notation of (\hyperref[I.10.11.3]{10.11.3})) canonical functorial
isomorphisms of topological $\mathscr{O}_{\mathfrak{X}}$-modules
(\hyperref[I.10.11.6]{10.11.6})
\[
  \label{I.10.11.7.1}
  \mathscr{F}\otimes_{\mathscr{O}_{\mathfrak{X}}}\mathscr{G}
  \xrightarrow{\sim}
  \varprojlim_n
  (\mathscr{F}_n\otimes_{\mathscr{O}_{X_n}}\mathscr{G}_n)
  \tag{10.11.7.1}
\]
\[
  \label{I.10.11.7.2}
  \mathscr{H}\!om_{\mathscr{O}_{\mathfrak{X}}}
  (\mathscr{F},\mathscr{G})
  \xrightarrow{\sim}
  \varprojlim_n
  \mathscr{H}\!om_{\mathscr{O}_{X_n}}(\mathscr{F}_n,\mathscr{G}_n)
  \tag{10.11.7.2}
\]
\end{proposition}

The existence of the isomorphism
(\hyperref[I.10.11.7.1]{10.11.7.1}) follows from the formula
\[
  \mathscr{F}_n\otimes_{\mathscr{O}_{X_n}}\mathscr{G}_n
  =
  (\mathscr{F}\otimes_{\mathscr{O}_{\mathfrak{X}}}
  \mathscr{O}_{X_n})
  \otimes_{\mathscr{O}_{X_n}}
  (\mathscr{G}\otimes_{\mathscr{O}_{\mathfrak{X}}}
  \mathscr{O}_{X_n})
  =
  (\mathscr{F}\otimes_{\mathscr{O}_{\mathfrak{X}}}\mathscr{G})
  \otimes_{\mathscr{O}_{\mathfrak{X}}}\mathscr{O}_{X_n}
\]
and from (\hyperref[I.10.11.3]{10.11.3}). The isomorphism
(\hyperref[I.10.11.7.2]{10.11.7.2}), where the two sides are
considered as sheaves of modules without topology, follows from the
definition of the sections of
$\mathscr{H}\!om_{\mathscr{O}_{\mathfrak{X}}}
(\mathscr{F},\mathscr{G})$ and
$\mathscr{H}\!om_{\mathscr{O}_{X_n}}(\mathscr{F}_n,\mathscr{G}_n)$ and from
the existence of the isomorphism
(\hyperref[I.10.11.4.1]{10.11.4.1}), applied to the prescheme induced on
an arbitrary Noetherian affine formal open subset of $\mathfrak{X}$. It remains
to prove that the isomorphism
(\hyperref[I.10.11.7.2]{10.11.7.2}) is bicontinuous over a
quasi-compact set, and we are therefore reduced to the case where
$\mathfrak{X}=\operatorname{Spf}(A)$, $A$ being a Noetherian adic ring, whence
(\hyperref[I.10.10.5]{10.10.5})
$\mathscr{F}=M^\Delta$, $\mathscr{G}=N^\Delta$, $M$, $N$ being
$A$-modules of finite type; taking into account
(\hyperref[I.10.10.2.1]{10.10.2.1}),
(\hyperref[I.10.10.2.3]{10.10.2.3}) and
(\hyperref[I.1.3.12]{1.3.12, (ii)}), we are reduced to showing that
the canonical isomorphism
$\operatorname{Hom}_A(M,N)\xrightarrow{\sim}
\varprojlim_n\operatorname{Hom}_{A_n}(M_n,N_n)$ (with
$M_n=M\otimes_A A_n$, $N_n=N\otimes_A A_n$) is continuous, which was
proved in (\hyperref[0.7.8.2]{0, 7.8.2}).

\begin{env}[10.11.8]
\label{I.10.11.8}
Since $\operatorname{Hom}_{\mathscr{O}_{\mathfrak{X}}}
(\mathscr{F},\mathscr{G})$ is the group of sections of the sheaf of
topological groups
$\mathscr{H}\!om_{\mathscr{O}_{\mathfrak{X}}}
(\mathscr{F},\mathscr{G})$, it is endowed with a group topology. If
$\mathfrak{X}$ is \emph{Noetherian}, it follows from
(\hyperref[I.10.11.7.2]{10.11.7.2}) that a fundamental system of
neighbourhoods of $0$ in this group is obtained by taking the subgroups
$\operatorname{Hom}_{\mathscr{O}_{\mathfrak{X}}}
(\mathscr{F},\mathscr{J}^n\mathscr{G})$ ($n$ arbitrary).
\end{env}

\begin{proposition}[10.11.9]
\label{I.10.11.9}
Let $\mathfrak{X}$ be a Noetherian formal prescheme, $\mathscr{F}$ and
$\mathscr{G}$ two coherent $\mathscr{O}_{\mathfrak{X}}$-modules. In
the topological group
$\operatorname{Hom}_{\mathscr{O}_{\mathfrak{X}}}
(\mathscr{F},\mathscr{G})$ the surjective (resp.
injective, bijective) homomorphisms form an open subset.
\end{proposition}

By virtue of (\hyperref[I.10.11.5]{10.11.5}), the set of
surjective homomorphisms in
$\operatorname{Hom}_{\mathscr{O}_{\mathfrak{X}}}
(\mathscr{F},\mathscr{G})$ is the inverse image under the continuous
map
$\operatorname{Hom}_{\mathscr{O}_{\mathfrak{X}}}
(\mathscr{F},\mathscr{G})\to
\operatorname{Hom}_{\mathscr{O}_{X_0}}(\mathscr{F}_0,\mathscr{G}_0)$
of a subset of the discrete group
$\operatorname{Hom}_{\mathscr{O}_{X_0}}(\mathscr{F}_0,\mathscr{G}_0)$,
whence the first assertion. To prove the second, let us cover
$\mathfrak{X}$ by finitely many Noetherian affine formal open subsets
$U_i$. For
$\theta\in\operatorname{Hom}_{\mathscr{O}_{\mathfrak{X}}}
(\mathscr{F},\mathscr{G})$ to be injective, it is necessary and sufficient
that all its images under the (continuous) restriction maps
$\operatorname{Hom}_{\mathscr{O}_{\mathfrak{X}}}
(\mathscr{F},\mathscr{G})\to
\operatorname{Hom}_{\mathscr{O}_{\mathfrak{X}}|U_i}
(\mathscr{F}|U_i,\mathscr{G}|U_i)$ be so; we are therefore reduced to the
affine case, and in that case this has already been proved in
(\hyperref[0.7.8.3]{0, 7.8.3}).

\oldpage[I]{206}

\subsection{Adic morphisms of formal preschemes.}
\label{subsection:I.10.12}

\begin{env}[10.12.1]
\label{I.10.12.1}
Let $\mathfrak{X}$, $\mathfrak{S}$ be two \emph{locally Noetherian}
formal preschemes; we will say that a morphism
$f:\mathfrak{X}\to\mathfrak{S}$ is \emph{adic} if there exists an ideal
of definition $\mathscr{J}$ of $\mathfrak{S}$ such that
$\mathscr{K}=f^*(\mathscr{J})\mathscr{O}_{\mathfrak{X}}$ is an ideal of
definition of $\mathfrak{X}$; we then also say that $\mathfrak{X}$ is
an \emph{adic $\mathfrak{S}$-prescheme} (for $f$). When this is
the case, for \emph{every} ideal of definition $\mathscr{J}_1$ of
$\mathfrak{S}$,
$\mathscr{K}_1=f^*(\mathscr{J}_1)\mathscr{O}_{\mathfrak{X}}$ is an ideal
of definition of $\mathfrak{X}$. Indeed, the question being local, we
may suppose $\mathfrak{X}$ and $\mathfrak{S}$ affine and Noetherian; there
therefore exists an integer $n$ such that $\mathscr{J}^n\subset\mathscr{J}_1$ and
$\mathscr{J}_1^n\subset\mathscr{J}$
(\hyperref[I.10.3.6]{10.3.6} and
\hyperref[0.7.1.4]{0, 7.1.4}), whence
$\mathscr{K}^n\subset\mathscr{K}_1$ and
$\mathscr{K}_1^n\subset\mathscr{K}$. The first of these relations proves
that $\mathscr{K}_1=\mathfrak{K}_1^\Delta$, where $\mathfrak{K}_1$ is an
open ideal of
$A=\Gamma(\mathfrak{X},\mathscr{O}_{\mathfrak{X}})$, and the second proves
that $\mathfrak{K}_1$ is an ideal of definition of $A$
(\hyperref[0.7.1.4]{0, 7.1.4}), whence our assertion.
\end{env}

It follows immediately from the above that if $\mathfrak{X}$ and
$\mathfrak{Y}$ are two adic $\mathfrak{S}$-preschemes, \emph{every
$\mathfrak{S}$-morphism $u:\mathfrak{X}\to\mathfrak{Y}$ is adic}:
indeed, if $f:\mathfrak{X}\to\mathfrak{S}$,
$g:\mathfrak{Y}\to\mathfrak{S}$ are the structural morphisms, and
$\mathscr{J}$ an ideal of definition of $\mathfrak{S}$, we have
$f=g\circ u$, so
$u^*(g^*(\mathscr{J})\mathscr{O}_{\mathfrak{Y}})
\mathscr{O}_{\mathfrak{X}}=f^*(\mathscr{J})
\mathscr{O}_{\mathfrak{X}}$ is an ideal of definition of
$\mathfrak{X}$, and by hypothesis
$g^*(\mathscr{J})\mathscr{O}_{\mathfrak{Y}}$ is an ideal of definition
of $\mathfrak{Y}$.

\begin{env}[10.12.2]
\label{I.10.12.2}
In what follows, we will suppose fixed a locally Noetherian formal
prescheme $\mathfrak{S}$ and an ideal of definition $\mathscr{J}$ of
$\mathfrak{S}$; we will set
$S_n=(\mathfrak{S},\mathscr{O}_{\mathfrak{S}}/\mathscr{J}^{n+1})$. The
(locally Noetherian) adic $\mathfrak{S}$-preschemes evidently form
a \emph{category}. We will say that an inductive system
$(X_n)$ of locally Noetherian (ordinary) $S_n$-preschemes is an
\emph{adic $(S_n)$-inductive system} if the structural morphisms
$f_n:X_n\to S_n$ are such that, for $m\leq n$, the diagrams
\[
  \label{I.10.12.2.1}
  \xymatrix{
    X_n \ar[d]_{f_n}
    & X_m \ar[l] \ar[d]^{f_m}\\
    S_n
    & S_m \ar[l]
  }
  \tag{10.12.2.1}
\]
are commutative and \emph{identify $X_m$ with the product
$X_n\times_{S_n}S_m=(X_n)_{(S_m)}$}. The adic inductive systems
form a \emph{category}: indeed, it suffices to define a morphism
$(X_n)\to(Y_n)$ of such systems as an \emph{inductive system of
$S_n$-morphisms $u_n:X_n\to Y_n$} such that $u_m$ is identified with
$(u_n)_{(S_m)}$ for $m\leq n$. This being so:
\end{env}

\begin{theorem}[10.12.3]
\label{I.10.12.3}
There is a canonical equivalence between the category of
adic $\mathfrak{S}$-preschemes and the category of
adic $(S_n)$-inductive systems.
\end{theorem}

The equivalence in question is obtained as follows: if
$\mathfrak{X}$ is an adic $\mathfrak{S}$-prescheme and
$f:\mathfrak{X}\to\mathfrak{S}$ the structural morphism,
$\mathscr{K}=f^*(\mathscr{J})\mathscr{O}_{\mathfrak{X}}$ is an ideal of
definition of $\mathfrak{X}$, and we assign to $\mathfrak{X}$ the
inductive system of the
$X_n=(\mathfrak{X},\mathscr{O}_{\mathfrak{X}}/\mathscr{K}^{n+1})$, the
structural morphism $f_n:X_n\to S_n$ corresponding to $f$
(\hyperref[I.10.5.6]{10.5.6}). Let us first show that $(X_n)$ is an
\emph{adic inductive system}: if $f=(\psi,\theta)$, we have
$\psi^*(\mathscr{J})\mathscr{O}_{\mathfrak{X}}=\mathscr{K}$, hence
$\psi^*(\mathscr{J}^n)\mathscr{O}_{\mathfrak{X}}=\mathscr{K}^n$ for every
$n$, and (by the exactness of the functor $\psi^*$)
$\mathscr{K}^{m+1}/\mathscr{K}^{n+1}
=\psi^*(\mathscr{J}^{m+1}/\mathscr{J}^{n+1})
(\mathscr{O}_{\mathfrak{X}}/\mathscr{K}^{n+1})$ for $m\leq n$; our
conclusion therefore follows from (\hyperref[I.4.4.5]{4.4.5}). It is
moreover immediate to verify that to an $\mathfrak{S}$-morphism
$u:\mathfrak{X}\to\mathfrak{Y}$ of adic $\mathfrak{S}$-preschemes
there corresponds (with evident notation)

\oldpage[I]{207}

an inductive system of $S_n$-morphisms $u_n:X_n\to Y_n$ such that $u_m$
is identified with $(u_n)_{(S_m)}$ for $m\leq n$.

The fact that we have indeed defined an equivalence in this way will
follow from the following more precise proposition:

\begin{proposition}[10.12.3.1]
\label{I.10.12.3.1}
Let $(X_n)$ be an inductive system of $S_n$-preschemes; suppose that the
structural morphisms $f_n:X_n\to S_n$ are such that the diagrams
(\hyperref[I.10.12.2.1]{10.12.2.1}) are commutative and identify
$X_m$ with $X_n\times_{S_n}S_m$ for $m\leq n$. Then the inductive system
$(X_n)$ satisfies conditions b) and c) of
(\hyperref[I.10.6.3]{10.6.3}); let $\mathfrak{X}$ be its inductive limit,
and $f:\mathfrak{X}\to\mathfrak{S}$ the morphism that is the inductive limit
of the inductive system $(f_n)$. Then, if $X_0$ is locally Noetherian,
$\mathfrak{X}$ is locally Noetherian and $f$ is an adic morphism.
\end{proposition}

Since the sheaf of ideals of $\mathscr{O}_{S_n}$ which defines the
sub-prescheme $S_m$ of $S_n$ is nilpotent, the same is true, by virtue of
(\hyperref[I.4.4.5]{4.4.5}), of the sheaf of ideals of
$\mathscr{O}_{X_n}$ defining the sub-prescheme $X_m$ of $X_n$, so the
conditions of (\hyperref[I.10.6.3]{10.6.3}) are indeed satisfied. The
question is consequently local on $\mathfrak{X}$ and $\mathfrak{S}$, and we
may suppose that $\mathfrak{S}=\operatorname{Spf}(A)$,
$\mathscr{J}=\mathfrak{J}^\Delta$, $A$ being a Noetherian
$\mathfrak{J}$-adic ring, and $X_n=\operatorname{Spec}(B_n)$; if
$A_n=A/\mathfrak{J}^{n+1}$, the hypothesis implies that $B_0$ is Noetherian
and that, if we set $\mathfrak{J}_n=\mathfrak{J}/\mathfrak{J}^{n+1}$,
$B_m=B_n/\mathfrak{J}_n^{m+1}B_n$. The kernel of $B_n\to B_0$ is therefore
$\mathfrak{K}_n=\mathfrak{J}_nB_n$ and the kernel of $B_n\to B_m$ is
$\mathfrak{K}_n^{m+1}$ for $m\leq n$; moreover, since $A_1$ is
Noetherian, $\mathfrak{J}_1$ is of finite type over $A_1$, hence
$\overline{\mathfrak{K}}_1=\mathfrak{K}_1/\mathfrak{K}_1^2$ is of finite type
over $B_1$, and \emph{a fortiori} over
$B_0=B_1/\mathfrak{K}_1$; the fact that $\mathfrak{X}$ is Noetherian
then follows from (\hyperref[I.10.6.4]{10.6.4}); if
$B=\varprojlim B_n$, we have $\mathfrak{X}=\operatorname{Spf}(B)$, and if
$\mathfrak{K}$ is the kernel of $B\to B_0$,
$B_n=B/\mathfrak{K}^{n+1}$. If
$\rho_n:A/\mathfrak{J}^{n+1}\to B/\mathfrak{K}^{n+1}$ is the homomorphism
corresponding to $f_n$, we therefore have
\[
  \mathfrak{K}/\mathfrak{K}^{n+1}
  =(B/\mathfrak{K}^{n+1})\rho_n(\mathfrak{J}/\mathfrak{J}^{n+1})
\]
since the homomorphism $\rho:A\to B$ corresponding to $f$ is equal to
$\varprojlim\rho_n$, the ideal $\mathfrak{J}B$ of $B$ is dense in
$\mathfrak{K}$, and since every ideal of $B$ is closed
(\hyperref[0.7.3.5]{0, 7.3.5}), we have $\mathfrak{K}=\mathfrak{J}B$. If
$\mathscr{K}=\mathfrak{K}^\Delta$, the relation
$f^*(\mathscr{J})\mathscr{O}_{\mathfrak{X}}=\mathscr{K}$ then follows from
(\hyperref[I.10.10.9]{10.10.9}) and completes the proof.

\begin{env}[10.12.3.2]
\label{I.10.12.3.2}
The preceding equivalence provides, for two adic $\mathfrak{S}$-preschemes
$\mathfrak{X}$, $\mathfrak{Y}$, a \emph{canonical bijection}
\[
  \operatorname{Hom}_{\mathfrak{S}}(\mathfrak{X},\mathfrak{Y})
  \xrightarrow{\sim}
  \varprojlim_n\operatorname{Hom}_{S_n}(X_n,Y_n)
\]
the projective limit being taken with respect to the maps
$u_n\to(u_n)_{(S_m)}$ for $m\leq n$.
\end{env}

\subsection{Morphisms of finite type.}
\label{subsection:I.10.13}

\begin{proposition}[10.13.1]
\label{I.10.13.1}
Let $\mathfrak{Y}$ be a locally Noetherian formal prescheme,
$\mathscr{K}$ an ideal of definition of $\mathfrak{Y}$, and
$f:\mathfrak{X}\to\mathfrak{Y}$ a morphism of formal preschemes. The
following conditions are equivalent:
\begin{enumerate}
  \item[a)] $\mathfrak{X}$ is locally Noetherian, $f$ is an adic morphism
    (\hyperref[I.10.12.1]{10.12.1}) and, if we set
    $\mathscr{J}=f^*(\mathscr{K})\mathscr{O}_{\mathfrak{X}}$, the morphism
    $f_0:(\mathfrak{X},\mathscr{O}_{\mathfrak{X}}/\mathscr{J})
    \to(\mathfrak{Y},\mathscr{O}_{\mathfrak{Y}}/\mathscr{K})$ deduced from
    $f$ is of finite type.
  \item[b)] $\mathfrak{X}$ is locally Noetherian, and is the inductive
    limit of an adic $(Y_n)$-inductive system $(X_n)$ such that the
    morphism $X_0\to Y_0$ is of finite type.

\oldpage[I]{208}

  \item[c)] Every point of $\mathfrak{Y}$ has a Noetherian affine formal
    open neighbourhood $V$ having the following property:
    \begin{enumerate}
      \item[(\textbf{Q})] $f^{-1}(V)$ is the union of a finite family
        of Noetherian affine formal open subsets $U_i$ such that the
        Noetherian adic ring $\Gamma(U_i,\mathscr{O}_{\mathfrak{X}})$ is
        topologically isomorphic to the quotient of an algebra of restricted
        formal power series (\hyperref[0.7.5.1]{0, 7.5.1}) over
        $\Gamma(V,\mathscr{O}_{\mathfrak{Y}})$ by a (necessarily
        closed) ideal.
    \end{enumerate}
\end{enumerate}
\end{proposition}

It is immediate that a) implies b), by virtue of
(\hyperref[I.10.12.3]{10.12.3}). To show that b) implies c), we
may, since the question is local on $\mathfrak{Y}$, suppose that
$\mathfrak{Y}=\operatorname{Spf}(B)$, where $B$ is a Noetherian adic ring; let
$\mathscr{K}=\mathfrak{K}^\Delta$, $\mathfrak{K}$ being an ideal of
definition of $B$. Since by hypothesis $X_0$ is of finite type over $Y_0$,
$X_0$ is a finite union of affine open subsets $U_i$ such that the ring
$A_{i0}$ of the affine scheme induced by $X_0$ on $U_i$ is an algebra of
finite type over the ring $B/\mathfrak{K}$ of $Y_0$
(\hyperref[I.6.3.2]{6.3.2}). By virtue of
(\hyperref[I.5.1.9]{5.1.9}), $U_i$ is also an affine open subset in each
of the Noetherian preschemes $X_n$, and if $A_{in}$ is the ring of the
affine scheme induced by $X_n$ on $U_i$, hypothesis b) implies that for
$m\leq n$, $A_{im}$ is isomorphic to
$A_{in}/\mathfrak{K}^{m+1}A_{in}$. Consequently, the formal prescheme
induced on $U_i$ by $\mathfrak{X}$ is isomorphic to $\operatorname{Spf}(A_i)$,
where $A_i=\varprojlim_n A_{in}$ (\hyperref[I.10.6.4]{10.6.4}); $A_i$ is a
$\mathfrak{K}A_i$-adic ring, and $A_i/\mathfrak{K}A_i$, isomorphic to
$A_{i0}$, is an algebra of finite type over $B/\mathfrak{K}$. We conclude
from this (\hyperref[0.7.5.5]{0, 7.5.5}) that $A_i$ is topologically
isomorphic to a quotient of an algebra of restricted formal power series
over $B$ (by a necessarily closed ideal, since such an algebra is Noetherian
(\hyperref[0.7.5.4]{0, 7.5.4})).

To prove that c) implies a), we may restrict to the case where
$\mathfrak{X}=\operatorname{Spf}(A)$ is also affine, $A$ being a Noetherian
adic ring isomorphic to the quotient of an algebra of restricted formal
power series over $B$ by a closed ideal. Then
(\hyperref[0.7.5.5]{0, 7.5.5}), $A/\mathfrak{K}A$ is an algebra of finite
type over $B/\mathfrak{K}$, and $\mathfrak{K}A=\mathfrak{J}$ is an ideal of
definition of $A$, so, by virtue of
(\hyperref[I.10.10.9]{10.10.9}), the conditions of a) are satisfied.

Note that if the conditions of Proposition
(\hyperref[I.10.13.1]{10.13.1}) are fulfilled, property a) holds
for \emph{every ideal of definition $\mathscr{K}$} of
$\mathfrak{Y}$ (by virtue of c)), and consequently, in property b),
\emph{all the $f_n$ are morphisms of finite type}.

\begin{corollary}[10.13.2]
\label{I.10.13.2}
If the conditions of (\hyperref[I.10.13.1]{10.13.1}) are satisfied,
every Noetherian affine formal open subset $V$ of $\mathfrak{Y}$ has
property (\textbf{Q}), and if $\mathfrak{Y}$ is Noetherian, so is
$\mathfrak{X}$.
\end{corollary}

This follows immediately from (\hyperref[I.10.13.1]{10.13.1}) and from
(\hyperref[I.6.3.2]{6.3.2}).

\begin{definition}[10.13.3]
\label{I.10.13.3}
When the equivalent properties a), b), c) of
(\hyperref[I.10.13.1]{10.13.1}) are satisfied, we say that the morphism
$f$ is \emph{of finite type}, or that $\mathfrak{X}$ is a
\emph{formal $\mathfrak{Y}$-prescheme of finite type}, or a
\emph{formal prescheme of finite type over $\mathfrak{Y}$}.
\end{definition}

\begin{corollary}[10.13.4]
\label{I.10.13.4}
Let $\mathfrak{X}=\operatorname{Spf}(A)$,
$\mathfrak{Y}=\operatorname{Spf}(B)$ be two Noetherian affine formal
schemes; for $\mathfrak{X}$ to be of finite type over
$\mathfrak{Y}$, it is necessary and sufficient that the Noetherian adic ring
$A$ be isomorphic to the quotient of an algebra of restricted formal power
series over $B$ by a closed ideal.
\end{corollary}

Indeed, with the notation of (\hyperref[I.10.13.1]{10.13.1}), if
$\mathfrak{X}$ is of finite type over $\mathfrak{Y}$,
$A/\mathfrak{K}A$ is then a $(B/\mathfrak{K})$-algebra of finite type by
virtue of (\hyperref[I.6.3.3]{6.3.3}), and $\mathfrak{K}A$ is an ideal of
definition of $A$ (\hyperref[I.10.10.9]{10.10.9}). We therefore conclude
by (\hyperref[0.7.5.5]{0, 7.5.5}).

\begin{proposition}[10.13.5]
\label{I.10.13.5}
\medskip\noindent
\begin{enumerate}
  \item[\rm{(i)}] The composite of two morphisms of formal preschemes which
    are of finite type is of finite type.

\oldpage[I]{209}

  \item[\rm{(ii)}] Let $\mathfrak{X}$, $\mathfrak{S}$,
    $\mathfrak{S}'$ be three locally Noetherian (resp.
    Noetherian) formal preschemes, and $f:\mathfrak{X}\to\mathfrak{S}$,
    $g:\mathfrak{S}'\to\mathfrak{S}$ two morphisms. If $f$ is of finite
    type, $\mathfrak{X}\times_{\mathfrak{S}}\mathfrak{S}'$ is locally
    Noetherian (resp. Noetherian) and is of finite type over $\mathfrak{S}'$.
  \item[\rm{(iii)}] Let $\mathfrak{S}$ be a locally Noetherian formal
    prescheme, and $\mathfrak{X}'$, $\mathfrak{Y}'$ two locally Noetherian
    formal $\mathfrak{S}$-preschemes such that
    $\mathfrak{X}'\times_{\mathfrak{S}}\mathfrak{Y}'$ is locally
    Noetherian. If $\mathfrak{X}$, $\mathfrak{Y}$ are locally Noetherian
    formal $\mathfrak{S}$-preschemes and
    $f:\mathfrak{X}\to\mathfrak{X}'$,
    $g:\mathfrak{Y}\to\mathfrak{Y}'$ two $\mathfrak{S}$-morphisms of
    finite type, then $\mathfrak{X}\times_{\mathfrak{S}}\mathfrak{Y}$ is
    locally Noetherian and $f\times_{\mathfrak{S}}g$ is an
    $\mathfrak{S}$-morphism of finite type.
\end{enumerate}
\end{proposition}

(iii) is deduced from (i) and (ii) by the formal argument of
(\hyperref[I.3.5.1]{3.5.1}), and it therefore suffices to prove (i) and (ii).

Let $\mathfrak{X}$, $\mathfrak{Y}$, $\mathfrak{Z}$ be three locally
Noetherian formal preschemes, and $f:\mathfrak{X}\to\mathfrak{Y}$,
$g:\mathfrak{Y}\to\mathfrak{Z}$ two morphisms of finite type. If
$\mathscr{L}$ is an ideal of definition of $\mathfrak{Z}$,
$\mathscr{K}=g^*(\mathscr{L})\mathscr{O}_{\mathfrak{Y}}$ is one for
$\mathfrak{Y}$ and
$\mathscr{J}=f^*(g^*(\mathscr{L}))\mathscr{O}_{\mathfrak{X}}$ is one
for $\mathfrak{X}$. Set
$X_0=(\mathfrak{X},\mathscr{O}_{\mathfrak{X}}/\mathscr{J})$,
$Y_0=(\mathfrak{Y},\mathscr{O}_{\mathfrak{Y}}/\mathscr{K})$,
$Z_0=(\mathfrak{Z},\mathscr{O}_{\mathfrak{Z}}/\mathscr{L})$, and let
$f_0:X_0\to Y_0$, $g_0:Y_0\to Z_0$ be the morphisms corresponding to $f$ and
$g$. Since by hypothesis $f_0$ and $g_0$ are of finite type, so is
$g_0\circ f_0$ (\hyperref[I.6.3.4]{6.3.4}), which corresponds to
$g\circ f$; hence $g\circ f$ is of finite type by
(\hyperref[I.10.13.1]{10.13.1}).

Under the conditions of (ii), $\mathfrak{S}$ (resp. $\mathfrak{X}$,
$\mathfrak{S}'$) is the inductive limit of a sequence $(S_n)$ (resp. $(X_n)$,
$(S'_n)$) of locally Noetherian preschemes, and we may suppose
(\hyperref[I.10.13.1]{10.13.1}) that
$X_m=X_n\times_{S_n}S_m$ for $m\leq n$. The formal prescheme
$\mathfrak{X}\times_{\mathfrak{S}}\mathfrak{S}'$ is then the inductive limit
of the preschemes $X_n\times_{S_n}S'_n$
(\hyperref[I.10.7.4]{10.7.4}), and we have
\[
  X_m\times_{S_m}S'_m
  =(X_n\times_{S_n}S_m)\times_{S_m}S'_m
  =(X_n\times_{S_n}S'_n)\times_{S'_n}S'_m.
\]
Moreover, $X_0\times_{S_0}S'_0$ is locally Noetherian since $X_0$
is of finite type over $S_0$ (\hyperref[I.6.3.8]{6.3.8}). We conclude
first (\hyperref[I.10.12.3.1]{10.12.3.1}) that
$\mathfrak{X}\times_{\mathfrak{S}}\mathfrak{S}'$ is locally
Noetherian; moreover, since $X_0\times_{S_0}S'_0$ is of finite type over
$S'_0$ (\hyperref[I.6.3.8]{6.3.8}), it follows from
(\hyperref[I.10.12.3.1]{10.12.3.1}) and from
(\hyperref[I.10.13.1]{10.13.1}) that
$\mathfrak{X}\times_{\mathfrak{S}}\mathfrak{S}'$ is of finite type over
$\mathfrak{S}'$, which completes the proof of (ii) (the assertion concerning
Noetherian preschemes being an immediate consequence of
(\hyperref[I.6.3.8]{6.3.8})).

\begin{corollary}[10.13.6]
\label{I.10.13.6}
Under the hypotheses of (\hyperref[I.10.9.9]{10.9.9}), if $f$ is a
morphism of finite type, so is its extension $\widehat f$
to the completions.
\end{corollary}

\subsection{Closed sub-preschemes of formal preschemes.}
\label{subsection:I.10.14}

\begin{proposition}[10.14.1]
\label{I.10.14.1}
Let $\mathfrak{X}$ be a locally Noetherian formal prescheme and
$\mathscr{A}$ a coherent sheaf of ideals of
$\mathscr{O}_{\mathfrak{X}}$. If $\mathfrak{Y}$ is the (closed) support of
$\mathscr{O}_{\mathfrak{X}}/\mathscr{A}$, the topologically ringed space
$(\mathfrak{Y},(\mathscr{O}_{\mathfrak{X}}/\mathscr{A})|\mathfrak{Y})$ is
a locally Noetherian formal prescheme, which is Noetherian if
$\mathfrak{X}$ is.
\end{proposition}

Note that $\mathscr{O}_{\mathfrak{X}}/\mathscr{A}$ is coherent by virtue of
(\hyperref[I.10.10.3]{10.10.3}) and
(\hyperref[0.5.3.4]{0, 5.3.4}), so its support $\mathfrak{Y}$ is closed
(\hyperref[0.5.2.2]{0, 5.2.2}). Let $\mathscr{J}$ be an ideal of definition
of $\mathfrak{X}$, and let
$X_n=(\mathfrak{X},\mathscr{O}_{\mathfrak{X}}/\mathscr{J}^{n+1})$; the
sheaf of rings $\mathscr{O}_{\mathfrak{X}}/\mathscr{A}$ is the projective
limit of the sheaves
$\mathscr{O}_{\mathfrak{X}}/(\mathscr{A}+\mathscr{J}^{n+1})
=(\mathscr{O}_{\mathfrak{X}}/\mathscr{A})
\otimes_{\mathscr{O}_{\mathfrak{X}}}
(\mathscr{O}_{\mathfrak{X}}/\mathscr{J}^{n+1})$
(\hyperref[I.10.11.3]{10.11.3}), which all have support
$\mathfrak{Y}$. The sheaf
$(\mathscr{A}+\mathscr{J}^{n+1})/\mathscr{J}^{n+1}$ is a
coherent $\mathscr{O}_{\mathfrak{X}}$-module, since
$\mathscr{J}^{n+1}$ is coherent, so
$(\mathscr{A}+\mathscr{J}^{n+1})/\mathscr{J}^{n+1}$ is also a
coherent $(\mathscr{O}_{\mathfrak{X}}/\mathscr{J}^{n+1})$-module
(\hyperref[0.5.3.10]{0, 5.3.10}); if $Y_n$ is the closed sub-prescheme
of $X_n$ defined by this sheaf of ideals, it is immediate that
$(\mathfrak{Y},(\mathscr{O}_{\mathfrak{X}}/\mathscr{A})|\mathfrak{Y})$

\oldpage[I]{210}
is the formal prescheme that is the inductive limit of the $Y_n$, and since
the conditions of (\hyperref[I.10.6.4]{10.6.4}) are satisfied, this proves
that this formal prescheme is locally Noetherian, and Noetherian if
$\mathfrak{X}$ is (since then $Y_0$ is, by virtue of
(\hyperref[I.6.1.4]{6.1.4})).

\begin{definition}[10.14.2]
\label{I.10.14.2}
We call a \emph{closed sub-prescheme} of a formal prescheme $\mathfrak{X}$
any formal prescheme
$(\mathfrak{Y},(\mathscr{O}_{\mathfrak{X}}/\mathscr{A})|\mathfrak{Y})$ where
$\mathscr{A}$ is a coherent ideal of $\mathscr{O}_{\mathfrak{X}}$; we say that
this prescheme is the closed sub-prescheme defined by $\mathscr{A}$.
\end{definition}

It is clear that the correspondence thus defined between
coherent ideals of $\mathscr{O}_{\mathfrak{X}}$ and closed sub-preschemes of
$\mathfrak{X}$ is bijective.

The morphism of topologically ringed spaces
$j=(\psi,\theta):\mathfrak{Y}\to\mathfrak{X}$, where $\psi$ is the injection
$\mathfrak{Y}\to\mathfrak{X}$ and $\theta^\sharp$ the canonical homomorphism
$\mathscr{O}_{\mathfrak{X}}\to\mathscr{O}_{\mathfrak{X}}/\mathscr{A}$, is
evidently (\hyperref[I.10.4.5]{10.4.5}) a morphism of formal preschemes,
which we call the \emph{canonical injection} of $\mathfrak{Y}$ into
$\mathfrak{X}$. Note that if $\mathfrak{X}=\operatorname{Spf}(A)$, where $A$
is a Noetherian adic ring, we have $\mathscr{A}=\mathfrak{a}^\Delta$, where
$\mathfrak{a}$ is an ideal of $A$ (\hyperref[I.10.10.5]{10.10.5}), and it
follows immediately from the above that we then have
$\mathfrak{Y}=\operatorname{Spf}(A/\mathfrak{a})$ up to isomorphism, and that
$j$ corresponds (\hyperref[I.10.2.2]{10.2.2}) to the canonical homomorphism
$A\to A/\mathfrak{a}$.

We say that a morphism $f:\mathfrak{Z}\to\mathfrak{X}$ of locally Noetherian
formal preschemes is a \emph{closed immersion} if it factors as
$\mathfrak{Z}\xrightarrow{g}\mathfrak{Y}\xrightarrow{j}\mathfrak{X}$, where $g$
is an isomorphism of $\mathfrak{Z}$ onto a closed sub-prescheme
$\mathfrak{Y}$ of $\mathfrak{X}$ and $j$ the canonical injection. Since $j$ is
a monomorphism of ringed spaces, $g$ and $\mathfrak{Y}$ are necessarily
\emph{unique}.

\begin{proposition}[10.14.3]
\label{I.10.14.3}
A closed immersion is a morphism of finite type.
\end{proposition}

We are immediately reduced to the case where $\mathfrak{X}$ is an affine
formal scheme $\operatorname{Spf}(A)$ and
$\mathfrak{Y}=\operatorname{Spf}(A/\mathfrak{a})$;
the proposition follows from (\hyperref[I.10.13.1]{10.13.1, c)}).

\begin{lemma}[10.14.4]
\label{I.10.14.4}
Let $f:\mathfrak{Y}\to\mathfrak{X}$ be a morphism of locally Noetherian
formal preschemes, and let $(U_\alpha)$ be a covering of
$f(\mathfrak{Y})$ by Noetherian affine formal open subsets of
$\mathfrak{X}$ such that the $f^{-1}(U_\alpha)$ are Noetherian affine formal
open subsets of $\mathfrak{Y}$. For $f$ to be a closed immersion,
it is necessary and sufficient that $f(\mathfrak{Y})$ be a closed subset of
$\mathfrak{X}$ and that, for every $\alpha$, the restriction of $f$ to
$f^{-1}(U_\alpha)$ correspond (\hyperref[I.10.4.6]{10.4.6}) to a
surjective homomorphism
$\Gamma(U_\alpha,\mathscr{O}_{\mathfrak{X}})\to
\Gamma(f^{-1}(U_\alpha),\mathscr{O}_{\mathfrak{Y}})$.
\end{lemma}

The conditions are evidently necessary. Conversely, if they are
fulfilled, and if we denote by $\mathfrak{a}_\alpha$ the kernel of
$\Gamma(U_\alpha,\mathscr{O}_{\mathfrak{X}})\to
\Gamma(f^{-1}(U_\alpha),\mathscr{O}_{\mathfrak{Y}})$, we define a coherent
sheaf of ideals $\mathscr{A}$ of $\mathscr{O}_{\mathfrak{X}}$ by taking
$\mathscr{A}|U_\alpha=\mathfrak{a}_\alpha^\Delta$, and by taking $\mathscr{A}$
equal to $\mathscr{O}_{\mathfrak{X}}$ in the complement of the union of the
$U_\alpha$. Indeed, since
$f(\mathfrak{Y})$ is closed and the support of
$(\mathscr{O}_{\mathfrak{X}}|U_\alpha)/\mathfrak{a}_\alpha^\Delta$ is
$U_\alpha\cap f(\mathfrak{Y})$, everything comes down
to verifying that $\mathfrak{a}_\alpha^\Delta$ and
$\mathfrak{a}_\beta^\Delta$ induce the same sheaf on a Noetherian affine
formal open subset $V\subset U_\alpha\cap U_\beta$. Now, the restriction of
$f$ to $f^{-1}(U_\alpha)$ being a closed immersion of this formal prescheme
into $U_\alpha$, $f^{-1}(V)$ is a Noetherian affine formal open subset in
$f^{-1}(U_\alpha)$ and the restriction of $f$ to $f^{-1}(V)$ is a closed
immersion; if $\mathfrak{b}$ is the kernel of the surjective homomorphism
$\Gamma(V,\mathscr{O}_{\mathfrak{X}})\to
\Gamma(f^{-1}(V),\mathscr{O}_{\mathfrak{Y}})$ corresponding to this
restriction, it is immediate (\hyperref[I.10.10.2]{10.10.2}) that
$\mathfrak{a}_\alpha^\Delta$ induces $\mathfrak{b}^\Delta$ on $V$. The sheaf
of ideals $\mathscr{A}$ being thus defined, it is then clear that
$f=j\circ g$, where $j:\mathfrak{Z}\to\mathfrak{X}$ is the canonical injection
of the closed sub-prescheme $\mathfrak{Z}$ of $\mathfrak{X}$ defined by
$\mathscr{A}$, and $g$ an isomorphism of $\mathfrak{Y}$ onto $\mathfrak{Z}$.

\begin{proposition}[10.14.5]
\label{I.10.14.5}
\begin{enumerate}
  \item[\rm{(i)}] If $f:\mathfrak{Z}\to\mathfrak{Y}$,
  $g:\mathfrak{Y}\to\mathfrak{X}$ are closed immersions of locally
  Noetherian formal preschemes, $g\circ f$ is a closed immersion.

\oldpage[I]{211}
  \item[\rm{(ii)}] Let $\mathfrak{X}$, $\mathfrak{Y}$, $\mathfrak{S}$ be three
  locally Noetherian formal preschemes,
  $f:\mathfrak{X}\to\mathfrak{S}$ a closed immersion, and
  $g:\mathfrak{Y}\to\mathfrak{S}$ a morphism. Then the morphism
  $\mathfrak{X}\times_{\mathfrak{S}}\mathfrak{Y}\to\mathfrak{Y}$ is a
  closed immersion.
  \item[\rm{(iii)}] Let $\mathfrak{S}$ be a locally Noetherian formal
  prescheme, and $\mathfrak{X}'$, $\mathfrak{Y}'$ two locally Noetherian
  formal $\mathfrak{S}$-preschemes such that
  $\mathfrak{X}'\times_{\mathfrak{S}}\mathfrak{Y}'$ is locally
  Noetherian. If $\mathfrak{X}$, $\mathfrak{Y}$ are locally Noetherian
  formal $\mathfrak{S}$-preschemes and
  $f:\mathfrak{X}\to\mathfrak{X}'$, $g:\mathfrak{Y}\to\mathfrak{Y}'$ two
  $\mathfrak{S}$-morphisms which are closed immersions, then
  $f\times_{\mathfrak{S}}g$ is a closed immersion.
\end{enumerate}
\end{proposition}

By virtue of (\hyperref[I.3.5.1]{3.5.1}), it again suffices to prove (i)
and (ii).

To prove (i), we may suppose that $\mathfrak{Y}$ (resp.
$\mathfrak{Z}$) is a closed sub-prescheme of $\mathfrak{X}$ (resp.
$\mathfrak{Y}$) defined by a coherent sheaf $\mathscr{J}$ (resp.
$\mathscr{K}$) of ideals of $\mathscr{O}_{\mathfrak{X}}$ (resp.
$\mathscr{O}_{\mathfrak{Y}}$); if $\psi$ is the injection
$\mathfrak{Y}\to\mathfrak{X}$ of the underlying spaces,
$\psi_*(\mathscr{K})$ is a coherent sheaf of ideals of
$\psi_*(\mathscr{O}_{\mathfrak{Y}})
=\mathscr{O}_{\mathfrak{X}}/\mathscr{J}$
(\hyperref[0.5.3.12]{0, 5.3.12}), hence also a
coherent $\mathscr{O}_{\mathfrak{X}}$-module
(\hyperref[0.5.3.10]{0, 5.3.10}); the kernel $\mathscr{K}_1$ of
$\mathscr{O}_{\mathfrak{X}}\to
(\mathscr{O}_{\mathfrak{X}}/\mathscr{J})/\psi_*(\mathscr{K})$ is therefore a
coherent sheaf of ideals of $\mathscr{O}_{\mathfrak{X}}$
(\hyperref[0.5.3.4]{0, 5.3.4}), and
$\mathscr{O}_{\mathfrak{X}}/\mathscr{K}_1$ is isomorphic to
$\psi_*(\mathscr{O}_{\mathfrak{Y}}/\mathscr{K})$, which proves that
$\mathfrak{Z}$ is isomorphic to a closed sub-prescheme of $\mathfrak{X}$.

To prove (ii), it is immediate that we may restrict to the case where
$\mathfrak{S}=\operatorname{Spf}(A)$,
$\mathfrak{X}=\operatorname{Spf}(B)$,
$\mathfrak{Y}=\operatorname{Spf}(C)$, $A$ being a Noetherian
$\mathfrak{J}$-adic ring, $B=A/\mathfrak{a}$, where $\mathfrak{a}$ is
an ideal of $A$, and $C$ an adic and Noetherian topological $A$-algebra.
Everything comes down to proving that the homomorphism
$C\to C\widehat{\otimes}_A(A/\mathfrak{a})$ is \emph{surjective}: now,
$A/\mathfrak{a}$ is an $A$-module of finite type, and its topology is the
$\mathfrak{J}$-adic topology; it then follows from
(\hyperref[0.7.7.8]{0, 7.7.8}) that
$C\widehat{\otimes}_A(A/\mathfrak{a})$ is identified with
$C\otimes_A(A/\mathfrak{a})=C/\mathfrak{a}C$, whence our assertion.

\begin{corollary}[10.14.6]
\label{I.10.14.6}
Under the hypotheses of (\hyperref[I.10.14.5]{10.14.5, (ii)}), let
$p:\mathfrak{X}\times_{\mathfrak{S}}\mathfrak{Y}\to\mathfrak{X}$,
$q:\mathfrak{X}\times_{\mathfrak{S}}\mathfrak{Y}\to\mathfrak{Y}$ be the
projections, so that the diagram
\[
  \xymatrix{
    \mathfrak{X} \ar[d]_{f}
    & \mathfrak{X}\times_{\mathfrak{S}}\mathfrak{Y}
      \ar[l]_{p} \ar[d]^{q}\\
    \mathfrak{S}
    & \mathfrak{Y} \ar[l]^{g}
  }
\]
is commutative. For every coherent $\mathscr{O}_{\mathfrak{X}}$-module
$\mathscr{F}$, we then have a canonical isomorphism of
$\mathscr{O}_{\mathfrak{Y}}$-modules
\[
  \label{I.10.14.6.1}
  u:g^*(f_*(\mathscr{F}))
  \xrightarrow{\sim}
  q_*(p^*(\mathscr{F})).
  \tag{10.14.6.1}
\]
\end{corollary}

To define a homomorphism
$g^*(f_*(\mathscr{F}))\to q_*(p^*(\mathscr{F}))$, we know that this amounts
to defining a homomorphism
$f_*(\mathscr{F})\to g_*(q_*(p^*(\mathscr{F})))
=f_*(p_*(p^*(\mathscr{F})))$
(\hyperref[0.4.4.3]{0, 4.4.3}): we will take for $u$ the homomorphism
corresponding by adjunction to $f_*(\rho)$, where $\rho$
is the canonical homomorphism
$\mathscr{F}\to p_*(p^*(\mathscr{F}))$
(\hyperref[0.4.4.3]{0, 4.4.3}). To see that $u$ is an isomorphism, we
are immediately reduced to the case where $\mathfrak{S}$, $\mathfrak{X}$,
$\mathfrak{Y}$ are formal spectra of Noetherian adic rings $A$,
$B$, $C$, with the conditions seen above in
(\hyperref[I.10.14.5]{10.14.5, (ii)}); we then have
$\mathscr{F}=M^\Delta$, where $M$ is an $(A/\mathfrak{a})$-module of finite
type (\hyperref[I.10.10.5]{10.10.5}), and the two sides of
(\hyperref[I.10.14.6.1]{10.14.6.1}) are identified respectively, by virtue
of (\hyperref[I.10.10.8]{10.10.8}), with
$(C\otimes_A M)^\Delta$ and with
$((C/\mathfrak{a}C)\otimes_{A/\mathfrak{a}}M)^\Delta$, whence the corollary,
since
$(C/\mathfrak{a}C)\otimes_{A/\mathfrak{a}}M
=(C\otimes_A(A/\mathfrak{a}))\otimes_{A/\mathfrak{a}}M$ is canonically
identified with $C\otimes_A M$.

\oldpage[I]{212}
\begin{corollary}[10.14.7]
\label{I.10.14.7}
Let $X$ be a locally Noetherian ordinary prescheme, $Y$ a closed
sub-prescheme of $X$, $j$ the canonical injection $Y\to X$, $X'$ a closed
subset of $X$, and $Y'=Y\cap X'$; then
$\widehat{j}:Y_{/Y'}\to X_{/X'}$ is a closed immersion, and for every
coherent $\mathscr{O}_Y$-module $\mathscr{F}$, we have
\[
  \widehat{j}_*(\mathscr{F}_{/Y'})=(j_*(\mathscr{F}))_{/X'}.
\]
\end{corollary}

Since $Y'=j^{-1}(X')$, it suffices to use
(\hyperref[I.10.9.9]{10.9.9}) and to apply
(\hyperref[I.10.14.5]{10.14.5}) and
(\hyperref[I.10.14.6]{10.14.6}).


\subsection{Separated formal preschemes.}
\label{subsection:I.10.15}

\begin{definition}[10.15.1]
\label{I.10.15.1}
Let $\mathfrak{S}$ be a formal prescheme, $\mathfrak{X}$ a formal
$\mathfrak{S}$-prescheme, and
$f:\mathfrak{X}\to\mathfrak{S}$ the structural morphism. We call the
morphism $(1_{\mathfrak{X}},1_{\mathfrak{X}})_{\mathfrak{S}}$ the
diagonal morphism
$\Delta_{\mathfrak{X}|\mathfrak{S}}:\mathfrak{X}\to
\mathfrak{X}\times_{\mathfrak{S}}\mathfrak{X}$ (also denoted
$\Delta_{\mathfrak{X}}$). We say that
$\mathfrak{X}$ is separated over $\mathfrak{S}$, or is a formal
$\mathfrak{S}$-scheme, or that $f$ is a separated morphism, if the image
under $\Delta_{\mathfrak{X}}$ of the underlying space $\mathfrak{X}$ is a
closed subset of the space underlying
$\mathfrak{X}\times_{\mathfrak{S}}\mathfrak{X}$. We say that a formal
prescheme $\mathfrak{X}$ is separated, or is a formal scheme, if it is
separated over $\mathbf{Z}$.
\end{definition}

\begin{proposition}[10.15.2]
\label{I.10.15.2}
Suppose that the formal preschemes $\mathfrak{S}$, $\mathfrak{X}$ are
respectively the inductive limits of the sequences $(S_n)$, $(X_n)$ of
ordinary preschemes, and that the morphism
$f:\mathfrak{X}\to\mathfrak{S}$ is the inductive limit of a sequence of
morphisms $f_n:X_n\to S_n$. For $f$ to be separated, it is necessary and
sufficient that the morphism $f_0:X_0\to S_0$ be so.
\end{proposition}

Indeed, $\Delta_{\mathfrak{X}|\mathfrak{S}}$ is then the inductive limit of
the sequence of morphisms $\Delta_{X_n|S_n}$
(\hyperref[I.10.7.4]{10.7.4}), and the image under
$\Delta_{\mathfrak{X}|\mathfrak{S}}$ of the underlying space
$\mathfrak{X}$ (resp. the underlying space
$\mathfrak{X}\times_{\mathfrak{S}}\mathfrak{X}$) is identical to the image
under $\Delta_{X_0|S_0}$ of the underlying space $X_0$ (resp. to the
underlying space $X_0\times_{S_0}X_0$); whence the conclusion.

\begin{proposition}[10.15.3]
\label{I.10.15.3}
In what follows, we suppose that all the formal preschemes (resp.
morphisms of formal preschemes) under consideration are inductive limits of
sequences of ordinary preschemes (resp. of morphisms of ordinary preschemes).
\begin{enumerate}
  \item[\rm{(i)}] The composite of two separated morphisms is separated.
  \item[\rm{(ii)}] If $f:\mathfrak{X}\to\mathfrak{X}'$,
  $g:\mathfrak{Y}\to\mathfrak{Y}'$ are two separated
  $\mathfrak{S}$-morphisms, $f\times_{\mathfrak{S}}g$ is separated.
  \item[\rm{(iii)}] If $f:\mathfrak{X}\to\mathfrak{Y}$ is a separated
  $\mathfrak{S}$-morphism, the $\mathfrak{S}'$-morphism
  $f_{(\mathfrak{S}')}$ is separated for every extension
  $\mathfrak{S}'\to\mathfrak{S}$ of the base formal prescheme.
  \item[\rm{(iv)}] If the composite $g\circ f$ of two morphisms is separated,
  $f$ is separated.
\end{enumerate}
\emph{(It is understood in this statement that if the same formal prescheme
$\mathfrak{Z}$ occurs several times in the same proposition, it is
regarded as the inductive limit of the \emph{same} sequence $(Z_n)$ of
ordinary preschemes wherever it appears, and the morphisms from
$\mathfrak{Z}$ to a formal prescheme (resp. from a formal prescheme to
$\mathfrak{Z}$) as inductive limits of morphisms from the $Z_n$ to ordinary
preschemes (resp. from ordinary preschemes to the $Z_n$)).}
\end{proposition}

With the notation of (\hyperref[I.10.15.2]{10.15.2}), we indeed have
$(g\circ f)_0=g_0\circ f_0$ and
$(f\times_{\mathfrak{S}}g)_0=f_0\times_{S_0}g_0$; the assertions of
(\hyperref[I.10.15.3]{10.15.3}) are then immediate consequences of
(\hyperref[I.10.15.2]{10.15.2}) and of the corresponding assertions of
(\hyperref[I.5.5.1]{5.5.1}) for ordinary preschemes.

We leave it to the reader to state, for the same kind of
formal preschemes and morphisms as in
(\hyperref[I.10.15.3]{10.15.3}), the propositions corresponding to
(\hyperref[I.5.5.5]{5.5.5}),

\oldpage[I]{213}
(\hyperref[I.5.5.9]{5.5.9}) and
(\hyperref[I.5.5.10]{5.5.10}) (replacing in them ``affine open subset'' by
``affine formal open subset satisfying condition b) of
(\hyperref[I.10.6.3]{10.6.3})'').

An analogous argument also shows that every \emph{Noetherian} affine formal
scheme is separated, which justifies the terminology.

\begin{proposition}[10.15.4]
\label{I.10.15.4}
Let $\mathfrak{S}$ be a locally Noetherian formal prescheme, and
$\mathfrak{X}$, $\mathfrak{Y}$ two locally Noetherian formal
$\mathfrak{S}$-preschemes such that $\mathfrak{X}$ or $\mathfrak{Y}$ is of
finite type over $\mathfrak{S}$ (so that
$\mathfrak{X}\times_{\mathfrak{S}}\mathfrak{Y}$ is locally Noetherian)
and that $\mathfrak{Y}$ is separated over $\mathfrak{S}$. Let
$f:\mathfrak{X}\to\mathfrak{Y}$ be an $\mathfrak{S}$-morphism; then the
graph morphism
$\Gamma_f=(1_{\mathfrak{X}},f)_{\mathfrak{S}}:\mathfrak{X}\to
\mathfrak{X}\times_{\mathfrak{S}}\mathfrak{Y}$ is a closed immersion.
\end{proposition}

We may suppose that $\mathfrak{S}$ is the inductive limit of a sequence
$(S_n)$ of locally Noetherian preschemes, $\mathfrak{X}$ (resp.
$\mathfrak{Y}$) the inductive limit of a sequence $(X_n)$ (resp. $(Y_n)$) of
$S_n$-preschemes, and $f$ the inductive limit of a sequence of
$S_n$-morphisms $f_n:X_n\to Y_n$; then
$\mathfrak{X}\times_{\mathfrak{S}}\mathfrak{Y}$ is the inductive limit of the
sequence $(X_n\times_{S_n}Y_n)$, and $\Gamma_f$ that of the sequence
$(\Gamma_{f_n})$ (\hyperref[I.10.7.4]{10.7.4}); by hypothesis, $Y_0$ is
separated over $S_0$ (\hyperref[I.10.15.2]{10.15.2}), so the space
$\Gamma_{f_0}(X_0)$ is a closed subspace of
$X_0\times_{S_0}Y_0$; since the underlying spaces of
$\mathfrak{X}\times_{\mathfrak{S}}\mathfrak{Y}$ (resp.
$\Gamma_f(\mathfrak{X})$) and $X_0\times_{S_0}Y_0$ (resp.
$\Gamma_{f_0}(X_0)$) are the same, we already see that
$\Gamma_f(\mathfrak{X})$ is a \emph{closed} subspace of
$\mathfrak{X}\times_{\mathfrak{S}}\mathfrak{Y}$. Let us now remark that
when $(U,V)$ runs through the set of pairs consisting of a Noetherian affine
formal open subset $U$ (resp. $V$) of $\mathfrak{X}$ (resp.
$\mathfrak{Y}$) such that $f(U)\subset V$, the open subsets
$U\times_{\mathfrak{S}}V$ form a covering of
$\Gamma_f(\mathfrak{X})$ in
$\mathfrak{X}\times_{\mathfrak{S}}\mathfrak{Y}$, and if
$f_U:U\to V$ is the restriction of $f$ to $U$,
$\Gamma_{f_U}:U\to U\times_{\mathfrak{S}}V$ is the restriction of
$\Gamma_f$ to $U$. If we show that $\Gamma_{f_U}$ is a closed
immersion, the same will therefore be true of $\Gamma_f$
(\hyperref[I.10.14.4]{10.14.4}); in other words, we are reduced to the
case where
$\mathfrak{S}=\operatorname{Spf}(A)$,
$\mathfrak{X}=\operatorname{Spf}(B)$,
$\mathfrak{Y}=\operatorname{Spf}(C)$ are affine ($A$, $B$, $C$ Noetherian
adic rings), $f$ corresponding to a continuous $A$-homomorphism
$\varphi:C\to B$; $\Gamma_f$ then corresponds to the unique continuous
homomorphism $\omega:B\widehat{\otimes}_A C\to B$ which, composed with the
canonical homomorphisms $B\to B\widehat{\otimes}_A C$ and
$C\to B\widehat{\otimes}_A C$, gives the identity and $\varphi$
respectively. Now, it is clear that $\omega$ is \emph{surjective}, whence
our assertion.

\begin{corollary}[10.15.5]
\label{I.10.15.5}
Let $\mathfrak{S}$ be a locally Noetherian formal prescheme, and
$\mathfrak{X}$ an $\mathfrak{S}$-prescheme of finite type; for
$\mathfrak{X}$ to be separated over $\mathfrak{S}$, it is necessary and
sufficient that the diagonal morphism
$\mathfrak{X}\to\mathfrak{X}\times_{\mathfrak{S}}\mathfrak{X}$ be a
closed immersion.
\end{corollary}

\begin{proposition}[10.15.6]
\label{I.10.15.6}
A closed immersion $j:\mathfrak{Y}\to\mathfrak{X}$ of locally Noetherian
formal preschemes is a separated morphism.
\end{proposition}

With the notation of (\hyperref[I.10.14.2]{10.14.2}),
$j_0:Y_0\to X_0$ is a closed immersion, hence a separated morphism, and it
suffices to apply (\hyperref[I.10.15.2]{10.15.2}).

\begin{proposition}[10.15.7]
\label{I.10.15.7}
Let $X$ be a locally Noetherian (ordinary) prescheme, $X'$ a closed
subset of $X$, and $\widehat{X}=X_{/X'}$. For $\widehat{X}$ to be separated,
it is necessary and sufficient that $\widehat{X}_{\mathrm{red}}$ be so, and
it suffices that $X$ be so.
\end{proposition}

Indeed, with the notation of (\hyperref[I.10.8.5]{10.8.5}), for
$\widehat{X}$ to be separated, it is necessary and sufficient that $X'_0$ be
so (\hyperref[I.10.15.2]{10.15.2}), and since
$\widehat{X}_{\mathrm{red}}=(X'_0)_{\mathrm{red}}$, it is equivalent to
saying that $\widehat{X}_{\mathrm{red}}$ is so
(\hyperref[I.5.5.1]{5.5.1, (vi)}).
